\chapter{梵网经}

\section{论游行者}

{\Large 1.} \textbf{如是我闻\footnote{此经旧译见\textbf{长阿含经}·梵动经，以及支谦的\textbf{佛说梵网六十二见经}。}。一时，世尊行于王舍城与那烂陀之间的旅途，与大比丘僧团五百比丘一起。游行者善念也行于王舍城与那烂陀之间的旅途，与弟子梵赐学童一起。于此，游行者善念以多种方法说佛陀的毁谤，说法的毁谤，说僧的毁谤，而游行者善念的弟子梵赐学童以多种方法说佛陀的赞誉，说法的赞誉，说僧的赞誉。如此，这师弟二人彼此的论说正相敌对，一直跟随在世尊与比丘僧团的背后。}

\begin{enumerate}\item 当这第一次大结集进行中，在结集律的终了，大迦叶尊者便问到经藏中首部首经之梵网——「朋友阿难！梵网是在何处所说」，在如是等所述之语的终了，为阐明在何处所说及就何所说等的一切，阿难尊者便说了「如是我闻」等。因此说「梵网之初的『如是我闻』，由阿难尊者在第一次大结集时所说，为因缘之初」\footnote{此引文即「开篇辞」的末句。}。
\item 这里，\textbf{如是}，即不变词。\textbf{我}等是名词。在\textbf{行}之中，paṭi 是前缀词，hoti 是动词。当知先以此法分析词语。
\item 而从意义上，先说\textbf{如是}一词，有譬喻、指示、高兴、呵责、领受话语、行相、显示、强调等多种义项。如在\begin{quoting}如是生为人，当作诸善事。{\bluefootnotesize 〔法句第 53 颂〕}\end{quoting}等处为譬喻，在\begin{quoting}你当如是进，你当如是返。{\bluefootnotesize 〔增支部第 4.122 段〕}\end{quoting}等处为指示，在\begin{quoting}此即如是，世尊！此即如是，善逝！{\bluefootnotesize 〔增支部第 3.66 段〕}\end{quoting}等处为高兴，在\begin{quoting}如是如是，这贱女在任何场合都要称赞那秃头沙门。{\bluefootnotesize 〔相应部第 1.187 段〕}\end{quoting}等处为呵责，在\begin{quoting}「如是，尊者！」彼诸比丘答世尊。{\bluefootnotesize 〔中部第 1.1 段〕}\end{quoting}等处为领受话语，在\begin{quoting}我实如是了知世尊所开示的法。{\bluefootnotesize 〔中部第 1.398 段〕}\end{quoting}等处为行相，在\begin{quoting}来！学童！去沙门阿难处！去后，以我的话语，问沙门阿难少病少恼、起居轻利、有力、安住：「可教之子、学童须婆问阿难您少病少恼、起居轻利、有力、安住？」且如是说：「善哉！恳请阿难您去可教之子、学童须婆的住处，怜悯故！」{\bluefootnotesize 〔长部第 1.445 段〕}\end{quoting}等处为显示，在\begin{quoting}迦罗摩！尔等以为如何，这些法是善还是不善？不善，尊者！有过还是无过？有过，尊者！为智者所呵责还是称赏？为智者所呵责，尊者！成就、受持者，导向非利与苦还是不然？尔等对此作如何？尊者！成就、受持者，导向非利与苦，我等对此作如是。{\bluefootnotesize 〔增支部第 3.66 段〕}\end{quoting}等处为强调，此处当视为行相、显示、强调。
\item 这里，以\textbf{行相}之义，\textbf{如是}一词即表明此义：彼世尊的话语具种种理法之微妙，等起多种意乐，具足义文，去除种种（敌对），于法、义、开示、通达等甚深，由随适一切有情各自的语言而来至耳路，谁能以一切品类了知此？而以所有力气为令生起欲闻，故说「如是我闻，我唯以一种行相得闻」。以\textbf{显示}之义，即以「我非自成者，此非由我证」为己开脱，显示现在当说的全部经文为「如是我闻，我如是得闻」。以\textbf{强调}之义，即为显明随适于世尊以\begin{quoting}诸比丘！我的声闻比丘中，多闻之最上，即此阿难，智慧、忆念、坚定、给侍之最上，即此阿难。{\bluefootnotesize 〔增支部第 1.223 段〕}\end{quoting}与法将以\begin{quoting}阿难尊者善巧于义，善巧于法，善巧于文，善巧于辞，善巧于前后。{\bluefootnotesize 〔增支部第 5.169 段〕}\end{quoting}所赞叹的自身的忆持力，令有情愿欲听闻，故说「如是我闻，且其于义、于文不增不减，唯当如是而见，而非其它」。\footnote{此段亦见于\textbf{经集}·耕田婆罗豆婆遮经注。}
\item \textbf{我}一词见于三义\footnote{这是在解释 me 字可兼作具格、为格、属格之用。}。如在\begin{quoting}我不应受用吟颂之物。{\bluefootnotesize 〔经集第 81 颂〕}\end{quoting}等处为「由我」之义，在\begin{quoting}善哉！尊者！请世尊为我略开示法！{\bluefootnotesize 〔相应部第 4.88 段〕}\end{quoting}等处为「为我」之义，在\begin{quoting}诸比丘！请作我法的继承者！{\bluefootnotesize 〔中部第 1.29 段〕}\end{quoting}等处为「我的」之义，而此处以「由我闻」及「我的闻」二义合适。
\item \textbf{闻}一词，带前缀与不带前缀，有服行、闻名、染污、积累、践行、耳所识者、随从耳门之所知等多种义项。如在「从事 \textit{pasuto} 戎行」等处为服行，在\begin{quoting}见所闻法者……{\bluefootnotesize 〔自说第 11 经〕}\end{quoting}等处义为著名之法，在\begin{quoting}心怀漏泄 \textit{avassutā}，（受用）心怀漏泄的……{\bluefootnotesize 〔波夷提第 657 段〕}\end{quoting}等处义为染污、不染污，在\begin{quoting}尔等产生 \textit{pasutaṃ} 之福德匪浅。{\bluefootnotesize 〔小诵第 7.12 颂〕}\end{quoting}等处义为积累，在\begin{quoting}智者修禅定。{\bluefootnotesize 〔法句第 181 颂〕}\end{quoting}等处义为践行禅那，在\begin{quoting}所见、所闻、所觉。{\bluefootnotesize 〔中部第 1.241 段〕}\end{quoting}等处义为耳所识者，在\begin{quoting}忆持所闻、积聚所闻。{\bluefootnotesize 〔中部第 1.339 段〕}\end{quoting}等处义为随从耳门之所知，而此处其义为因随从耳门之所理解或能理解。因为当\textbf{我}字义为「由我」时，作「如是由我所闻」，则因随从耳门之所理解适合，当义为「我的」时，作「我的所闻如是」，则因随从耳门之能理解适合。
\item 如是，在此三词中，\textbf{如是}显示耳识等识的作用，\textbf{我}显示据有所说之识的人，\textbf{闻}由排除无闻之相，显示无减无增无颠倒之把握。同样，\textbf{如是}阐明经由这因随从耳门转起的识路、以种种品类于所缘所转之相，\textbf{我}阐明自我，\textbf{闻}阐明法。此中的略义为：以种种品类于所缘转起的识路，我不造作其它，而造作此，此法即所闻。
\item 同样，\textbf{如是}阐明应予说明之法，\textbf{我}阐明人，\textbf{闻}阐明人之作用。这即是说：我将说明所闻，它由我如是闻。
\item 同样，对此心相续的种种行相之所转，而有种种义、文之把握，\textbf{如是}即说明此种种行相——因为「如是」即这行相的概念，\textbf{我}说明作者，\textbf{闻}说明境域。至此，以种种行相转起之心相续，对据有此者，其作者、境域之把握业已决定。
\item 或者，\textbf{如是}即说明人的作用，\textbf{闻}说明识的作用，\textbf{我}说明与二作用相应的人。而此中的略义为：由我——据有听的作用之识的人——以得称为听的作用之识所闻。
\item 这里，「如是」与「我」依真实义与第一义是不存在的概念。因为此中从第一义上存在什么可以得说为「如是」或「我」呢？「闻」是存在的概念。因为凡是此中以耳获得的，从第一义上它就存在。同样，「如是」与「我」由执取彼彼而说，为执取的概念，「闻」由相较于见等而说，为相较的概念。
\item 且此中，以「如是」之语表明不迷乱。因为迷乱者不堪种种品类之通达。以「闻」之语表明对所闻的不忘失。因为对所闻忘失者，则过时不能自认「由我所闻」。如此，以不迷乱有慧的成就，而以不忘失有念的成就。这里，以慧为前导的念堪能明确文句，以念为前导的慧堪能通达意义，对具足结合义、文这二种堪能者，由堪能守护法藏故，有法的司库的成就。
\item 另一法：以「如是」之语表明如理作意。因为对不如理作意者，无有种种品类的通达故。以「闻」之语表明不散乱，由心散乱者无有听闻故。同样，因为心散乱的人即便当以一切成就被说，还会说「我尚未闻，请再说」。且此中，以如理作意证明自身正誓愿及于先前曾作福德，由无正誓愿或先前未作福德者无有此故。以不散乱证明听闻正法及依止善人。因为心散乱者不能听闻，且不依止善人则无听闻。
\item 另一法：因为说「对此心相续的种种行相之所转，而有种种义、文之把握，『如是』即说明此种种行相」，且这如是贤善的行相，非无正誓愿或先前未作福德者所有，所以「如是」以贤善的行相表明自己后二轮的成就，「闻」以致力于听闻表明前二轮\footnote{前二轮、后二轮：即四轮，见\textbf{增支部}第 4:31 经，依次为住于适宜处、依止善人、自身正誓愿、先前曾作福德。}的成就。因为对住于不适宜处者，或无有善人可依止者，则无有听闻。如此，以后二轮的成就有意乐清净的成就，以前二轮的成就有加行清净，且以此意乐清净有精通证得的成就，以加行清净有精通阿笈摩的成就。如此加行与意乐清净、阿笈摩与证得具足者的话语，好似太阳升起前的曙光，又好似善业的如理作意，应得作为世尊言语的前导，基于此，便置诸因缘，说了「如是我闻」等。
\item 另一法：以「如是」这表明通达种种品类的话语，表明义、辩无碍解成就的自性，以「闻」这表明通达当闻之分类的话语，表明法、词无碍解成就的自性。且在说「如是」这表明如理作意的话语时，表明「这些法由我以意熟虑、以见善通达」，在说「闻」这表明致力于听闻的话语时，表明「许多法由我听闻、忆持，以语串习」。亦以此二者表明义、文的圆满，令于听闻生起尊敬。因为不以尊敬听闻义、文圆满的法，则撇离巨大的利益，所以此法当令生起尊敬，恭敬地听闻。
\item 而以「如是我闻」这一整句，阿难尊者以不把如来宣说之法归于自己而超迈了非善人之地，以自认弟子而跃入善人之地。同样，令心从非善法出起，住立于善法。以表明「这全都由我所闻，实为世尊的话语」为己开脱，指认大师，引入胜者之语，令住立于法之准绳。
\item 复次，以「如是我闻」自认不由自己所创，揭示前语——「这彼世尊、四无所畏、持十力、居牛王位、作狮子吼、一切有情之最上、法自在、法王、法主、法洲、法皈依、转善法高贵之轮、正等正觉者的话语，由我当面受纳，对此中的义、法、句、文不当有疑惑或疑虑」，令一切人天于此法消除不信，生起信的成就。因此说：\begin{quoting}消除不信，于教法增长信，\\如是作「如是我闻」论者，为乔达摩弟子。\end{quoting}
\item \textbf{一}说明数目的限定，\textbf{时}说明被限定者，\textbf{一时}是不定的阐述。这里，「时」一词\begin{quoting}于结合、刹那、时间、集合、因、见、\\获得、舍弃及通达处出现。\end{quoting}如在\begin{quoting}也许明天我们得了时机会前往。{\bluefootnotesize 〔长部第 1.447 段〕}\end{quoting}等处义为结合，在\begin{quoting}对住梵行者，唯有一个时机。{\bluefootnotesize 〔增支部第 8.29 段〕}\end{quoting}等处为刹那，在\begin{quoting}热时，烧时……{\bluefootnotesize 〔波夷提第 358 段〕}\end{quoting}等处为时间，在\begin{quoting}在林中有大集会。{\bluefootnotesize 〔长部第 2.332 段〕}\end{quoting}等处为集合，在\begin{quoting}跋陀利！你未通达时机——「世尊住舍卫国，世尊亦将了知我，名为跋陀利的比丘于大师教法之学未能圆满」，跋陀利！你未通达这时机。{\bluefootnotesize 〔中部第 2.135 段〕}\end{quoting}等处为因，在\begin{quoting}尔时，游行者乌伽诃摩那、沙门文祁之子，进入教义议论处、异色柿林、一间堂之茉莉园。{\bluefootnotesize 〔中部第 2.260 段〕}\end{quoting}等处为见，在\begin{quoting}现法中的义利，与来世的义利，\\贤人现观义利，便被称为智者。\footnote{现观 \textit{abhisamaya}：这里解释带前缀 abhi 的 samaya，下文解释「舍弃、通达」处亦同。}{\bluefootnotesize 〔相应部第 1.128 段〕}\end{quoting}等处为获得，在\begin{quoting}彻底现观了慢，他便尽了苦边。{\bluefootnotesize 〔增支部第 7.9 段〕}\end{quoting}等处为舍弃，在\begin{quoting}苦的逼迫义、有为义、热恼义、变坏义、现观义。{\bluefootnotesize 〔无碍解道第 108 段〕}\end{quoting}等处为通达，而此处义为时间。以此表明年、季、月、半月、日夜、上午、中午、哺时、初中后夜、瞬间等作为时间种类的「时」中的一时。
\item 这里，虽然在这些年等的时中，彼彼经于彼彼年、季、月、半月的夜分或日分所说，这一切长老都以慧善知晓、善确定，然而，因为当说「如是我闻，某年、某季、某月、某半月的某夜分或某日分」时，不易于持、诵或教人诵，且需说许多，所以唯以一句汇集其义，而说「一时」。
\item 或者，在这些入胎之时、出生之时、悚惧之时、出离之时、行苦行之时、战胜魔罗之时、现等觉之时、现法乐住之时、开示之时、般涅槃之时如是等世尊于人天中极其闪耀的多个时段的时中，表明被称为开示之时的「一时」。且在智与悲作用之时中为悲作用之时，在自利、利他行道之时中为利他行道之时，在为聚会者应作的二时中为说法之时\footnote{\textbf{疏}云「智作用之时、自利行道之时为\textbf{现等觉之时}，圣默然之时为\textbf{现法乐住之时}，悲作用、利他行道、说法之时为\textbf{开示之时}」，则「为聚会者应作的二时」的另一当为圣默然之时。}，在开示与行道之时中为开示之时，就这些时中的某一时而说「一时」。
\item 那为什么此处不同于阿毗达摩中「在尔时 \textit{samaye}，欲界……」{\bluefootnotesize 〔法集论第 1 段〕}，以及此外的经句中「在尔时，诸比丘！比丘已离诸欲……」以依格作说明，也不同于律中「以尔时 \textit{samayena}，佛世尊……」以具格，而以业格「一时」作说明呢？由彼处与此处意义之形成不同故。
\item 因为在彼处的阿毗达摩及此外的经句中，形成起因义，并以状态形成状态之特相义。起因者，时间义及集合义为「时」，而对彼彼处所说的触等诸法，以被称为刹那、结合、因的「时」的状态，它们的状态得以明确，所以，为阐明此义，在彼处以依格作说明。
\item 而在律中，形成因义及造作义。因为凡学处施设之时，即便舍利弗等也难以知晓，在作为因、作为造作的尔时，当施设学处并关切学处施设之因时，世尊便住于彼彼处，所以，为阐明此义，在彼处以具格作说明。
\item 而在此处及其它如是种类处，形成绝对系缚义。因为世尊开示此或其它经之时，此时唯绝对以悲住而住，所以，为阐明此义，在此处以业格作说明。因此说：\begin{quoting}虑及彼彼之义，以依与具\\说别处之时，此处则以业。\end{quoting}然而，古人们解释道「在尔时，或以尔时，或一时，这只是言谈的差别，一切处都是依（格）之义」，所以即便说「一时」，当知为「在一时」之义。
\item \textbf{世尊}，即老师。因为他们称世间的老师为世尊，且他以一切功德殊胜为一切有情的老师，所以当知为世尊。古人们也说：\begin{quoting}世尊是最胜之语，世尊是最上之语，\\这老师应受尊重，因此被称为世尊。\end{quoting}复次，依于此颂——\begin{quoting}具吉祥，具破，与卓越相应，且具分别，\\具亲近，离弃于诸有而行，所以为世尊。\end{quoting}可知该词的详义，且其在「\textbf{清净道论}·说佛随念」中已述。
\item 且至此，此中以「如是我闻」之语显明所闻之法，使世尊的法身成为现量，以「这不是大师之后的宣教，这即尔等的大师」安抚因未见大师而不满的人。以「一时，世尊」之语显明在那时世尊之不存在性，证明色身之般涅槃，以「如是种类圣法的开示者、持十力、身等同金刚聚的彼世尊尚且般涅槃，还因何对活命生起希望」令仅迷醉于活命的人们悚惧，并令其于正法生起勇猛。且在说「如是」时，说明开示的成就，「我」为弟子的成就，「一时」为时间的成就，「世尊」为开示者的成就。
\item \textbf{王舍城与那烂陀之间}，「之间」一词见于原因、刹那、心、中间、间隔等处。在\begin{quoting}除却如来，谁能知晓其间？{\bluefootnotesize 〔增支部第 6.44 段〕}\end{quoting}与\begin{quoting}人们聚集后，谈论着我、你与其间。{\bluefootnotesize 〔相应部第 1.228 段〕}\end{quoting}等处「之间」一词作原因，在\begin{quoting}尊者！某个正在洗器皿的女人在电光之间看见了我。{\bluefootnotesize 〔中部第 2.149 段〕}\end{quoting}等处作刹那，在\begin{quoting}若其中间无诸忿恨。{\bluefootnotesize 〔自说第 20 经〕}\end{quoting}等处作心，在\begin{quoting}他便中道而废。{\bluefootnotesize 〔小品第 350 段〕}\end{quoting}等处作中间，在\begin{quoting}复次，诸比丘！这温泉从二大地狱之间而来。{\bluefootnotesize 〔波罗夷第 231 段〕}\end{quoting}等处作间隔，此处则作间隔合适，所以当知此中之义为「王舍城与那烂陀的间隔」。而由与「之间」一词连接故，用了业格。且在类此等处，晓韵律者以「行于村落与河流之间」，唯用一个「之间」\footnote{义注这样说，是因为「王舍城与那烂陀之间」的原文用了两个「之间 \textit{antarā}」。}，它还应与第二词相连，若未被连接，则不成业格，而此处即是在连接后而说。
\item \textbf{行于旅途}，即行于被称为旅行的道路，意即长路。因为在「经分别」中，对行于旅途之时，有「我将走半由旬而食」{\bluefootnotesize 〔波夷提第 218 段〕}等语故，半由旬亦为旅途，而从王舍城到那烂陀适一由旬。
\item \textbf{与大比丘僧团一起}，大，即因功德大为大，数目大为大。因为这比丘僧团以功德而为大，由具足少欲等的功德故，亦以数目为大，由数有五百故。众比丘的僧团为「比丘僧团」，与此比丘僧团，即与被称为「见、戒共同之聚」的沙门众之义。\textbf{一起}，即一同。
\item \textbf{五百比丘}，即以「彼等以五为量」为五。量被称为度，所以，好比当说「饮食知量」时，意即「他于饮食知量、知度」，如是，此处也当视其义为「这些百比丘以五为量、以五为度」。众比丘之百即百比丘，即与此五百比丘。
\item \textbf{游行者善念}，善念是其名。\textbf{也}字意为联系同分行路之人。\textbf{Kho} 字作词的连接，因文辞的平顺而说。\textbf{游行者}，即删阇夜\footnote{删阇夜·毗罗胝子 \textit{Sañjaya Belaṭṭhaputta}：六师外道之一，\textbf{杂阿含经}作「先阇那」。}的弟子，穿衣游行者。这即是说「当世尊行于此旅途时，游行者善念也曾行」—— hoti 一词在此意为过去时。
\item 在\textbf{与弟子梵赐学童一起}中，住在一边者为弟子，意即行于附近、侍于跟前、学生。梵赐是其名。\textbf{学童}，即有情、盗贼、青年之称\footnote{有情、盗贼之义，似从「人类 \textit{manu}」而来。}。因为在\begin{quoting}那些放逸的学童，为天使所督促，\\他们长夜忧伤，经历卑劣的生命。{\bluefootnotesize 〔中部第 3.271 段〕}\end{quoting}等中称有情为学童，在\begin{quoting}与业已犯事或尚未犯事的学童们相会。{\bluefootnotesize 〔中部第 2.149 段〕}\end{quoting}等中为盗贼，在\begin{quoting}阿摩昼学童，鸯伽学童。{\bluefootnotesize 〔长部第 1.316 段〕}\end{quoting}等中称青年为学童，亦即此处之义。这即是说：与名为梵赐的青年弟子一起。
\item \textbf{于此}，即在此旅途中，或于此二人中。\textbf{Sudaṃ} 只是不变词。\textbf{多种方法}，先说方法一词，适于顺次、开示、原因等处。即在\begin{quoting}阿难！今天是谁的顺次去教诫众比丘尼？{\bluefootnotesize 〔中部第 3.398 段〕}\end{quoting}等处，方法一词适于顺次，在\begin{quoting}你应受持此蜜丸方法！{\bluefootnotesize 〔中部第 1.205 段〕}\end{quoting}等处为开示，在\begin{quoting}以此方法，王官！应当如是……{\bluefootnotesize 〔长部第 2.411 段〕}\end{quoting}等处为原因，而它在此处适于原因，所以此中之义即是说，以多种种类的原因、以许多原因。
\item \textbf{说佛陀的毁谤}，即便对无可毁谤、具足无量赞誉的佛世尊，也以「对世间年耆者应作的问讯等礼待之业，被称为和合之味，沙门乔达摩无此，所以样貌无味，沙门乔达摩离财，是无作论者、断灭论者、嫌厌者、虚无论者、苦行者、去胎者。沙门乔达摩无上人法、无足称圣智见的殊胜。沙门乔达摩开示以寻思推导、以考察竞逐、由自身理解的法。沙门乔达摩非一切知，非世间解，非无上士，非最上人」等，说彼彼非因为因，如是如是说毁谤、过失、非难。
\item 且如对佛陀，亦如是对法说彼彼非因为因——沙门乔达摩的法是恶说、难通达、非出离、非导向寂止——如是如是说毁谤。
\item 且如对法，亦如是对僧说彼彼非因为因——沙门乔达摩的声闻僧团是邪行道、曲行道，行敌对、不顺、顺于非法之道——如是如是说毁谤。
\item 而其弟子却说：「我们的阿阇黎执取不应执取者，踩踏不应踩踏者，他好像吞下了火，好像用手碰了刀锋，好像想用拳头砸碎须弥卢，好像在锯齿边缘嬉戏，好像用手去制服发情的恶象，说着值得赞誉的三宝的毁谤，将至于厄难。而当阿阇黎踩到粪便、或火堆、或荆棘、或黑蛇，或踏上枪尖，或吃了猛烈的毒药，或绊倒在灰水河\footnote{灰水河 \textit{khārodaka}：为地狱中的专名，或可按字面理解为「碱水」。PTS 本此处作「碎裂的指甲 \textit{nakha-bhedaka}」。}，或跌落地狱的崖壁，弟子不应效仿这一切。因为有情是业之自身，他们唯去到与业相适之趣。父既不能以子之业而去，子亦不能以父之业，母不能以子，子不能以母，兄弟不能以姊妹，姊妹不能以兄弟，阿阇黎不能以弟子，弟子不能以阿阇黎之业而去。我的阿阇黎说三宝的毁谤，而指责圣人有大罪过。如是经如理运思，站在阿阇黎的论说之上，唯从原因陈述正确的原因，以多种方法去说三宝的赞誉，如同智者样的族姓子一般。」因此说：\textbf{而游行者善念的弟子梵赐学童以多种方法说佛陀的赞誉，说法的赞誉，说僧的赞誉}。
\item 这里，\textbf{赞誉}，赞誉一词见于形相、出生、色处、原因、度量、功德、称赏等处。这里，在\begin{quoting}变化为大蛇王的样貌。{\bluefootnotesize 〔相应部第 1.142 段〕}\end{quoting}等处被称为形相，在\begin{quoting}唯婆罗门是最胜的种姓，其他则为卑劣的种姓。{\bluefootnotesize 〔中部第 2.402 段〕}\end{quoting}等处为出生，在\begin{quoting}具足第一肤色之美。{\bluefootnotesize 〔长部第 1.303 段〕}\end{quoting}等处为色处，在\begin{quoting}我不带走，也不破坏，远远地嗅莲花，\\那又以何原因被称为偷香者？{\bluefootnotesize 〔相应部第 1.234 段〕}\end{quoting}等处为原因，在\begin{quoting}有三种钵色。{\bluefootnotesize 〔波罗夷第 602 段〕}\end{quoting}等处为度量，在\begin{quoting}长者！你何时编集了这沙门乔达摩的功德？{\bluefootnotesize 〔中部第 2.77 段〕}\end{quoting}等处为功德，在\begin{quoting}说值得赞誉者的赞誉。{\bluefootnotesize 〔增支部第 2.135 段〕}\end{quoting}等处为称赏，此处为功德和称赏。据说，他唯陈述彼彼真实之因，以多种方法说了三宝所有的功德和赞誉。
\item 这里，当知以\begin{quoting}彼世尊亦即是阿罗汉、正等正觉者……{\bluefootnotesize 〔波罗夷第 1 段〕}\end{quoting}等方法，与\begin{quoting}诸比丘！于佛陀净喜者，于最上净喜。{\bluefootnotesize 〔增支部第 4.34 段〕}\end{quoting}等，以及\begin{quoting}诸比丘！一人出现于世……无等，无与相等。{\bluefootnotesize 〔增支部第 1.174 段〕}\end{quoting}等方法赞誉佛陀。当知以\begin{quoting}法是世尊所善说……{\bluefootnotesize 〔长部第 2.159 段〕}\end{quoting}\begin{quoting}根除执著，断绝流转。{\bluefootnotesize 〔增支部第 4.34 段〕}\end{quoting}\begin{quoting}诸比丘！于八支圣道净喜者，于最上净喜。\end{quoting}等方法赞誉法。当知以\begin{quoting}声闻僧团是世尊的善行道者……{\bluefootnotesize 〔长部第 2.159 段〕}\end{quoting}\begin{quoting}诸比丘！于僧净喜者，于最上净喜。{\bluefootnotesize 〔增支部第 4.34 段〕}\end{quoting}等方法赞誉僧。
\item 然而，有能力的论法者应在深入五部、九分大师教、八万四千法蕴后，阐明佛等的赞誉。因为在此基础上，当阐明佛等的功德时，不能说论法者是以邪径唐突。唯在类此等处，当知论法者的力量。而学童梵赐只是跟随传闻等，以自己的力量说三宝的赞誉。
\item \textbf{如此，这师弟二人}，即如是，这二阿阇黎与弟子。\textbf{彼此}，即两两。\textbf{论说正相敌对}，即无丝毫回避，径直作种种敌对的论说、多番作相违的论说之意。因为当阿阇黎说三宝的毁谤时，弟子说赞誉，再一个毁谤、一个赞誉，如是，阿阇黎好比在坚实的木板敲入毒树之楔一般，反复说三宝的毁谤，而弟子好比以金银摩尼制成的楔子排挤出那楔一般，反复说三宝的赞誉。因此说「论说正相敌对」。\textbf{一直跟随在世尊与比丘僧团的背后}，即一直在世尊与比丘僧团的后面，不离视线，以相同的威仪跟随，瞻首随行之意。
\item 那么为什么世尊要行此旅途？善念又为什么跟随？他又为什么说三宝的毁谤？首先，世尊在那时在王舍城周围的十八大寺中的某处住了之后，早上料理了身体，在行乞之时，为比丘僧团围绕，在王舍城行乞。他在那天令比丘僧团的乞食易得，饭后从乞食返回，教比丘僧团拿了衣钵「我要去那烂陀」，便从王舍城离开，行于这旅途。
\item 而善念在那时也在王舍城周围的某处游行者园中住了之后，为游行者围绕，在王舍城行乞。他在那天令游行者会众的乞食易得，吃了早餐，教游行者拿了游行者的资具「我要去那烂陀」，并不知晓世尊已行于此路而跟随。假如他知晓，便不会跟随。他正不知晓而行时，延颈而望，便看见世尊正以佛光闪耀，好似围了红毯而移动的金山之巅。
\item 据说，在那时，六色光芒从十力的身体发散，在周围八十肘之量的地方扑朔，如宝珠、宝环、宝粉散落，如宝彩的金布敷展，如赤金的汁液挥洒，如百颗流星乱坠，如无间散落的翅子树花\footnote{翅子树 \textit{kaṇikāra}，其花为金色。}，如为风势吹散的红铅粉，又如闪烁的彩虹、电弧、星团、光簇的聚合，无间而然。
\item 而世尊的身体又为八十随形好所照亮，如盛开的青莲，如摇曳着满树花朵的波利质多树，如璀璨的星光，如以辉光笑对着苍穹，周匝一寻身光之光彩与其三十二高贵相之鬘相交织，如以排列的三十二月鬘、三十二日鬘的辉光依次胜过了排列的三十二转轮王、三十二帝释天王、三十二大梵的辉光。
\item 而且，围绕彼世尊站立的比丘们全都是少欲、知足、远离、无交际、敦促者、恶之指责者、演说者、接受劝诫者\footnote{接受劝诫 \textit{vacanakkhama}：字面意思是「忍受言语」。}，具足戒，具足定、慧、解脱、解脱智见。在彼等中间的世尊，如被红毯之墙包围的金柱，如行至红莲深处的金船，如被珊瑚栏楯包围的火聚，如被星团围绕的满月，甚至怡悦鸟兽之眼，遑论人天？
\item 而在那天，多数八十大长老把乌云之色的粪扫衣偏覆一肩，持了拄杖，如披挂善甲的香象，离去过失，吐出过失，破除烦恼，解开缚结，断除束缚，围绕着世尊。他自己离贪，为离贪者围绕，自己离嗔，为离嗔者围绕，自己离痴，为离痴者围绕，自己离渴爱，为离渴爱者围绕，自己无烦恼，为无烦恼者围绕，自己觉悟，为随觉者围绕，如被花瓣围绕的花丝，如被花丝围绕的莲蓬，如被八千大象围绕的六牙象王，如被九万天鹅围绕的持国鹅王，如被军行围绕的转轮王，如被天众围绕的帝释天王，如被梵众围绕的诃梨多大梵，以不可度量之时积聚的福力所转起、不可思议而无比的佛陀光彩，如月行空中般，行于这道路。
\item 于是，游行者见到如是以无比的佛陀光彩而行的世尊，以及目光下视、诸根寂静、心意寂静的众比丘礼敬着如住立于上空的满月般的世尊，便回看自己的会众。他们挑了扁担，满载了绑扎过的破旧板凳、三脚架、孔雀翎、土钵、包袱、水罐等诸多资具的重担，以「某人的手美，某人的脚美」等无义之语而饶舌，言语散漫，不足以观，非净喜者。他见到此后，便生起后悔。
\item 现在，他应当说世尊的赞誉，但因为利养恭敬的减少以及徒众的减少故，他始终嫉妒世尊。外道在佛未出世前，一直有利养恭敬，但从佛陀出世以后，利养恭敬便衰减了，好比在日出时萤火虫便归黯淡一般。且优波提舍与拘律陀\footnote{优波提舍与拘律陀：分别是舍利弗与目犍连的名字。}在删阇夜跟前出家时，游行者的会众甚大，但当这二人离开时，他俩的这会众也解散了。如此，当知因为这二重原因故，这游行者还始终嫉妒世尊，所以，为吐露嫉妒之毒，便只说三宝的毁谤。\end{enumerate}

{\Large 2.} \textbf{于是，世尊在王家芒果树园度过了一宿，与比丘僧团一起。游行者善念也在王家芒果树园度过了一宿，与弟子梵赐学童一起。于此，游行者善念以多种方法说佛陀的毁谤，说法的毁谤，说僧的毁谤，而游行者善念的弟子梵赐学童以多种方法说佛陀的赞誉，说法的赞誉，说僧的赞誉。如此，这师弟二人彼此的论说正相敌对而住。}

\begin{enumerate}\item \textbf{于是，世尊在王家芒果树园度过了一宿，与比丘僧团一起}，当世尊以此佛陀的光彩而行时，渐次到达芒果树园的门口，观察了太阳，想「现在前行非时，太阳将落」，便在王家芒果树园度过了一宿。
\item 这里，\textbf{芒果树园}是国王的园林。据说，在其门前有棵芒果幼树，人们便说它是「芒果树园」。园林因在其不远处，便也得称为芒果树园。它林泉具足，墙垣围绕，善结门户，如宝函般善加守护。人们为了国王能在此嬉戏，造了饰以壁画的家舍，它便被称为「王家」。
\item \textbf{善念也……}，善念也在此处观察了太阳，想「现在前行非时，有许多大大小小的游行者，且此道路有许多盗贼、凶猛的夜叉野兽的危险，而既然沙门乔达摩进了园林，且在沙门乔达摩所住之处有天人守护，那我也在此度过一宿，明天再去」，便进了这园林。随后，比丘僧团展示了对世尊的义务，便去观察各自的住处。游行者也在园林的一侧卸下游行者的资具，与自己的会众一起住下。但依经文中所说：\textbf{与}自己的\textbf{弟子梵赐学童一起}。
\item 而如是住下后，这游行者便在夜分观察十力。且在此时，周围灯火燃烧，如散落的众星，世尊坐在中间，比丘僧团围绕着世尊。在此，甚至没有一位比丘有手的不安、脚的不安，或发出咳嗽声、喷嚏声。因为此会众以自身业已修学与对大师的尊重等二重原因，如无风处的灯焰一般，毫无动摇地共坐。游行者看到这盛况，便观察自己的会众。在此，有人动手，有人动脚，有人悲叹，有人摇舌吐沫，磨牙打鼾，咕噜咕噜喘气而睡。他即便应当说三宝的功德之美，但因嫉妒，仍又开始毁谤，梵赐则仍以所说之法赞誉。因此而说「\textbf{于此，游行者善念……}」等当说的一切。这里，\textbf{于此}，即在此，即在芒果树的园林之义。\end{enumerate}

{\Large 3.} \textbf{于是，众多比丘在夜的黎明时分起身，于圆亭共坐聚集，便生起此名法：「希有！朋友！未曾有！朋友！彼世尊、知者、见者、阿罗汉、正等正觉者是何等善知晓有情的种种倾向！此游行者善念以多种方法说佛陀的毁谤，说法的毁谤，说僧的毁谤，而游行者善念的弟子梵赐学童以多种方法说佛陀的赞誉，说法的赞誉，说僧的赞誉。如此，这师弟二人彼此的论说正相敌对，一直跟随在世尊与比丘僧团的背后。」}

\begin{enumerate}\item \textbf{众多}，即许多。这里，依律的方法，三人被称为「众多」，过此为僧团，但依经的方法，三即是三，自此以后为众多，此处当知依经的方法为众多。\textbf{圆亭}，有些地方用了两片屋顶，以天鹅之翼般覆盖\footnote{以天鹅之翼般覆盖 \textit{haṃsa-vaṭṭaka-cchannena}：词典中 vaṭṭaka 为车，vaṭṭakā 为鹌鹑，并无「翼」义，考虑到 vaṭṭa 有回转义，故将其派生的 vaṭṭaka 作此译，取「有亭翼然」之意也。}所造的尖顶会堂被称为圆亭，有些地方用了一片屋顶，围以柱列所造的前厅也被称为圆亭，而此处当知以坐具之堂为圆亭。以坐下为\textbf{共坐}。以集合为\textbf{聚集}。\textbf{此名法}，谈论被称为「名」，即谈论之法之义。\textbf{便生起}，即已生起。那这是什么？即如「希有！朋友……」等。
\item 这里，如盲人登山般不常有为\textbf{希有}，这先是语源之法\footnote{语源之法 \textit{saddanaya}：义注这样说，是在以「盲人 \textit{andha}」和「常 \textit{nicca}」拼凑成「希有 \textit{acchariya}」一字。}。而义注之法为：适合弹指为\textbf{希有}，即适合击节之义。先前未曾有的存在为\textbf{未曾有}。此二者都是惊叹的同义语。\textbf{是何等}，即以此善知晓性显示不可量性。
\item \textbf{彼世尊、知者……善知晓}，此中的略义为：彼世尊圆满了完整的三十波罗蜜，粉碎了一切烦恼，现无上正等正觉，彼世尊\textbf{知}彼彼有情的意乐、随眠，\textbf{见}一切应知之法，如置于手掌中的阿摩勒果一般。
\item 复次，以宿住等\textbf{知}，以天眼而\textbf{见}。以三明或六通而\textbf{知}，于一切处以无碍之普眼而\textbf{见}。或者，以能知一切法之慧而\textbf{知}，对超过一切有情肉眼境域，乃至对墙后所及之色，以极清净的肉眼而\textbf{见}。或者，以成就自身利益、以三摩地为足处的通达之慧而\textbf{知}，以成就他人利益、以悲为足处的开示之慧而\textbf{见}。
\item 由破贼故及由应受资具等故，为\textbf{阿罗汉}。由完全且由自己觉悟一切法故，为\textbf{正等正觉者}。或者，知障碍之法，见出离之法，由破烦恼贼故为阿罗汉，由完全且由自己觉悟一切法故为正等正觉者，如是依四无所畏\footnote{四无所畏 \textit{catuvesārajja}：见\textbf{增支部}第 4:8 经等，\textbf{俱舍论}作「正等觉无畏、漏永尽无畏、说障法无畏、说出道无畏」。}，以四种行相赞美，是何等\textbf{善知晓有情的种种倾向}、种种意乐。
\item 现在，为显明其善知晓性，说了「此游行者……」等。这即是说：如世尊以\begin{quoting}诸比丘！有情依界而认同、体认，倾向下劣者与倾向下劣者认同、体认，倾向善者与倾向善者认同、体认。诸比丘！在过去时，有情亦曾依界而认同、体认……诸比丘！在未来时，有情仍将依界而认同、体认……诸比丘！现今，在现在时，有情依界而认同、体认，倾向下劣者与倾向下劣者认同、体认，倾向善者与倾向善者认同、体认。\end{quoting}所说的有情的种种倾向、种种意乐、种种见、种种偏爱、种种喜好，如以升计量、以秤称量一般，以种种倾向之智、一切知之智知晓，这是何等的善知晓！在世间，甚至同一意乐的两个有情也难得。当一欲行时，一则欲站，当一欲饮时，一则欲食。而在此二师弟中，\textbf{此游行者善念……一直跟随在世尊与比丘僧团的背后}。这里，\textbf{如此}，即「如是，此」之义。其余则如已述。\end{enumerate}

{\Large 4.} \textbf{于是，世尊知晓了彼众比丘的这名法，便往圆亭而去，去后，便坐在设好的坐处。世尊坐下后，便告众比丘：「诸比丘！现今你们因何谈论共坐聚集？你们又中断了什么闲谈？」}

\textbf{如是说已，彼众比丘便对世尊说：「在此，尊者！我们在夜的黎明时分起身，于圆亭共坐聚集，便生起此名法：『…………\footnote{连续两个省略号，表示译文省略了重复的经文，下例同。}』尊者！此即我们中断的闲谈，然后世尊到来。」}

\begin{enumerate}\item 在「\textbf{于是，世尊知晓了彼众比丘的这名法……}」中，\textbf{知晓}，即以一切知智了知。因为世尊在某处以肉眼见而了知，如\begin{quoting}世尊便看到一大段木材随恒河的水流漂荡。{\bluefootnotesize 〔相应部第 4.241 段〕}\end{quoting}等处，在某处以天眼见而了知，如\begin{quoting}世尊以超越人类的清净天眼，便看见那一千天人占据了波吒梨村内的基地。{\bluefootnotesize 〔长部第 2.152 段〕}\end{quoting}等处，在某处以自然之耳听而了知，如\begin{quoting}世尊便听到阿难尊者和善贤游行者的这会谈。{\bluefootnotesize 〔长部第 2.213 段〕}\end{quoting}等处，在某处以天耳听而了知，如\begin{quoting}世尊以超越人类的清净天耳界，便听到先陀那长者和尼拘律陀游行者的这会谈。{\bluefootnotesize 〔长部第 3.54 段〕}\end{quoting}等处，而此处是以一切知智听而了知。
\item 在作什么时了知？后夜义务。且此义务分有义利与无义利两种。这里，无义利的义务，已由世尊在菩提座上以阿罗汉道根除。唯有义利的义务是世尊的义务。它分五种：饭前义务、饭后义务、前夜义务、中夜义务、后夜义务。
\item 此处\footnote{以下介绍的几种佛陀的义务，基本同\textbf{经集}·耕田婆罗豆婆遮经的义注，唯后夜分作三分，较之\textbf{经集义注}少了以果定度过的一分。}，这饭前义务为：世尊清早起来后，为摄受侍者及身体安泰，作好洗脸等的洗漱，直至行乞之时，于远离的坐处等待，到了行乞之时，便著下衣、扎腰带、披覆上衣、持钵，时而独自，时而为比丘僧团围绕，进入村镇行乞，有时自然前去，有时则现起种种神变。
\item 例如，在世间之依怙进入行乞时，柔风在正前方开道，扫净地面。雨云洒下水珠，让路上的尘埃落定，在上方形成华盖。又一阵风收集了花朵，撒到路上。隆起的地面下沉，而下沉者隆起，使落足之时地面平整，或有乐触的莲花承受双足。当右足踏入因陀柱内，六色光芒从身上放出，使楼阁、尖顶等有如金黄的色泽，又如被彩布包裹，处处散射。象马飞鸟等于各自的所在发出甜美的声音，鼓、琵琶等乐器以及人们随身的璎珞也同样。人们以此征兆了知「今天世尊来此乞食」。
\item 他们整理衣裳，持了香、花等，从家出来，行到路边，以香、花等恭敬地供养世尊，顶礼后请求道「尊者！给我们十个比丘，二十个、五十个、一百个比丘」，乃至取了世尊的钵，设好座位，以食物恭敬地敬候。世尊食事已毕，观察了这些有情的心相续，开示以相应的法，以使有些人住于皈依，有些于五戒，有些于须陀洹、斯陀含、阿那含中的任一果，有些出家已于最上的阿罗汉果。如是，如此摄受大众后，便从坐起，去往寺庙。到达此后，他坐在圆亭中设好的最上佛座，等待诸比丘食事终了。随后，当诸比丘食事终了，侍者便报告世尊。于是，世尊便进入香房。以上即\textbf{饭前义务}。
\item 于是，世尊如是完成了饭前义务，坐于香房的前厅，清洗了双足，置于足凳，教诫比丘僧团「诸比丘！应以不放逸而成就，在世间佛陀出世难得、获得人身难得、成就难得、出家难得、听闻善法难得」。于此，有些问取世尊业处，世尊也依他们的性行授予业处。随后，他们全都礼拜了世尊，去到各自的夜住、昼住处，有些到林野，有些到树下，有些到山中等的某处，有些到四大王天的居处……有些到自在天的居处。随后，世尊进入香房，倘若需要，便具念正知，以右胁作片刻狮子卧。
\item 于是，恢复体力起来后，他在第二分中观察世间。第三分中，在他所依于居住的村镇，大众在饭前作了布施后，在饭后整理衣裳，持了香、花等，在寺庙集合。随后，世尊以适当的神变去到聚集的会众中，坐在法堂设好的最上佛座，开示适时、适量的法。然后，知时已到，即遣散会众，人们礼拜了世尊后离开。此即\textbf{饭后义务}。
\item 他如是完结了饭后义务，若欲擦拭肢体，即从佛座起，进入浴室，用侍者准备好的水，令肢体温暖。侍者则取了佛座，在香房的隔间设好。世尊著了善加染色的双层下衣，扎了腰带，在一肩穿了上衣，便回到那里坐下，独自宴坐片刻。然后，众比丘从各处赶来给侍世尊。于此，有些人提出问题，有些请求业处，有些请求闻法。世尊遂其所愿，即度过前夜。此即\textbf{前夜义务}。
\item 而在前夜义务的终了，当众比丘礼拜了世尊离开时，整个一万世界的天人获得机会，去到世尊处，提出如其准备、乃至四音节的问题。世尊解答这些天人的问题，即度过中夜。此即\textbf{中夜义务}。
\item 而在后夜分三部分，为舒缓从饭前起、因坐而受压迫的身体的疲惫，在一分以经行度过。在第二分，进入香房，具念正知，以右胁作狮子卧。在第三分起身而坐，为见在过去诸佛跟前以布施、持戒等曾作服务的补特伽罗，以佛眼观察世间。此即\textbf{后夜义务}。
\item 而在这天，世尊在王舍城完成了饭前义务，在饭后上路，在前夜为众比丘谈论了业处，在中夜为天人解答了问题，在后夜上了经行道经行，以一切知智听到五百比丘就这一切知智转起的谈论而了知。因此说：在作后夜义务时了知。
\item 且了知后，他便想：「这些比丘就我的一切知智谈论功德，而一切知智的作用对他们不明了，唯对我明了。当我离去后，他们将不断宣说自己的谈论，因此，我将对他们以此作为事件的发生，分别三种戒，于六十二处作不可逆转的狮子吼，汇集缘起的行相，令佛陀的功德明了，如撑起须弥卢一般，又如以金峰敲打天空一般，以阿罗汉为顶点，开示震动一万世界的梵网经，这开示即便在我般涅槃后五千年，仍将成为有情甘露大涅槃之导向。」如是思惟已，\textbf{便往圆亭而去}。\textbf{往}，即往此应去的方向。或者这是依格意义上的具格，此处之义即去到圆亭所在之处。
\item \textbf{便坐在设好的坐处}，据说在佛时，即便一个比丘所住之处，也总是设一佛座。为什么？据说，世尊对在自己跟前取了业处、住于安乐之处者作意道：「某人在我跟前取了业处而去，是否能转起殊胜呢？」然后，看到他舍了业处，起不善寻，因而想：「在像我这样的大师跟前取了业处而住，不善寻如何能征服这族姓子，令其轮回于无始的流转之苦？」为摄受他，即于此处显现自身，教诫了这族姓子，跃入空中，再回到自己的住处。于是，那些比丘如是受教诫已，想到「大师知晓我们的心意，来后只能站在我们附近显现自身，那时再以『尊者！请坐这里、请坐这里』去寻求坐处则成负担」，他们便设好坐处而住。若有椅子，便设椅子。若没有，便设床或木板、木头、岩石、沙堆。不得这些，则收集了枯叶，在上敷展粪扫衣而置。但在此处有国王的坐处，那些比丘便拍打设好而围绕，就世尊的倾向之智赞扬着功德而坐。就此而说「便坐在设好的坐处」。
\item 而如是坐已，他即便知晓，为发起谈论，仍问众比丘。他们便对其说了一切。因此说「\textbf{世尊坐下后}」等等。这里，\textbf{你们因何} \textit{kāya nuttha}，即你们因何\textbf{谈论共坐}之义。文本也作 kāya nettha，即「因何在此」之义，文本也作 kāya nottha，仍同前义。
\item \textbf{闲谈}，即业处的作意、总说、遍问等之间的另一谈论。\textbf{中断}，即缘于我的前来而未终结、未至顶点。以此显明什么？「我不是为了打断你们的谈论而来，我将以一切知完结你们的谈论，令其至于顶端」，刚坐下就作了一切知的邀请。
\item \textbf{此即我们中断的闲谈，然后世尊到来}，此处的旨趣为：尊者！我们的这场就世尊一切知智的功德论被中断，而非王论等的旁论，然后世尊到来，现在，请完结并开示它！
\item 至此，好比为了易于进入荷花、睡莲闪耀且清澈甜美之水的池塘，而用无垢的岩石排列、优雅装饰了珍宝的阶梯，津渡的淡黄地面铺满沙子，宛如表面散落了珍珠，又好比为了易于登上墙垣善加规划、围以彩色栏楯、仿佛欲触银河、体量雄伟的高贵楼阁，而用象牙制成、板材光滑柔软、金藤缠绕、因摩尼聚的光芒聚集闪耀而洁净的阶梯，又好比易于进入夹杂了金手镯、脚镯等碰撞的响声和欢声笑语、家人徘徊、显赫自在、财富闪耀的豪宅，而用有着令金银摩尼、珍珠珊瑚等光辉失色、辉煌稳固而优雅的门柱的大门，为了易于深入这部具足义文、彰显佛陀功德威力的契经，这由阿难尊者所说的由时间、地点、开示者、事由、会众、场景组成的因缘，其释义业已完成。\end{enumerate}

{\Large 5.} \textbf{诸比丘！\footnote{因缘部分既了，以下对佛陀的话便不再加引号，为便利计。}他人也许会说我的毁谤，或说法的毁谤，或说僧的毁谤，对此，你们不应嫌恨，不应郁闷，不应心怀不喜。诸比丘！他人也许会说我的毁谤，或说法的毁谤，或说僧的毁谤，对此，你们若受干扰，或不满意，这只会因此成为你们的障碍。诸比丘！他人也许会说我的毁谤，或说法的毁谤，或说僧的毁谤，对此，你们若受干扰，或不满意，你们还能了解他人的善说、恶说吗？}

\textbf{——不也，尊者！}

\textbf{诸比丘！他人也许会说我的毁谤，或说法的毁谤，或说僧的毁谤，对此，你们应从不实澄清不实：『如此则不实、如此则不如是，我们无有此、我们不存有此。』}

\begin{enumerate}\item 现在，轮到对世尊以「诸比丘！他人也许会说我的毁谤」等方法所安立的经的解释。而这即此经的解释。因为在考察经的安立后，所说便能明了，所以我们将首先考察经的安立。有四种经的安立：自身的意乐、他人的意乐、问题的主导、事件的发生。
\item 这里，若世尊非为他人祈请，仅仅以自身的意乐而说的经，如：若他希望经\footnote{若他希望经：即\textbf{中部}第 6 经。}、布经\footnote{布经：即\textbf{中部}第 7 经。}、大念处\footnote{大念处经：即\textbf{长部}第 22 经。}、大六入处分别经\footnote{六入处分别经：即\textbf{中部}第 137 经。}、圣种经\footnote{圣种经：即\textbf{增支部}第 4:28 经。}，正勤经系列，神足、根、力、觉支、道支经系列，如是等等，它们的安立为自身的意乐。
\item 若如\begin{quoting}能令罗睺罗解脱成熟之法已经成熟，我何不调伏罗睺罗至更高的漏尽？{\bluefootnotesize 〔相应部第 4.121 段〕}\end{quoting}考虑到他人的意乐、希望、心意、志向、可觉性，因他人的意乐而说，如：小教诫罗睺罗经\footnote{小教诫罗睺罗经：即\textbf{中部}第 147 经。}、大教诫罗睺罗经\footnote{大教诫罗睺罗经：即\textbf{中部}第 62 经。}、转法轮\footnote{转法轮经：即\textbf{相应部}第 56:11 经。}、界分别经\footnote{界分别经：即\textbf{中部}第 140 经。}，如是等等，它们的安立为他人的意乐。
\item 而四众、四种姓、龙、金翅鸟、乾达婆、阿修罗、夜叉、大王、三十三天等的诸天、大梵，如是等等，去到世尊处后，以「尊者！说有觉支、觉支」、「尊者！说有诸盖、诸盖」、「尊者！这些是五取蕴吗」、「什么是人在此世最上的财富」等方法提了问题。如是被提问的世尊所说的觉支相应\footnote{觉支相应：即\textbf{相应部}第 46 相应。}等，或其它天相应\footnote{天相应：即\textbf{相应部}第 1 相应。}、魔罗相应\footnote{魔罗相应：即\textbf{相应部}第 4 相应。}、梵相应\footnote{梵相应：即\textbf{相应部}第 6 相应。}、帝释所问\footnote{帝释所问经：即\textbf{长部}第 21 经。}、小毗陀罗\footnote{小毗陀罗经：即\textbf{中部}第 44 经。}、大毗陀罗\footnote{大毗陀罗经：即\textbf{中部}第 43 经。}、沙门果\footnote{沙门果经：即\textbf{长部}第 2 经。}、旷野\footnote{旷野经：即\textbf{经集}第 1:10 经。}、针毛粗毛经\footnote{针毛经：即\textbf{经集}第 2:5 经。}等，它们的安立为问题的主导。
\item 而凡是缘于已生起的事由所说的，如：法的继承者\footnote{法的继承者经：即\textbf{中部}第 3 经。}、小狮子吼\footnote{小狮子吼经：即\textbf{中部}第 11 经。}、月喻、子肉喻、木聚喻、火聚喻、泡团喻、波利质多树喻\footnote{月喻、子肉喻、泡团喻见\textbf{相应部}，木聚喻、火聚喻、波利质多树喻见\textbf{增支部}。}，如是等等，它们的安立为事件的发生。
\item 如是，在这四种安立中，此经的安立为事件的发生。由事件的发生，世尊安立此。什么事件的发生？于赞誉和毁谤。阿阇黎说三宝的毁谤，弟子则赞誉。如此，善巧于开示的世尊以此赞誉和毁谤为事件的发生，便开始以「诸比丘！他人也许会说我的毁谤」作开示。这里，\textbf{我}作属格，即「我的」之义。\textbf{或}字为可选之义。\textbf{他人}，即对立的有情。\textbf{对此}，即对那些说毁谤者。
\item 以「不应嫌恨」等，虽然这些比丘们并无嫌恨，但为了防止将来的族姓子们在这样的场合生起不善，而设置法的准绳。这里，触及心为\textbf{嫌恨}，即忿恨的同义语。已经不高兴，因此不满足、不喜悦，为\textbf{郁闷}，即忧的同义语。既不满足自己，也不满足他人的利益，为\textbf{不喜}，即忿恨的同义语。如是，此中以二词说行蕴，以一说受蕴，说了二蕴。依于这些，其余相应法的出现亦被拒斥。
\item 如是，以第一法遮止了意的嗔恚，为显明此处的过患，便以第二法说了「对此，你们若受干扰，或不满意，这只会因此成为你们的障碍」。这里，\textbf{对此，你们若……}，即「你们若对这些说毁谤者，或对这毁谤，你们如果……」之义。以忿怒\textbf{受干扰}，以忧而\textbf{不满意}，\textbf{这只会因此成为你们的障碍}，即这只会对你们，因此忿怒与此不满意，成为初禅等的障碍。
\item 如是，以第二法显明了过患，为显明甚至不能了解语义，便以第三法说了「你们还能了解他人的……」等。这里，\textbf{他人}，即任何人。因为受干扰者不能了解佛、辟支佛、圣弟子、父母、怨敌的善说、恶说之义。如说：\begin{quoting}忿怒者不知义利，忿怒者不得见法，\\若忿怒征服了人，此时便黑暗盲目。\\忿怒是非义之母，忿怒是心的动荡，\\人尚不能理解，这由内而生的怖畏。{\bluefootnotesize 〔增支部第 7.64 段〕}\end{quoting}
\item 如是，在一切处制止了对毁谤的意的嗔恚，现在，为显明应行的行相，便说了「对此，你们应从不实澄清不实」等。这里，\textbf{对此，你们}，即你们对此毁谤。\textbf{应从不实澄清不实}，即对于不实，应以不实之相驱除之。如何？即以「如此则不实……」等方法。此处的连结为：当听到「你们的大师非一切知，法是恶说，僧是恶行道」等，不应默然，而应当说：\textbf{如此则不实}，凡你们所说的，以此原由不实，以此原由\textbf{不如是，我们无有此、我们不存有此}，我们的大师就是一切知，法是善说，僧是善行道，对此，有这般这般原由。且此中，当知第二句、第四句为第一句、第三句的同义异语。且此澄清应只对毁谤而作，非一切处。因为如果当被说到「你是恶戒，你的阿阇黎恶戒，这般这般由你所作，由你的阿阇黎所作」时，应默然承受，保持怀疑。所以，应不起意的嗔恚，澄清毁谤。而对于以「你是骆驼、你是牛」等方法以十种骂詈事谩骂的人，在此也应超然视之，承受、忍辱。\end{enumerate}

{\Large 6.} \textbf{诸比丘！他人也许会说我的赞誉，或说法的赞誉，或说僧的赞誉，对此，你们不应欢喜，不应喜悦，不应心中飘然。}

\textbf{诸比丘！他人也许会说我的赞誉，或说法的赞誉，或说僧的赞誉，对此，你们若欢喜、喜悦、飘然，这只会因此成为你们的障碍。}

\textbf{诸比丘！他人也许会说我的赞誉，或说法的赞誉，或说僧的赞誉，对此，你们应从实肯定实：『如此则实、如此则如是，我们有此、我们存有此。』}

\begin{enumerate}\item 如是显明了在毁谤处的如如之相，现在，为显明在赞誉处，说了「诸比丘！他人也许会说我的赞誉」等。这里，\textbf{他人}，即任何净信的人天。以之而欢喜者为\textbf{欢喜}，即喜的同义语。悦意的状态为\textbf{喜悦}，即心之乐的同义语。得意的状态为\textbf{飘然}。什么的飘然？\textbf{心}的。即带来掉举、令起得意之喜的同义语。此处也以二词说行蕴，以一说受蕴。
\item 如是，以第一法遮止了飘然，为显明此处的过患，便以第二法说了「对此，你们若……」等。此处，亦当知「\textbf{这只会因此成为你们的障碍}」，即这只会对你们，因此飘然，成为初禅等的障碍。为什么这么说？世尊不是以\begin{quoting}称扬佛者在身中生起了喜，\\这喜甚至胜过整个阎浮提。\end{quoting}\begin{quoting}称扬法者在身中生起了喜，\\这喜甚至胜过整个阎浮提。\end{quoting}\begin{quoting}称扬僧者在身中生起了喜，\\这喜甚至胜过整个阎浮提。\end{quoting}\begin{quoting}诸比丘！于佛陀净喜者，于最上净喜。\end{quoting}如是等数百经赞誉对三宝的喜悦吗？赞誉是真实的，但这是依于出离，而此处的意趣则以「我们的佛、我们的法」等方法，与车匿尊者生起的依于俗家的喜悦相似。而这对获得禅那等构成障碍。因此，车匿尊者在佛陀未般涅槃前不能转起殊胜，在般涅槃时，受施设的梵罚威吓，舍弃了这喜悦，便转起殊胜。所以，当知就构成障碍而说此。因为这喜与贪俱行，而贪仍与忿怒相似。如说：\begin{quoting}贪婪者不知义利，贪婪者不得见法，\\若贪婪征服了人，此时便黑暗盲目。\\贪婪是非义之母，贪婪是心的动荡，\\人尚不能理解，这由内而生的怖畏。{\bluefootnotesize 〔如是语第 88 经〕}\end{quoting}
\item 而第三番，在此处虽未出现，当知从意义上也已出现。因为正如同忿怒者，如是贪婪者也不能知义。\footnote{第三番：即上文所说的「……你们还能了解他人的善说、恶说吗」。}
\item 在显明应行的行相番，其连结为：当听到「你们的大师是一切知、阿罗汉、正等正觉者，法是善说，僧是善行道」等，不应默然，而应如是肯定：\textbf{如此则实}，凡你们所说的，以此原由则实，以此原由\textbf{如是}——因为彼世尊亦即阿罗汉、正等正觉者，法即善说、自见，僧即善行道、正直行道。当被问到「你具戒吗」，如果具戒，则应肯定「我是具戒者」。如果被问到「你是得初禅者……阿罗汉吗」，则应只对俱同分的比丘肯定。因为如是既可避免恶欲，又能阐明教法不虚。其余当知仍如已述。\end{enumerate}

\section{小戒}

{\Large 7.} \textbf{诸比丘！这说着如来赞誉的凡夫所能说的，只是小量、浅量的戒量。诸比丘！什么是这说着如来赞誉的凡夫所能说的小量、浅量的戒量？}

\begin{enumerate}\item 「诸比丘！……只是小量」之适用为何？此经与二句相关：赞誉和毁谤。这里，毁谤，即「如此则不实、如此则不如是」，如火刚触及水边即退。而赞誉，即「应从实肯定实——如此则实」，仍如是继续。且它有二种：梵赐所说的赞誉，和比丘僧团以「希有！朋友」等方法开始的赞誉。其中，比丘僧团所说赞誉之适用，将在之后的「空性阐释」中显明。而此处，为显明梵赐所说赞誉之适用，以「诸比丘！……只是小量」开始开示。
\item 这里，\textbf{小量}，即有限之名，\textbf{浅量}只是其异名。度量被称为量，其量为小即小量，其量为浅即浅量。戒即\textbf{戒量}。这即是说：诸比丘！\textbf{凡夫}即便鼓起热情——我要说如来的赞誉，\textbf{说着赞誉所能说的}，只是小量、浅量的戒量。
\item 这里，设问：这戒难道不是修行者最上的严饰吗？如古人们说：\begin{quoting}戒是修行者的庄严，戒是修行者的装饰，\\以诸戒庄严的修行者，得至装饰的最上。\end{quoting}世尊也在数百经中以戒为大而作谈论。如说：\begin{quoting}诸比丘！若比丘希望「愿我为同梵行者喜爱、悦意、尊重、尊敬」，他唯应圆满诸戒。{\bluefootnotesize 〔中部第 1.65 段〕}\end{quoting}以及\begin{quoting}如松鸦之于蛋，如牦牛之于尾，\\如对爱子，如对独眼，\\亦应如是守护戒，\\始终善加检点、带着尊重。\end{quoting}以及\begin{quoting}花香不逆风而行，旃檀、山马茶、茉莉亦否，\\但善人之香逆风而行，善人飘扬到四方。\end{quoting}\begin{quoting}旃檀、山马茶，抑或青莲及茉莉花，\\在这些香类之中，戒香为无上。\end{quoting}\begin{quoting}此香为少量：即此山马茶、旃檀，\\而具戒者之香最上，飘扬至天间。\end{quoting}\begin{quoting}对具足戒者、住不放逸者、\\正智解脱者，魔不识其途。{\bluefootnotesize 〔法句第 54-57 颂〕}\end{quoting}以及\begin{quoting}俱慧之人既住于戒，修习着心与慧，\\热忱、贤明的比丘，他能解开此结。{\bluefootnotesize 〔相应部第 1.23 段〕}\end{quoting}以及\begin{quoting}好比，诸比丘！任何至于增长、成长、广大的种子、生物，它们全都依于地，住立于地，如是，这些种子、生物得至增长、成长、广大。如是，诸比丘！比丘依于戒，住立于戒，当修习七觉支、多作七觉支时，得以圆满诸法的增长、成长、广大。{\bluefootnotesize 〔相应部第 5.150 段〕}\end{quoting}还有其它许多经可见。如是在数百经中唯以戒为大而作谈论，为什么在此处说其为小量？由比较更上的功德故。因为戒不能圆满定，定不能圆满慧，所以由比较最上而名为较下、浅量。
\item 如何戒不能圆满定？因为世尊在从现等觉起的第七年，在舍卫城门外的茎芒果树下，坐在十二由旬宝亭内的量一由旬的宝座上，于被撑起的三由旬的天白伞下，为对十二由旬的会众彰显自身之建立、摧破外道，显示以「从上身转起火聚，从下身转起水流……从一一毛孔转起火聚，从一一毛孔转起水流，六色……」等方法转起的双神变。从其金色身中放出金色的光芒，直趋有顶，整个一万轮围仿佛到了庄严的时刻。
\item 后后的光芒好似与前前成双成对，在同一刹那转起。而无有二心在同一刹那转起。然而对诸佛世尊，以有分停留之轻快性，与五行相之宿习自在性，它们好似在同一刹那转起。虽然各个光线的转向、遍作、决意仍是各别的。因为为了青色光芒，世尊入青遍，为黄色光芒而入黄遍，为赤、白色光芒而入赤、白遍，为火聚入火遍，为水流则入水遍。而「大师经行，所化现者或站、或坐、或躺卧」等的一切，当需详述。此中，无一戒的作用，一切唯是定的作用。如是，戒不能圆满定。
\item 世尊既圆满了四阿僧祇又十万劫的波罗蜜，在二十九岁时，从有着转轮王光辉的居处出离，在无劣河畔出家，行了六年的精勤修行后，当毗舍佉的满月日，受用了优楼频螺村中善生所施的已投入天之食素的乳糜，晡时沿南北进入菩提座，右绕菩提树王三匝，立于东北方，敷展草垫，结好三连接之跏趺，以四支具足的慈业处为前导，决意于精进，（他）在十四肘的高贵床座上，背后五十肘高的菩提树干好似置于金椅上的银杆，上方撑起的菩提树枝好似摩尼伞，当金色的衣上落满如同珊瑚的菩提嫩芽、日色将暮时，驱赶了魔军，初夜忆念了宿住，中夜净化了天眼，在黎明时，落实了一切诸佛宿习的缘相之智，转起入出息的四禅，即以此为基础，增长了毗婆舍那，依道的次第，以证得的第四道摧毁了一切烦恼，便通达了一切诸佛的功德，此即慧的作用。如是，定不能圆满慧。
\item 这里，好比手中之水不能满碟中之水，碟中之水不能满瓶中之水，瓶中之水不能满瓮中之水，瓮中之水不能满罐中之水，罐中之水不能满大釜之水，大釜之水不能满小池之水，小池之水不能满沟渠之水，沟渠之水不能满小川之水，小川之水不能满五大河之水，五大河之水不能满轮围大海之水，轮围大海之水不能满须弥卢山脚下大海之水。较之碟中之水，手中之水为有限……较之须弥卢山脚下大海之水，轮围大海之水为有限，如此，较之上上之大水，下下之水为有限。如是，较之上上之功德，下下之戒当知为小量、浅量。因此说「诸比丘！这……只是小量、浅量的戒量」。
\item 在「\textbf{凡夫所……}」中，\begin{quoting}有二凡夫为佛、日种所说：\\一为盲凡夫，一为善凡夫。\end{quoting}这里，若于蕴、界、处等无有学习、遍问、听闻、忆持、省察，此即盲凡夫。若有此等，斯即善凡夫。
\item 此另有二种：\begin{quoting}由众多之生起等原因故而为凡夫，\footnote{这里是将「凡夫 \textit{puthujjana}」拆开为「众多 \textit{puthu}、生起 \textit{janana}」两词来解释。}\\由隶属于凡夫故，以其人不同故。\footnote{义注这样说，是因为 puthu 兼有「不同、差异」之意。}\end{quoting}因为他由众多种种品类的烦恼等之生起等原因故，而为凡夫。如说：\begin{quoting}以生起众多烦恼为凡夫，未破除众多有身见者为凡夫，以仰众多大师之口为凡夫，以未从众多的一切趣出起为凡夫，以造作众多种种行作为凡夫，以被众多种种暴流冲掠、众多热火燃烧、众多热恼焦灼，在众多种种五欲中染著、贪求、交结、沉迷、沉瘾、固著、阻碍、障碍，而为凡夫，以被众多五盖所覆、所盖、所覆障、所闭、所掩、所包裹，而为凡夫。{\bluefootnotesize 〔大义释第 51 段〕}\end{quoting}且由隶属于众多超出算数、背对圣法、践行劣法之人故而为凡夫，其唯以不同、个别而得称，不与相应于戒、闻等功德的圣人交际，即为凡夫。
\item \textbf{如来}，以八种原因，世尊为如来：以如是而来为如来，以如是而行为如来，以来至如实之相为如来，以确切地等觉如实之法为如来，以如实之见为如来，以如实之语为如来，以如是之行为如来，以征服之义为如来。
\item 世尊如何\textbf{以如是而来为如来}？好比为了一切世间的利益而兴起热情的先前的众正等正觉者而来，好比毗婆尸世尊、尸弃世尊、毗舍婆世尊、拘留孙世尊、拘那含世尊、迦叶世尊而来。这说的是什么？这些世尊以此志向而来，我们的世尊也以此而来。
\item 或者，好比毗婆尸世尊……迦叶世尊圆满了布施波罗蜜，圆满了戒、出离、慧、精进、忍辱、真实、决意、慈、舍波罗蜜，圆满了这十波罗蜜、十随波罗蜜、十第一义波罗蜜等整三十波罗蜜，遍舍了肢体，遍舍了眼、财产、王位、妻儿等这五大遍舍，圆满了宿行、前行、法的言谈知义等诸行，证得了觉悟行的顶点而来，我们的世尊也如是而来。
\item 或者，好比毗婆尸世尊……迦叶世尊修习、增长了四念处、四正勤、四神足、五根、五力、七觉支、八支圣道而来，我们的世尊也如是而来。如是，以如是而来为如来。\begin{quoting}正好比在此世间，毗婆尸等牟尼来至于一切知性，\\这释迦牟尼也如是而来，具眼者因此被称为如来。\end{quoting}
\item 如何\textbf{以如是而行为如来}？好比刚出生的毗婆尸世尊……迦叶世尊而行。那彼世尊又如何行？即他刚出生，便以平等的双足站立于地，面朝北方，以交错的七步而行。如说：\begin{quoting}阿难！刚出生的菩萨以平等的双足站立，面朝北方，以交错的七步而行，在被撑起的白伞下环顾四方，说牛王之语——我是世间最上，我是世间最尊，我是世间最胜，此是最后一生，现在已无再有。{\bluefootnotesize 〔长部第 2.31 段〕}\end{quoting}且其行如实、不虚妄，以其为诸多殊胜证得之征兆。
\item 即他刚出生便以平等的双足站立，这是他获得四神足的征兆。而面朝北方，为一切世间最上\footnote{北方与最上的巴利文相同，即 uttara。}的征兆。七步交错，为获得七觉支之宝的。而「金柄拂尘上下翻飞」中所说的拂尘的扬起，为碾压一切外道的。撑起白伞，为获得最上阿罗汉解脱之无垢白伞的。七步过后站定环顾四方，为获得一切知无碍智的。述说牛王之语，为转起不可逆转的无上法轮的征兆。
\item 彼世尊亦如是而行，且其行如实、不虚妄，以其亦为彼等殊胜证得之征兆。因此古人们说：\begin{quoting}甫一出生，便如牛王，\\以双足触地，\\这乔达摩便跨了七步，\\众风神随而撑起白伞。\\这乔达摩行完了七步，\\便环顾周遭四方，\\发出八支具足的话语，\\如狮子站在高山之上。\end{quoting}如是，以如是而行为如来。
\item 或者，好比毗婆尸世尊……迦叶世尊，此世尊亦如是以出离舍弃欲贪而行，以无嗔舍弃嗔，以光明想舍弃昏沉睡眠，以不散乱舍弃掉举，以法之确定舍弃疑，以智破碎无明，以愉悦除遣不喜，以初禅开启（五）盖之窗，以二禅平息寻伺，以三禅脱离喜，以四禅舍弃苦乐，以空无边处定超越色想、有对想、种种想，以识无边处定超越空无边处想，以无所有处定超越识无边处想，以非想非非想处定超越无所有处想而行。
\item 以无常随观舍弃常想，以苦随观舍弃乐想，以无我随观舍弃我想，以厌离随观舍弃喜，以离贪随观舍弃贪，以灭随观舍弃集，以舍遣随观舍弃执取，以尽随观舍弃密集想，以衰减随观舍弃追逐，以变易随观舍弃坚牢想，以无相随观舍弃相，以无愿随观舍弃愿，以空性随观舍弃执著，以增上慧法观舍弃对执取坚实的执著，以如实知见舍弃对迷妄的执著，以过患随观舍弃对阿赖耶的执著，以省察随观舍弃无省察，以还灭随观舍弃对连接的执著，以须陀洹道粉碎与见一处的烦恼，以斯陀含道舍弃粗重的烦恼，以阿那含道根除与微细俱行的烦恼，以阿罗汉道彻底斩断一切烦恼而行。如是，亦以如是而行为如来。\footnote{以上两段略同\textbf{清净道论}·说戒品第 140 段。}
\item 如何\textbf{以来至如实之相为如来}？地界以坚硬性之相而如实不虚，水界以滴淌之相，火界以热性之相，风界以支持之相，空界以不可触之相，识界以了知之相，色以变坏之相，受以经验之相，想以识别之相，诸行以行作之相，识以了知之相，寻以运用之相，伺以细思之相，喜以遍满之相，乐以幸福之相，心一境性以不散乱之相，触以触摸之相，信根以胜解之相，精进根以策励之相，念根以照顾之相，定根以不散乱之相，慧根以识知之相，信力以在不可信处不可动摇之相，精进力以在懈怠处，念力以在忘失处，定力以掉举处，慧力以在无明处不可动摇之相，念觉支以照顾之相，择法觉支以简择之相，精进觉支以策励之相，喜觉支以遍满之相，轻安觉支以平息之相，定觉支以不散乱之相，舍觉支以省察之相，正见以知见之相，正思惟以运用之相，正语以彻晓之相，正业以等起之相，正命以洁净之相，正精进以策励之相，正念以照顾之相，正定以不散乱之相，无明以无知之相，诸行以思之相，识以了知之相，名以倾向之相，色以变坏之相，六入处以领域之相，触以触摸之相，受以经验之相，爱以因之相，取以把握之相，有以追逐之相，生以发生之相，老以衰老之相，死以亡殁之相，诸界以空性之相，诸处以领域之相，念处以照顾之相，正勤以精勤之相，神足以成就之相，根以增上之相，力以不可动摇之相，觉支以出离之相，道以因之相，谛以如实之相，奢摩他以不散乱之相，毗婆舍那以随观之相，止观以一味之相，双运以不违越之相，戒清净以防护之相，心清净以不散乱之相，见清净以知见之相，灭智以正断之相，无生智以安息之相，欲以根本为相，作意以令起为相，触以汇集为相，受以汇合为相，定以首领为相，念以主导为相，慧以较之更上为相，解脱以坚实为相，隶属不死的涅槃以终了为相而如实不虚。如是，以智趣来至如实之相，无所错失而来、而到达，即为如来。如是，以来至如实之相为如来。
\item 如何\textbf{以确切地等觉如实之法为如来}？如实之法者，即四圣谛。如说：\begin{quoting}诸比丘！这四者如实、不虚、不异。四者为何？「此是苦」，诸比丘！这为如实，这为不虚，这为不异……{\bluefootnotesize 〔相应部第 5.1090 段〕}\end{quoting}所详述。而这些为世尊所等觉，所以，由等觉如实之法故，被称为如来。此中的「来 \textit{gata}」一词，即等觉之义。
\item 另外，对于老死，以生为缘生成、兴起之义如实、不虚、不异……对于诸行，以无明为缘生成、兴起之义如实、不虚、不异，同样，无明为诸行的缘之义、诸行为识的缘之义……生为老死的缘之义如实、不虚、不异。世尊已等觉彼等一切，所以，亦由等觉如实之法故，被称为如来。如是，以确切地等觉如实之法为如来。
\item 如何\textbf{以如实之见为如来}？对在俱有天……天人的人世间中，但凡存在来至无量世界中无量有情眼门领域的色所缘，世尊从一切行相能知能见。且如是知见时，他或以可意、不可意等，或以于所见、所闻、所觉、所知中所得之句——\begin{quoting}什么色为色处？若四大种所造之色，俱有光泽，有见有对，青、黄……{\bluefootnotesize 〔法集论第 616 段〕}\end{quoting}等方法，以众多名称、十三番、五十二法\footnote{这些均详见\textbf{法集论}·色章。}所分别者，唯是如实，无有虚妄。来至耳门等领域的声等亦同。且世尊说：\begin{quoting}诸比丘！这俱有天……天人的人世间的所见、所闻、所觉、所知、所得、所遍求、意之所随伺者，我皆知之。我既证知此，对此如来之所知，如来便不顾念。{\bluefootnotesize 〔增支部第 4.24 段〕}\end{quoting}如是，以如实之见为如来。这里，当知在「如实之见」之义上，而有「如来」一词的生成。
\item 如何\textbf{以如实之语为如来}？在世尊以不可战胜之跏趺坐于菩提座、粉碎了三魔罗\footnote{魔罗有五种，即蕴、烦恼、行作、死、天子魔罗，这里的三魔罗当指「烦恼、行作、天子」三者。}之头首、等觉无上正等正觉之夜，与在娑罗双树间以无取涅槃界般涅槃之夜，两者之间的四十五年期间，无论在觉悟初期、中期、后期由世尊所说的契经、应颂……毗陀罗，这一切在义、文上无过，无减无增，一切行相圆满，磨碎贪之㤭慢，磨碎嗔、痴之㤭慢。其间连毫尖之量的差错也无，一切如以一符契所标记，如以一升斗所量取，如以一天枰所称重，唯是如实、不虚、不异。因此说：\begin{quoting}均陀！在如来等觉无上正等正觉之夜，与以无余取涅槃界般涅槃之夜，两者之间所说、所谈、所说明的一切，唯是如实、不异，所以被称为「如来」。{\bluefootnotesize 〔增支部第 4.23 段〕}\end{quoting}此中的「来 \textit{gata}」一词，即「说 \textit{gada}」之义。如是，以如实之语为如来。
\item 另外，言语 \textit{āgada} 即言说 \textit{vacana} 之义，以其言语无有三种颠倒\footnote{三种颠倒：即想、心、见的颠倒，见\textbf{增支部}第 4.49 段。}，将 da 字作 ta 字，便为「如来 \textit{Tathāgata}」，如是当知基于此义而有词的成立。
\item 如何\textbf{以如是之行为如来}？因为世尊的身随顺语，语亦随顺身，所以行与语同，且语与行同。且对如是作为者，语如是，身亦如是而来、而转起，身如是，语亦如是而来、而转起，即如来之义。因此即说：\begin{quoting}诸比丘！如来行与语同，语与行同。以此行与语同、语与行同，所以被称为「如来」。{\bluefootnotesize 〔增支部第 4.23 段〕}\end{quoting}如是，以如是之行为如来。
\item 如何\textbf{以征服之义为如来}？以上至有顶、下至无间为界，在旁及无量的世界中，以戒、定、慧、解脱、解脱知见征服一切有情，而无有与其等量者，无等、不可量、无上、王中胜王、天中之天、帝释中之胜帝释、梵天中之胜梵天。因此说：\begin{quoting}诸比丘！在俱有天……天人的人世间中，如来征服，不被征服，是仅见者、自在者，所以被称为「如来」。\end{quoting}
\item 此处，当知词的成立如是。Agada 作药解。那这（药）是什么？优雅所成之开示与积聚之福德。因为这大威力的医生以此征服一切其它论点和俱有天的世间，如以天药征服蛇一般。当如此征服一切世间时，如实、不颠倒的优雅所成之开示与积聚之福德便即是药。将 da 字作 ta 字，当知即「如来 \textit{Tathāgata}」。如是，以征服之义为如来。
\item 另外，以如实而来亦为如来，如实地来亦为如来。来，即理解、超越、证得、行道之义。这里，以度遍知之如实而来理解整个世间为如来。以断遍知之如实而来超越世间之集为如来。以证得之如实而来证得世间之灭为如来。对趣向世间灭之行道，如实地来行道为如来。因此世尊说：\begin{quoting}世间，诸比丘！为如来所等觉，如来不与世间相应。世间之集，诸比丘！为如来所等觉，如来已断世间之集。世间之灭，诸比丘！为如来所等觉，如来已证世间之灭。趣向世间灭之行道，诸比丘！为如来所等觉，如来已修趣向世间灭之行道。诸比丘！这俱有天……的世间等的一切，都为如来所等觉，所以被称为「如来」。{\bluefootnotesize 〔增支部第 4.23 段〕}\end{quoting}其义当知亦为如是。且这在阐明如来的如来之相中只是入门，而唯有如来能以一切行相解释如来的如来之相。
\item \textbf{诸比丘！什么是这……}，即是问：以此小量、浅量的戒量说着如来赞誉的凡夫所能说的，是什么呢？这里，问有五种：显明未见之问、交流所见之问、消除疑虑之问、赞同之问、欲说之问。
\item 这里，什么是\textbf{显明未见之问}？其相本来未知、未见、未量、未度、未明了、未阐明，为了知、见、量、度、明了、阐明而问问题，此即显明未见之问。
\item 什么是\textbf{交流所见之问}？其相本来已知、已见、已量、已度、已明了、已阐明，为了与其他智者一起交流而问问题，此即交流所见之问。
\item 什么是\textbf{消除疑虑之问}？本来陷入疑惑、陷入疑虑、生起两可——「如是？或否？是何？如何？」他为了消除疑虑而问问题，此即消除疑虑之问。
\item 什么是\textbf{赞同之问}？世尊为了比丘们的赞同而问问题——\begin{quoting}你们认为如何，诸比丘！色是常还是无常？\\——无常，尊者！\\若无常者，是苦还是乐？\\——是苦，尊者！{\bluefootnotesize 〔大品第 21 段〕}\end{quoting}等一切当说，此即赞同之问。
\item 什么是\textbf{欲说之问}？世尊为了欲对比丘们说而问问题——\begin{quoting}诸比丘！有此四念处，何等为四？……有此八支道，何等为八？\end{quoting}此即欲说之问。
\item 如此，在这五种提问中，首先，由无有任何未见之法故，如来无显明未见之问。由不会生起「对此，在与其他沙门婆罗门之智者一起交流后，我再开示」般的存念故，亦无交流所见之问。而因为诸佛连一法的疑念、犹豫也无，已在菩提座上断除一切疑惑，所以亦无消除疑虑之问。诸佛有剩余的二问，此即其中的欲说之问。\end{enumerate}

{\Large 8.} \textbf{『舍弃了杀生，戒离杀生，沙门乔达摩放下棍杖、放下刀剑，有耻，俱有怜悯，哀愍一切群生而住』，诸比丘！说着如来赞誉的凡夫或如此说。}

\textbf{『舍弃了不与取，戒离不与取，沙门乔达摩取所与、期待所与，以不盗、成为清净之自我而住』，诸比丘！说着如来赞誉的凡夫或如此说。}

\textbf{『舍弃了非梵行，作梵行，沙门乔达摩作远离行，戒离媾和、村法』，诸比丘！说着如来赞誉的凡夫或如此说。}

\begin{enumerate}\item 现在，为说此欲说之问的所问之义，便说了「舍弃了杀生」等。
\item 这里，杀害生类为\textbf{杀生}，即是说杀戮生类、毁灭生类。此中的「生类」即俗称之有情、第一义之命根，而对此生类作生类想者，为断绝命根发起准备，转起身语门中的任一门的杀戮之思，即杀生。对于无有功德的畜生等生类，对小生类为小罪，对大躯体为大罪。为什么？由加行大故，当加行相等，亦由依处大故。对于俱有功德的人类等，对功德小的生类为小罪，对大功德者为大罪。而当身躯、功德相等时，当知对烦恼和准备均柔和者为小罪，对粗重者为大罪。
\item 其要素有五：生类、作生类想者、杀心、准备、因此而死。实施有六：亲手、命令、放舍、静置、明咒所成、神变所成。而此义若加详述，将极冗繁，所以我们不再详述，它处亦如此。有需求者可查看律藏义注「一切善见」把握。
\item \textbf{舍弃}，即舍弃了这被称为杀生之思的恶戒。\textbf{戒离}，即从舍弃之时起，便停止、戒离此恶戒。对其尚无以「我将违犯」的眼、耳所了知之法，遑论以身？在其它类似等句中，其义亦当以此法而知。
\item \textbf{沙门}，即世尊因止恶所得之称。\textbf{乔达摩}则因族姓。且不仅世尊戒离杀生，比丘僧团亦戒离——尽管开示从一开始如是出现，但当阐明意义时，以比丘僧团来阐明也适当\footnote{义注在这里区分了「开示 \textit{desanā}」和「意义 \textit{attha}」。开示，即相当于文本。}。
\item \textbf{放下棍杖、放下刀剑}，意即不存在为了伤害他人拿起棍杖或刀剑，即为放下棍杖、放下刀剑。且此中，除了棍杖，其余所有的器具，由恼害有情故，当知为「刀剑\footnote{这是从语源来解释「刀剑 \textit{sattha}」，即来自有情 \textit{satta}、恼害 \textit{viheṭhana} 二词。}」。而当比丘拿起拐杖、齿木或剪刀而行时，并非为了伤害他人，所以仍可称为「放下棍杖、放下刀剑」。
\item \textbf{有耻}，即具有以对恶的嫌厌为相的耻。\textbf{俱有怜悯}，即俱有怜悯之慈心。\textbf{哀愍一切群生}，即以利益怜悯一切群生，意即以此所俱之怜悯性，对一切群生心怀利益。\textbf{住}，即举止、存活、存续、守护。\textbf{诸比丘……或如此}，即「诸比丘……或如是」。「或」字是考虑到后文的「舍弃了不与取」等而说的可选之义，如是在一切处，当知为考虑到前后文的可选性。
\item 而此中的略义为：诸比丘！说着如来赞誉的凡夫会如是说——「沙门乔达摩不杀生类，不教杀，对此不许可，戒离此恶戒，哎！咄！佛陀的功德实大！」在鼓起大热情后欲说赞誉，也只能说小量、浅量的正行之戒量，不能就更上的不共性而说赞誉。且何止凡夫，连须陀洹、斯陀含、阿那含、阿罗汉、辟支佛也不能够，唯如来能够，我将在下文对你们说。这即此中连同旨趣的释义。此后我们将只解释先前未及的词句。
\item 在「舍弃了不与取」中，取不与者即\textbf{不与取}，即是说收集他物、盗窃、劫掠。这里，所不与者，即他人之所有，但当他人随意处置时，则不应受罚且无罪。对此他人之所有作「他人之所有」想者，为取其而发起准备的盗窃之思，即不与取。当他人的财产低贱时为小罪，高贵时为大罪。为什么？由依处的高贵故。当依处相等时，对德高者的财产依处为大罪，相较于彼彼德高者，对较之德低者的财产依处为小罪。
\item 其要素有五：他人之所有、作他人之所有想者、盗心、准备、因此而取。实施有六：即亲手等。彼等又依次以「盗窃窃取、征服窃取、覆蔽窃取、遍计窃取、孤沙窃取\footnote{孤沙窃取：\textbf{律注}说「更换孤沙（草之标识）的窃取被称为孤沙窃取 \textit{kusaṃ saṅkāmetvā pana avaharaṇaṃ kusāvahāro ti vuccati}」。}」等窃取而转起。这是略说，而详说如「一切善见」中所述。
\item 唯取所与者，即\textbf{取所与}。心中唯期待所与，即\textbf{期待所与}。盗窃为盗，不以盗即\textbf{不盗}。正由不盗而\textbf{成为清净}。\textbf{自我}，即自体。即是说，作为不盗、成为清净之自我而住。其余仍以在第一学处中所说之法连接。且如此处，在一切处均如是。
\item \textbf{非梵行}，即非胜行。作梵、胜之正行为\textbf{作梵行}。\textbf{作远离行}，即作远离非梵行之行。\textbf{媾和}，即由贪染现行相似故，得称为「媾和者」，由（媾和者）所受用故，得名「媾和」的非善法。\textbf{村法}，即村民之法。\end{enumerate}

{\Large 9.} \textbf{『舍弃了妄语，戒离妄语，沙门乔达摩作真实语，致力真实，坚牢，可信，不欺骗世间』，诸比丘！说着如来赞誉的凡夫或如此说。}

\textbf{『舍弃了离间语，戒离离间语，沙门乔达摩在此处听闻后，不会为了分裂此等在彼处说，或在彼处听闻后，为了分裂彼等对此等说。如此，或调停分裂，或增进团结，乐于和合，喜于和合，欢喜和合，说令和合之语』，诸比丘！说着如来赞誉的凡夫或如此说。}

\textbf{『舍弃了粗恶语，戒离粗恶语，沙门乔达摩说凡是无瑕、悦耳、可亲、入心、礼貌、众人喜爱、众人悦意这般的话语』，诸比丘！说着如来赞誉的凡夫或如此说。}

\textbf{『舍弃了绮语，戒离绮语，沙门乔达摩是时语者、实语者、义语者、法语者、律语者，适时地说具有价值、有着理由、具有边界、伴有义利之语』，诸比丘！说着如来赞誉的凡夫或如此说。}

\begin{enumerate}\item 在「舍弃了妄语」中，\textbf{妄}，即预设欺骗的破坏义利的语行或身行。而以欺骗之意趣，发起欺骗他人的身语行之思，为\textbf{妄语}。另一法：\textbf{妄}，即不实、不如此之事。\textbf{语}，即将其令人了知作如实、如此。但从相上来说，将不如此之事欲令他人了知作如此，发起如是之表的思，即妄语。视其所破坏的义利，小为小罪，大为大罪。
\item 另外，对在家人，因不欲布施自己的财产，而以「没有」等方法所转起者为小罪，作为见证，为了破坏义利之所说为大罪。对出家人，得了少许的油或酥，以嬉笑之意趣、夸大之辞——我说，今天村中的油如河水流淌——所转起者为小罪，虽未见，却以「已见」等方法所说者为大罪。
\item 其要素有四：不如此之事、欺骗之心、为此的努力、他人了知此义。实施唯一，即亲手。它可被理解为以身，或以依赖于身，或以语，作出欺骗他人的行为。若他人因此行为了知其义，在发起行为之思的刹那，他便缚于妄语之业。而因为如同以身、依赖于身、语欺骗他人，同样，以「你对他作此说」等命令也可，写在纸上放舍面前也可，以「此义作如是观」等写在墙上静置也可，所以在此，命令、放舍、静置等实施也合适，但由在义注中未提及故，应在考察后把握\footnote{命令、放舍、静置，即「杀生」中提到的六种实施中的三种。}。
\item 说真实者，即\textbf{作真实语}。以真实与真实相连、连接为\textbf{致力真实}，意即不夹杂地说妄语。因为若人有时妄语，有时实语，由其以妄语夹杂故，不以真实连接真实，所以非致力真实者。而他并非如此，即便因性命也不妄语，唯以真实与真实相连，为致力真实。\textbf{坚牢}，即坚定，意即言语坚定。因为有人好似姜黄所染，好似埋于稃堆的树桩，如置于马背的葫芦，非言语坚定者。有人则好似石刻，好似因陀柱，为言语坚定者，即便以剑砍头，曾无二辞，此即坚牢。
\item \textbf{可信}，即可信赖，意即可相信。因为有人不可信，当说「这是某某所说」时，势必有「别相信他的话」。有人可信，当说「这是某某所说」时，势必有「既然由他所说，此即为量，现在无需考察，此即如是」，此即可信。\textbf{不欺骗世间}，意即以此真实语性，不欺骗世间。
\item 在「舍弃了离间语」等中，对其说此言语，在其心中对自己起喜爱相，对他人起空无相，此即\textbf{离间语}。而以之令自己和他人粗恶，此语本身粗恶，既不悦耳，又不入心，此即\textbf{粗恶语}。以之作无聊、无义的闲谈，此即\textbf{绮语}。作为彼等根本的思，仍得名离间语等，且其即此处之所指。
\item 这里，心杂染者为了分裂他人或取悦自己发起的身语行之思，即离间语。若其造成分裂者的功德少为小罪，功德大为大罪。其要素有四：将被分裂之他人，以「如此，他们将成数个、将会分离」之分裂前导或以「如此，我将成可喜、亲厚」之取悦，为此的努力，他人了知此义。
\item \textbf{为了分裂此等}，即为了分裂被称为\textbf{此处}、在其跟前曾作听闻者。\textbf{调停分裂}，即对二友人，或对同一亲教师者，或对以无论何种原因而分裂者，前往彼彼处，说「这对你们——出生在这样的家族、又如是多闻者——不合适」等等，作、随作和解。\textbf{增进}，即增进和解，意即见到二人和合，说「这对你们——出生在这样的家族、具足这样的功德者——才合适」等等，而作巩固。
\item 和合则乐，即\textbf{乐于和合}，意即在不和合处，甚至不愿居住。文本或作 samaggarāmo，仍同此义。\textbf{喜于和合}，即在和合众中则喜，意即舍弃彼等，甚至不愿去往别处。或见或闻和合则欢喜，即\textbf{欢喜和合}。\textbf{说令和合之语}，即唯说令有情和合、阐明和合功德之语，而非其它。
\item 发起戳中他人要害之身语行的极端粗恶之思，为\textbf{粗恶语}。为其彰明而说此事。据说，有一男孩不听母亲的话，跑进林野，母亲不能制止他，便骂道：「让凶恶的水牛追着你！」于是，水牛便立刻出现在林野中。男孩便立下誓言：「愿我母亲口中所说的不成，愿她心中所想的得成！」水牛便像被缚在那里一样站着。如是，虽然是戳中要害之行，但因心地柔软，不成为粗恶语。因为父母有时会对孩子说：「让强盗把你们撕成碎片！」却连一片青莲叶落在他们身上也不愿意。而阿阇黎、亲教师有时对依止者说：「这些无惭无愧的逛什么？速速赶走他们！」但也希望他们有阿含和证悟的成就。且正如因心地柔软便不成为粗恶语，如是因言语柔软也不成为非粗恶语。因为欲谋害者的话语——让他好睡——并非粗恶语，但因心地粗恶，它便是粗恶语。就其所转起者而言，功德小者为小罪，功德大者为大罪。其要素有三：将被骂詈的他人、忿怒之心、骂詈。
\item \textbf{无瑕}，过失被称为瑕，其无有瑕者即无瑕，意即无过失，如在「部件无瑕，伞盖洁白」{\bluefootnotesize 〔自说第 65 经〕}中所说的无瑕一般。\textbf{悦耳}，即因字句甜美，双耳悦乐，非如针刺令生耳痛。因语意甜美，全身不生忿恨而生亲爱，即\textbf{可亲}。进入心中不受妨碍，悄然入心，即\textbf{入心}。因功德圆满，如城中之状，为\textbf{礼貌}，好似城中长大的女子般细致，亦为礼貌，城内所有者亦为礼貌，意即城民们的谈论。因为城民们的谈论得体，称父辈为父，兄辈为兄，母辈为母，这样的谈论为众人喜爱，即\textbf{众人喜爱}。唯因喜爱，令众人悦意、心意增长，即\textbf{众人悦意}。
\item 发起令知晓无义之身语行的不善思，即\textbf{绮语}。它对习行钝者为小罪，对习行大者为大罪。其要素有二：以无义之论为前导的婆罗多战争、劫持悉多等，且说出这样的谈论。
\item 适时地说，即\textbf{时语者}，意即了解了当说的合适时机而说。唯说实、真实、如实、实相，即\textbf{实语者}。唯依现法、来世之义利而说，即\textbf{义语者}。唯依九出世间法而说，即\textbf{法语者}。唯依防护律仪、舍断律仪而说，即\textbf{律语者}。置放的空间被称为贮藏，对其应贮藏，即\textbf{具有价值}，意即\textbf{说}在心中适合贮藏\textbf{之语}。\textbf{适时地}，意即即便在以「我将说具有价值的话」说这样的话时，不在非时说，唯寻找合适的时机而说。\textbf{有着理由}，即有着比较，意即有着原因。\textbf{具有边界}，意即显示界限后，令其界限可知而说。\textbf{伴有义利}，即以多种方法亦难以穷尽的义利之具足。或者，当义说者说义利时，由伴有此义故，说伴有义利之语，即是说不舍一说一。\end{enumerate}

{\Large 10.} \textbf{『沙门乔达摩戒离破坏种子、生物』，诸比丘……}\\
\textbf{『沙门乔达摩是一食者，节制于夜，戒离非时食』……}\\
\textbf{『沙门乔达摩戒离歌舞伎乐、观看演出』……}\\
\textbf{『沙门乔达摩戒离基于严饰的鬘香涂油之持戴、装饰』……}\\
\textbf{『沙门乔达摩戒离高床、大床』……}\\
\textbf{『沙门乔达摩戒离接受金、银』……}\\
\textbf{『沙门乔达摩戒离接受生谷』……}\\
\textbf{『沙门乔达摩戒离接受生肉』……}\\
\textbf{『沙门乔达摩戒离接受女人、女童』……}\\
\textbf{『沙门乔达摩戒离接受女奴、男奴』……}\\
\textbf{『沙门乔达摩戒离接受山羊、绵羊』……}\\
\textbf{『沙门乔达摩戒离接受鸡、豚』……}\\
\textbf{『沙门乔达摩戒离接受象、牛、马、牝马』……}\\
\textbf{『沙门乔达摩戒离接受田、地』……}\\
\textbf{『沙门乔达摩戒离从事信使、差役』……}\\
\textbf{『沙门乔达摩戒离买卖』……}\\
\textbf{『沙门乔达摩戒离秤的虚假、铜的虚假、称量虚假』……}\\
\textbf{『沙门乔达摩戒离索贿、欺诈、虚诳的邪路』……}\\
\textbf{『沙门乔达摩戒离砍、杀、缚、潜伏、掠夺、暴行』，诸比丘！说着如来赞誉的凡夫或如此说。}

\begin{enumerate}\item \textbf{破坏种子、生物}，意即\textbf{戒离}以砍、断、煮等方式破坏根种子、茎种子、节种子、插种子、种子种子等五种种子，以及任何青草绿树等生物。
\item \textbf{一食者}，有早餐之食和晚餐之食两种食，其中，早餐之食以正午前为界，另一则从正午以后到明相以前，所以，在正午前就算吃了十次，也是一食者，就此而说「一食者」。夜间之食即夜，节制此即\textbf{节制于夜}。过了正午直到日落之食名为非时食，由戒离此故，为\textbf{戒离非时食}。在何时戒离？从在无劣河畔出家之日起。
\item 由不随顺教法故，观看作为相违的演出，为观看演出。依自己舞蹈、令人舞蹈等的舞蹈、歌唱、伎乐，与观看依孔雀舞等转起的舞蹈等的演出，为\textbf{歌舞伎乐、观看演出}。因为舞蹈等，无论自己去做、让他人做，或是去看所做，对比丘和比丘尼都不合适。
\item 在「鬘」等中，\textbf{鬘}，即任何的花。\textbf{香}，即任何种类的香。\textbf{涂油}，即作用于皮肤。这里，当穿戴时，名为\textbf{持戴}，当补足缺陷时，名为\textbf{装饰}，当因香和皮肤而受用时，名为\textbf{严饰}。\textbf{基于}，即是称原因。所以，意即戒离于此大众以恶习之思所作的持戴花鬘等。
\item 超量者被称为\textbf{高床}。\textbf{大床}，即不被允许的铺盖。意即戒离于此。
\item \textbf{金}，即黄金。\textbf{银}，即钱币、铜币、硬币、木币等进入流通者。戒离此两者的接受，意为既不受持、不教人受持，也不受用被存放者。
\item \textbf{接受生谷}，即接受被称为稻、米、大麦、小麦、黍、豆、稗的七类生谷。且不仅接受这些，连触碰对比丘也不合适。
\item 在\textbf{接受生肉}中，除了指定许可外，唯接受生肉和鱼对比丘不合适，而非触碰。
\item 在\textbf{接受女人、女童}中，女人，即已与男人交往者，其余名为女童，对她们的接受和触碰都不被允许。
\item 在\textbf{接受女奴、男奴}中，只有以「女奴、男奴」接受彼等不合适，但当作如是说「让我给一名作净者、让我给一名园工」时，则合适。
\item 在「山羊、绵羊……田、地」中，允不允许之法当依律藏考察。这里，\textbf{田}者，即主粮生长处。\textbf{地}者，即蔬菜生长处。或者，两者生长处为田，为此而尚未开垦的地区为地。且此中，以田地为首，池塘、水沼等亦被包含在内。
\item 使者之事被称为\textbf{信使}，即得了在家人交待的书信或口信，四处奔走。\textbf{差役}，即从家到家被派遣的杂役。\textbf{从事}者，即行此二者。所以，当知此中之义为从事信使与差役。
\item \textbf{买卖}，即买与卖。
\item 在「秤的虚假」等中，\textbf{虚假}，即欺诈。这里，\textbf{秤的虚假}有外形虚假、肢体虚假、握持虚假、覆蔽虚假等四种。这里，\textbf{外形虚假}者，先做两杆外形相同的秤，在收取时以大秤收，给出时以小秤给。\textbf{肢体虚假}者，在收取时以手压秤的后部，给出时压前部。\textbf{握持虚假}者，在收取时握绳的根部，给出时握头部。\textbf{覆蔽虚假}者，把秤掏空后装入铁粉，在收取时令其居后部，给出时令其居头部。
\item 铜是说金盘，以之欺诈即\textbf{铜的虚假}。如何？先作一金盘，再将另外两三个铜盘做成金色，随后去到一国，进入任一富庶之家，说「你们买金器吗」，当被询问价格时，以较低的价格令其易出。随后，当他们说「如何知道这些是真金」时，说「你们先验了再收」，先把金盘在石头上摩擦，再给出所有盘子而去。
\item \textbf{称量虚假}者，依中缺、尖缺、绳缺而有三种。这里，\textbf{中缺}在称量酥、油时使用。即在收取彼等时，以底部穿孔的量器，说「慢慢地灌」，让内瓶装了许多而收，在给出时则堵上缺孔，速速装满而给。\textbf{尖缺}在称量胡麻、米粒时使用。即在收取彼等时，慢慢地堆起尖顶而收，在给出时急急装满，抹去尖顶而给。\textbf{绳缺}在丈量田地时使用。即在未得贿赂时，将小田作大而量。
\item 在「索贿」等中，\textbf{索贿}，即收取贿赂，将无主之物变为有主。\textbf{欺诈}，即以彼彼方法欺诈他人。此处有一事例：据说，有一猎人捕了鹿和鹿崽而来，一无赖便对他说：「先生！鹿值多少？鹿崽呢？」当被告知「鹿二钱，鹿崽一钱」时，便给了一钱，牵了鹿崽后没走几步即返，说：「先生！我说的不是鹿崽，给我鹿！」「那么，就给二钱！」他便说：「先生！我刚才不是给了你一钱？」「是，已经给了。」「你牵这鹿崽，那么那一钱和这价值一钱的鹿崽，就是二钱。」他考虑到「他说了理由」，便牵了鹿崽，给了鹿。\textbf{虚诳}，即以修行或幻术把非首饰变为首饰，非摩尼变为摩尼，非黄金变为黄金，以伪装欺诈。\textbf{邪路}，即歪路，为以上索贿等的总名。所以当视此中之义为索贿邪路、欺诈邪路、虚诳邪路等。有些人说，在展示某物后替换为另一物为邪路，但这已由欺诈所摄。
\item 在「砍」等中，\textbf{砍}即砍手等，\textbf{杀}即谋杀，\textbf{缚}即以绳索等束缚。\textbf{潜伏}有雪中潜伏、丛林潜伏两种。在落雪时以雪覆蔽，抢劫行路之人，此即雪中潜伏。以丛林等覆蔽而抢劫，此即丛林潜伏。\textbf{掠夺}是说对村镇等行劫掠。\textbf{暴行}，即粗暴之行。进入家舍后，用刀抵住人们的胸膛，取走所欲之物。如是，沙门乔达摩戒离这些砍、杀、缚、潜伏、掠夺、暴行，诸比丘！说着如来赞誉的凡夫或如此说。\end{enumerate}

\section{中戒}

{\Large 11.} \textbf{或者，『好比有些沙门婆罗门尊者受用了信施之食，他们从事破坏这样的种子、生物而住——例如根种子、茎种子、节种子、插种子、种子种子为第五——如此，沙门乔达摩戒离这样的破坏种子、生物』，诸比丘！说着如来赞誉的凡夫或如此说。}

\begin{enumerate}\item 现在，为详述中戒，说了「或者，好比有些尊者」等。此处为不明了之词的解释。\textbf{信施}，即信业、果、此世、后世等所施者，意即非如以「此系吾亲、吾友」，或以「他将回报、他先前曾行此」之所施。因为如是所施者不名为信施。\textbf{食}只是开示的纲目，从意义上则是说受用信施之食、穿衣、坐床座、服病药等一切。
\item \textbf{例如}是不变词。其义为：什么是这他们所从事破坏的种子、生物？随后，为显示此，说了「根种子」等。这里，\textbf{根种子}者，即姜黄、生姜、鸢尾根、白菖蒲、麦冬、黑根铁筷子、香根草、香附子等等。\textbf{茎种子}者，即菩提树、榕树、黄葛树、无花果树、红椿、林檎等等。\textbf{节种子}者，即甘蔗、芦苇、竹子等等。\textbf{插种子}者，即罗勒、墨角兰、须芒草等等。\textbf{种子种子}者，即主粮、蔬菜等等。凡此一切，从树上脱离仍能成长的，称为种子，而未从树上脱离、未干枯的，称为生物。这里，当知破坏生物为波逸提事，破坏种子为恶作事。\end{enumerate}

{\Large 12.} \textbf{或者，『好比有些沙门婆罗门尊者受用了信施之食，他们从事这样的受用贮藏物而住——例如贮藏食物、贮藏饮品、贮藏衣服、贮藏车乘、贮藏卧具、贮藏香料、贮藏财货——如此，沙门乔达摩戒离这样的受用贮藏物』，诸比丘！说着如来赞誉的凡夫或如此说。}

\begin{enumerate}\item \textbf{受用贮藏物}，即贮藏物之受用。这里，有两种说法，依律和依损减。先依律说，今天接受的任何食物，在翌日便成贮藏物，其受用为波逸提。但将自己所得的给了沙弥，让他们存放所得，在第二天受用则合适，但不成损减。
\item 于\textbf{贮藏饮品}亦如此法。这里，饮品者，即芒果饮品等八种饮品，以及与之相符者。彼等的抉择在「一切善见」中已述。
\item 于\textbf{贮藏衣服}，未决意、未分配即为贮藏，并触犯损减，这是概括之谈。若从字面\footnote{概括 \textit{pariyāya}、字面 \textit{nippariyāya}，也可作「间接、直接」。}，则应满足于三衣，得到第四件应施与他人。如果无人可施，而欲施与之人为了诵戒或遍问而离开，则刚一回来便应施与，不施即不合适。当衣不足时，对存有期望者，可存放至许可的时间\footnote{此句当指舍波逸提之三，许可的时间即一个月。}。因不得针线、裁缝，则过此更作律的羯磨后适合存放。但以「当此旧时，我又从何而得如此者」便不适合存放，成为贮藏，并触犯损减。
\item 于\textbf{贮藏车乘}，车乘者，即乘、马车、车、战车、轿、椅轿等，此非出家人的车乘，唯有鞋子是出家人的车乘。而对一比丘，一为林野，一为洗足，最多适合两双鞋。得到第三双应施与他人。即以「当此旧时，我又从何而得其它」便不适合存放，成为贮藏，并触犯损减。
\item 于\textbf{贮藏卧具}，卧具即床。对一比丘，一在内室，一在昼处，最多适合二床。过此而得，应施与其他比丘或众人，不施即不合适，成为贮藏，并触犯损减。
\item 于\textbf{贮藏香料}，当比丘患有痒疥等皮肤病时，适合用香料。让人带来这些香料后，当病愈时，或施与其他患者，或供于门前五指的熏屋处等。以「再生病时还会有吗」便不适合存放，成为贮藏，并触犯损减。
\item \textbf{财货}，当视为上述之余。例如，兹有比丘，以「在那样的场合会有帮助」，让人带来胡麻、米、豆、菜豆、椰子、盐、鱼、肉、干肉、酥、油、砂糖和器皿后存放。他在雨季早早就让沙弥们煮粥，受用后吩咐道：「沙弥！村中泥泞难入，去！到了某家，告知我在寺内坐着！从某家取了凝乳等来！」当比丘们说「尊者！你为什么不入村」时，说：「朋友！现在村子难入。」他们说：「好吧！尊者！您等着！我们觅了乞食带来！」便去了。然后，沙弥取了凝乳等，备好食物和调味后供上，正享用时，护持们又送来食物，他便又从中享用美味。然后，比丘们带了团食而返，仍延颈从中享用美味。如是便度过四个月。这被称为——比丘以秃头地主之命而活，非以沙门之命。像这样的便名为贮藏财货。
\item 然而，在比丘的住处储存一升米、一块糖、四分量的酥等这些则合适，为了在非时到来的盗贼。因为他们连这些财货的招待也得不到的话，将会害命，所以若是没有这些，让人带来存放则合适。且在不安乐时，自己受用此中允许的也合适。在允许的寮房里，即便存放许多也不名为贮藏。但对如来，无论是一升米等中的任何，或四肘量的布片，都不会以「这在今天或明天将是我的」而存放。\end{enumerate}

{\Large 13.} \textbf{或者，『好比有些沙门婆罗门尊者受用了信施之食，他们从事观看这样的演出而住——例如歌舞、伎乐、戏剧、故事、手鼓、铙钹、陶筒、美术、旃陀罗、竹戏、清洗、象斗、马斗、水牛斗、牛斗、羊斗、公羊斗、鸡斗、鹌鹑斗、械斗、拳斗、摔跤、演练、队列、军阵、观兵——如此，沙门乔达摩戒离这样的观看演出』，诸比丘！说着如来赞誉的凡夫或如此说。}

\begin{enumerate}\item 在「观看演出」中，\textbf{舞}者，即任何的舞，即便行于道路，延颈观望也不合适。而此中详细的抉择，当如在「一切善见」中所述之法而知。如是，在一切与学处相关的经句处皆如此。此后更连这些也不说，我们仅随处解释所需者。
\item \textbf{戏剧}，即舞者之集会。\textbf{故事}，即婆罗多战争等。若在某处讲述，则不适合去往那里。\textbf{手鼓}，即铜锣，也有说是手锣。\textbf{铙钹}，即槌锣，也有说是以咒术令死尸立起。\textbf{陶筒}，即四方槽形的锣，也有说是陶罐之声。\textbf{美术}，即舞者的挥洒，或作美化，即是说壁画。\textbf{旃陀罗}，即铁球戏，也说是旃陀罗洗麻布的游戏。\textbf{竹戏}，即先架设竹子再游戏。\textbf{清洗}，即洗骨。据说，在有些地方，当人死时，亲人们不作荼毗，而是掘土安葬。然后，知晓这些尸体已经腐烂，便挖出来经清洗、涂香后放置。他们在节日时将骨头置于一处，在另一处置酒等，嚎啕悲泣，而后饮酒。如下所述：\begin{quoting}诸比丘！在南方诸国有所谓洗骨者，那里有食物、饮品、嚼食、主食、舔食、啜食、歌舞、伎乐。诸比丘！我说有此洗骨，非无有此。{\bluefootnotesize 〔增支部第 10.107 段〕}\end{quoting}有人则说，以幻术洗骨为清洗。
\item 在「\textbf{象斗}」等中，比丘既不适合与象等斗，也不适合令它们斗，或看它们斗。\textbf{摔跤}，即角力。\textbf{演练}，即对抗发生之处。\textbf{队列}，即点兵之处。\textbf{军阵}，即军队的驻扎，军队依车阵等驻扎。\textbf{观兵}，即以\begin{quoting}三头象是最少的象军。{\bluefootnotesize 〔波逸提第 324 段〕}\end{quoting}等方法所说的观摩军队。\end{enumerate}

{\Large 14.} \textbf{或者，『好比有些沙门婆罗门尊者受用了信施之食，他们从事这样的基于放逸的博戏而住——例如八格、十格、空盘、避路、接近、骰盘、棒戏、筹尺、骰子、叶笛、弯犁、空翻、风车、叶斗、车戏、弓戏、猜字、臆测、效残——如此，沙门乔达摩戒离这样的基于放逸的博戏』，诸比丘！说着如来赞誉的凡夫或如此说。}

\begin{enumerate}\item 处于放逸之中为基于放逸，此博戏与基于放逸，即\textbf{基于放逸的博戏}。每行各有八格，为\textbf{八格}。\textbf{十格}也同理。\textbf{空盘}同八格、十格一样，只是在空中游戏。\textbf{避路}，即在地上划好许多应随处避让的路线，避让着路线而游戏。\textbf{接近}，即接近的游戏。不得动摇叠在一起的棋子或石子，只用指甲移开或移近，如果那里有任何动摇，便算输，即如此游戏的同义语。\textbf{骰盘}，即在博盘上的骰子游戏。\textbf{棒戏}是说用长棍击打短棍的游戏。\textbf{筹尺}，即在筹尺上沾染染料或深红色或浆水，以「它会成为什么」而甩在地上或墙上，显现象、马等形状的游戏。\textbf{骰子}，即骰块的游戏。
\item \textbf{叶笛}是说叶管，吹奏它作游戏。\textbf{弯犁}，即村童游戏的小犁。\textbf{空翻}是说翻筋斗游戏，即是说或在空中拿了棍子，或以头置地，以上下颠倒翻筋斗游戏。\textbf{风车}是说用多罗树叶等制成，靠风吹旋转的轮子。\textbf{叶斗}是说叶管，以之量取沙子等来游戏。\textbf{车戏}，即小车。\textbf{弓戏}，即小弓。\textbf{猜字}是说在空中或背上识别字母的游戏。\textbf{臆测}者，即以意猜测所思的游戏。\textbf{效残}者，即举凡独眼、跛足、驼背等的残疾，一一效仿表演的游戏。\end{enumerate}

{\Large 15.} \textbf{或者，『好比有些沙门婆罗门尊者受用了信施之食，他们从事这样的高床、大床而住——例如高椅、跏趺座、长毛毯、彩毯、白毯、花毯、棉垫、绘毯、上毛毯、单毛毯、迦絺萨、丝品、鸠多迦、象垫、马垫、车垫、羚羊毡、羚鹿高级铺盖、带顶罩者、双红枕——如此，沙门乔达摩戒离这样的高床、大床』，诸比丘！说着如来赞誉的凡夫或如此说。}

\begin{enumerate}\item \textbf{高椅}，即过量的坐处。考虑到「从事而住」，在所有句中都作业格。\textbf{跏趺座}，即在四脚雕刻猛兽图案后制成。\textbf{长毛毯}，即长毛的大毯，据说其毛四指有余。\textbf{彩毯}，即杂色编织的羊毛制毯。\textbf{白毯}，即羊毛制的白毯。\textbf{花毯}，即密集花样的羊毛制毯，也被称为余甘子叶。\textbf{棉垫}，以三种棉花中的任一种填充者，即棉垫。\textbf{绘毯}，即有狮虎等形象的彩色羊毛制毯。\textbf{上毛毯}，即双面可见的羊毛制毯，有人说是一面凸起的花样。\textbf{单毛毯}，即一面可见的羊毛制毯，有人说是双面凸起的花样。\textbf{迦絺萨}，即缝制了宝石、以丝和迦絺萨制成的铺盖。\textbf{丝品}，仍是缝制了宝石、丝线制成的铺盖。在律藏中说纯丝是合适的，但在长部义注中说，除了棉垫，所有缝制了宝石的长毛毯都不合适。
\item \textbf{鸠多迦}，即适于十六位舞女站立其上起舞的羊毛制毯。\textbf{象垫}、\textbf{马垫}，即垫在象、马背上的垫子。\textbf{车垫}同理。\textbf{羚羊毡}，即以羚羊皮按床的尺寸缝制的毡。\textbf{羚鹿高级铺盖}，有所谓的羚鹿皮，以其制作的高级铺盖，即最上铺盖之义。据说，它是在白布上铺上羚鹿皮缝制而成。\textbf{带顶罩者}，即与顶罩一起，意即与系在上面的红色华盖一起。即便是白色华盖，当下方有不允许的铺盖时，也不合适，但没有则合适。\textbf{双红枕}，即头枕与脚枕等在床两头的红色枕头，这不允许。但一个枕头的两侧是红色，或莲色、彩色，若尺寸合适，则合适。大枕头则被排除。非红色者，即便两个也合适。过此而得，应施与他人。当无法施与时，横铺于床，其上施以铺盖方可躺卧。但在「高椅」等处，可按所说之法而行，如下所述：\begin{quoting}诸比丘！我允许：截断高椅之足而受用，拆除跏趺座之兽而受用，解开棉垫而用作枕头，其余用作地毯。{\bluefootnotesize 〔小品第 297 段〕}\end{quoting}\end{enumerate}

{\Large 16.} \textbf{或者，『好比有些沙门婆罗门尊者受用了信施之食，他们从事这样的基于严饰的装饰而住——例如涂身、按捏、洗浴、按摩、镜子、眼膏、鬘香涂油、脸粉、脸油、手串、顶串、杖、筒、剑、伞、彩鞋、头巾、摩尼、羽扇——如此，沙门乔达摩戒离这样的基于严饰的装饰』，诸比丘！说着如来赞誉的凡夫或如此说。}

\begin{enumerate}\item 在「涂身」等中，从母胎出生的孩子的体味要到十二岁时才消失，为除去他们身体的臭味，便用香粉等涂身，像这样的\textbf{涂身}不合适。而让具福德的孩子趴在大腿上，用油涂抹后，为了身材端正而按捏手足、大腿、肚脐等，像这样的\textbf{按捏}也不合适。\textbf{洗浴}，仍是为这些孩子们用香等洗浴。\textbf{按摩}，即如对大力士一般，用棍棒敲打手足，增长力气。\textbf{镜子}，即不适合携带任何镜子。\textbf{眼膏}，仅指庄严用的眼膏。\textbf{鬘}，即绑扎或未绑扎的花鬘。\textbf{涂油}，即任何作用于皮肤者。\textbf{脸粉}、\textbf{脸油}，即为了去掉脸上的黑痘等施以泥粉，当血色动摇时，施以芥子粉，当毒素消尽时，施以芝麻粉，当血色安定时，施以姜黄粉，当肤色恢复时，以脸粉涂脸，这一切都不合适。
\item 在「\textbf{手串}」等中，在手上系了彩色的螺贝、龟甲等游荡，这或其它所有手上的璎珞都不合适。另有人系在头\textbf{顶}游荡，围以金带、珠链等，这一切都不合适。另有人拿了四肘长的\textbf{杖}，或其它有装饰的杖而游荡，同样，左胁悬挂绘有男女图案等外表华丽的药\textbf{筒}，另有人穿戴剑鞘镶有耳珰宝石、极其锋利的\textbf{剑}，用五色线缝制、饰以鱼齿等的\textbf{伞}，金银杂色、饰以孔雀翎等的\textbf{鞋}，有人露出一肘宽、四指高的发际线，在额前系上如云边光带般的\textbf{头巾}布，戴着宝冠\textbf{摩尼}，拿着拂尘\textbf{羽扇}，这一切都不合适。\end{enumerate}

{\Large 17.} \textbf{或者，『好比有些沙门婆罗门尊者受用了信施之食，他们从事这样的旁论而住——例如王论、贼论、大臣论、军队论、怖畏论、战争论、食物论、饮品论、衣服论、卧具论、鬘论、香论、亲族论、车乘论、村论、镇论、城论、国论、女子论、勇士论、街道论、市井论、先亡论、差异论、世间传说、海洋传说、如此诸有论——如此，沙门乔达摩戒离这样的旁论』，诸比丘！说着如来赞誉的凡夫或如此说。}

\begin{enumerate}\item 由不能出离故，作为天界、解脱之道的旁门之论，而为\textbf{旁论}。这里，就诸国王而以「摩诃三摩多、曼达多\footnote{摩诃三摩多、曼达多：见\textbf{经集}·正游行经的义注。}、法阿育有如是大威德」等方法转起的论为\textbf{王论}。\textbf{贼论}等处同理。其中，以「某某国王英俊可观」等方法、依于俗家之论，即旁论，以「他即便如是大威德，也行归于尽」等转起，则仍属业处的范畴。于诸\textbf{贼}中，以「根本天有如是大威德，云鬘有如是大威德」等，就彼等的业说「哎！英雄」的依于俗家之论，即旁论。在\textbf{战争}中，以「在婆罗多战争等中，某某被某某如是杀害、如是击中」，唯依爱欲之味的论，即旁论，而以「他们也行归于尽」等转起的，则于一切处仍属业处。
\item 另外，在「食物」等中，不适合依爱欲之味，以「我们曾吃过、尝过如是有色有香有味、具足口感的」作谈论，但作有意义的「我们先前曾施与具戒者具足色等的\textbf{食物、饮品、衣服、卧具、鬘、香}，我们曾供养支提」，则适合谈论。而在\textbf{亲族论}等中，以「我们的亲族有勇、堪能」或「我们先前曾以如是多彩的车乘游荡」，依味著而说，则不适合，作有意义的「我们的亲族，他们也行归于尽」或「我们先前曾布施这样的鞋子给僧团」，则适合谈论。\textbf{村论}，也不适合依好住、难住、好乞食、难乞食等，或以「某某村民有勇、堪能」等依味著而说，但作有意义的「有信、净喜」或「行归灭尽」，则适合说。\textbf{镇、城、国论}等处同理。
\item \textbf{女人论}，也不适合就容貌、身材等依于味著，以「有信、净喜，行归灭尽」等则适合。\textbf{勇士论}，也不适合依味著而说「喜友是战士、勇士」，唯适合「他有信，行归于灭」。\textbf{街道论}，也不适合依味著而说「某某街道好住、难住，（居民）有勇、堪能」，唯适合「有信、净喜，行归灭尽」。\textbf{市井论}，也被称为水边论、水津论，或挑水女工论，也不适合依味著而说「明艳照人，能歌善舞」，唯适合以「有信、净喜」等方法。\textbf{先亡论}，即过去亲族之论。这里的抉择与在世的亲族论相同。
\item \textbf{差异论}，即脱离前后所说，其余有差异的无义之论。\textbf{世间传说}，即如「这世间由谁所创？由某某所创」、「鸦是白色的，因为骨头是白色，鹤是红色的，因为血是红色」等顺世、诡辩的谈论。\textbf{海洋传说}者，即如「为什么叫海洋、海？由海神所挖，所以叫海\footnote{海神 \textit{Sāgaradeva}：为摩诃三摩多 \textit{Mahāsammata} 的后裔，其父为「海 \textit{Sāgara}」，其子为「婆罗多 \textit{Bharata}」。}，以『我所挖』的手印宣告自身，所以叫海洋\footnote{这是语源学的解释，即以「自身 \textit{sayaṃ}」和「印 \textit{muddā}」为「海洋 \textit{samudda}」。}」等无义的海洋传说之论。「有」即增长，「无有」即减损，凡是说「如此有、如此无有」等无义的原因所转起之论，即\textbf{如此诸有论}。\end{enumerate}

{\Large 18.} \textbf{或者，『好比有些沙门婆罗门尊者受用了信施之食，他们从事这样的诤论而住——例如：你不知道这法律，我知道这法律，你怎么会知道这法律？你是邪行道，我是正行道，我的一致，你的不一致，应该先说的你后说，应该后说的你先说，你所深思的已被推翻，你的观点已被指摘，你已被折服，去发表观点吧！请澄清，要是你能够的话——如此，沙门乔达摩戒离这样的诤论』，诸比丘！说着如来赞誉的凡夫或如此说。}

\begin{enumerate}\item \textbf{诤论}，即争吵之论、愤激之论。这里，\textbf{我的一致}，即我的言语一致、平顺、符合义理、符合原因之义。\textbf{你的不一致}，即你的言语不一致、不平顺。\textbf{你所深思的已被推翻}，即你以长时习行而极练达者，经我一语便被推翻，更改而住，你一无所知之义。\textbf{你的观点已被指摘}，即你的过失已被我指摘。\textbf{去发表观点吧}，即为了摆脱过失\footnote{义注将此句与上句中的「观点 \textit{vāda}」都解释为「过失 \textit{dosa}」。}而行、而周游，去到各处修学之义。\textbf{请澄清，要是你能够的话}，即要是你自己能够，现在就请澄清！\end{enumerate}

{\Large 19.} \textbf{或者，『好比有些沙门婆罗门尊者受用了信施之食，他们从事这样的信使、差役而住——例如对国王、王大臣、刹帝利、婆罗门、长者、童子们说：去此处！来到彼处！带走这！从那里带来这——如此，沙门乔达摩戒离这样的信使、差役』，诸比丘！说着如来赞誉的凡夫或如此说。}

\begin{enumerate}\item 在信使之论中，\textbf{去此处}，即从这里去名为某某之处。\textbf{来到彼处}，即从此来到名为某某之处。\textbf{带走这}，即从此带走这个。\textbf{从那里带来这}，即从某处带这个到此。而略言之，则此信使者，除了五同法者、事关资助三宝者，其它俗家的口信概不适合。\end{enumerate}

{\Large 20.} \textbf{或者，『好比有些沙门婆罗门尊者受用了信施之食，他们成为诡诈者、虚谈者、现相者、压榨者\footnote{压榨者 \textit{nippesikā}，叶均作「嗔骂示相」。}、以利求利者，如此，沙门乔达摩戒离这样的诡诈、虚谈』，诸比丘！说着如来赞誉的凡夫或如此说。}

\begin{enumerate}\item 在「诡诈者」等中，以三种诡诈事诡诈、惊异世间者，即\textbf{诡诈者}。希求利养恭敬而作虚谈者，即\textbf{虚谈者}。以现相为彼等之习者，即\textbf{现相者}。以压榨为彼等之习者，即\textbf{压榨者}。以利养企求、追寻、遍求者，即\textbf{以利求利者}。即具足这些诡诈、虚谈、现相、压榨、以利求利之人的同义语。这在此是略说。而若详说，则此诡诈等在「清净道论·说戒品」中已援引圣典和义注，作了阐明。\end{enumerate}

\section{大戒}

{\Large 21.} \textbf{或者，『好比有些沙门婆罗门尊者受用了信施之食，他们藉由这样的旁道，以邪命营生——例如察肢、占相、预兆、占梦、特相\footnote{此「特相 \textit{lakkhaṇa}」即大人相之相。}、鼠啮、火供、勺供、稃供、糠供、粒供、酥供、油供、口供、血供、肢咒、地咒、政术、平安咒、鬼咒、土咒、蛇咒、毒咒、蝎咒、鼠咒、鸟咒、鸦咒、虑朽、箭护卫、兽轮——如此，沙门乔达摩戒离藉由这样旁道的邪命』，诸比丘！说着如来赞誉的凡夫或如此说。}

\begin{enumerate}\item 以下是大戒。\textbf{察肢}，即以「在手足等处具足如此这般的肢体者便长寿、闻名」等方法转起的察肢之学。\textbf{占相}，即占相之学。据说，盘楼王在拳头里放了三颗珍珠，问占相者：「我手里是什么？」他便四处观望，恰在此时，一只被壁虎捕获的苍蝇挣脱了，他便说「珍珠\footnote{此中「挣脱」和「珍珠」的原文均为 muttā。}」。又被问是「几个」，听到鸡鸣三次之声，他便说「三」。如是将彼彼关联，从事占相而住。\textbf{预兆}，即看到打雷等的大事发生，联想到「将会有此、将会如此」。\textbf{占梦}，即以「若人在上午做梦，则有如是果报，若人做此梦，则会发生此事」等方法从事占梦而住。\textbf{特相}，即具足此特相者为国王，具足彼则为王储，等等。\textbf{鼠啮}，即被老鼠所咬。因为藉此，对新布或旧布会联想到「从此开始，当被如是啮时，将会如此」。
\item \textbf{火供}，即「用这样的木头这样来供奉，将会如此」而供奉火。\textbf{勺供}等一如火供，唯以「用这样的勺子、用这样的糠来供奉，将会如此」等转起各别之说。这里，\textbf{糠}，即米糠。\textbf{粒}，即稻等及草类的粒实。\textbf{酥}，即牛酥等。\textbf{油}，即芝麻油等。用口含了芥子喷到火里，或者念完咒语再供奉，即\textbf{口供}。以右锁骨、膝盖的血来供奉，即\textbf{血供}。
\item \textbf{肢咒}，先前仅以察看肢体后记说，称为察肢，此处则在察看指骨后念诵咒语，以「他是否是族姓子、是否具足祥瑞」等记说，称为肢咒。\textbf{地咒}，即观察屋地、园地等的功过之咒。因为即便只见到粘土等的差别，经念咒语，也能看到地下三十肘、上空八十肘范围内的功过。\textbf{政术}，即 Abbheyya、Māsurakkha 之王论等的学问\footnote{此句费解，暂将 Abbheyya、Māsurakkha 作书名或人名处理。PTS 本作「Aṅgeyya、Māsurakkha 等的政治学 \textit{nītisattha}」。}。\textbf{平安咒}，即进入墓地后起宁静作用的咒语，也有说是豺呼之咒。\textbf{鬼咒}，即鬼医的咒术。\textbf{土咒}，即住在土屋者应学习的咒术。\textbf{蛇咒}，即治疗蛇啮的咒语和唤蛇的咒语。\textbf{毒咒}，即以之守护旧毒，或制造新毒，或仅令具毒。\textbf{蝎咒}，即治疗蝎子咬伤的咒语。\textbf{鼠咒}处亦同理。\textbf{鸟咒}，即与有翼、无翼、二足、四足者的鸣叫有关的鸟智。\textbf{鸦咒}，即乌鸦鸣叫之智，它是特别的学问，所以分开来说。
\item \textbf{虑朽}，即思及衰老。意即「现在，他能活如许、他能如许」等转起的预见之智。\textbf{箭护卫}，即箭的防卫，好让其不落于自己身上的咒语。\textbf{兽轮}，此总摄一切，依所有鸟类、四足者的鸣叫之智而说。\end{enumerate}

{\Large 22.} \textbf{或者，『好比有些沙门婆罗门尊者受用了信施之食，他们藉由这样的旁道，以邪命营生——例如摩尼相、衣相、棍相、剑相、刀相、箭相、弓相、武器相、女相、男相、童子相、童女相、男奴相、女奴相、象相、马相、水牛相、公牛相、牛相、羊相、公羊相、鸡相、鹌鹑相、蜥蜴相、耳环相、龟相、兽相——如此，沙门乔达摩戒离藉由这样旁道的邪命』，诸比丘！说着如来赞誉的凡夫或如此说。}

\begin{enumerate}\item 在\textbf{摩尼相}等中，意即以「这样的摩尼好，这样的不好，是或不是主人健康自在之因」等，依色泽、形态等从事摩尼等的相而住。这里，\textbf{武器}，即除了刀等的其余武器。\textbf{女相}等，当以这些男女等所住之家的兴衰而知。而在\textbf{羊相}等中，当知其差别为「这样的羊等的肉可吃，这样的不可吃」。
\item 另外，于此中的\textbf{蜥蜴相}，当知其差别为「在彩绘、珠宝等中，当有这样的蜥蜴存在，将会如此」。且对此有一事：据说，在一寺庙的彩绘中曾作有吹火的蜥蜴。自此以后，比丘们便生起大争论。一来访的比丘见到后，便将其抹去。自此以后，争论渐息。\textbf{耳环相}，当知可作珠宝之耳环，也可作房舍之屋顶\footnote{原文 kaṇṇikā 兼有耳环和屋顶之义。}。\textbf{龟相}和蜥蜴相相同。\textbf{兽相}总摄一切，以所有四足者的特相而说。\end{enumerate}

{\Large 23.} \textbf{或者，『好比有些沙门婆罗门尊者受用了信施之食，他们藉由这样的旁道，以邪命营生——例如国王们将要出发、国王们将不出发，内王们将到来、外王们将离开，外王们将到来、内王们将离开，内王们将胜利、外王们将失败，外王们将胜利、内王们将失败，如此将此胜而彼败——如此，沙门乔达摩戒离藉由这样旁道的邪命』，诸比丘！说着如来赞誉的凡夫或如此说。}

\begin{enumerate}\item \textbf{国王们将要出发}，即在某某日，随某某节日，名为某某的国王将要出发，如是记说国王们的外出之行。一切处均同此理。只是此中的「\textbf{不出发}」为离开后又回来。\textbf{内王们将到来、外王们将离开}，即城内我们的国王将迎战敌对的外王，随后他将回来，如是记说国王们的到来与离开。第二句亦同此理。\textbf{胜利}、\textbf{失败}则自明。\end{enumerate}

{\Large 24.} \textbf{或者，『好比有些沙门婆罗门尊者受用了信施之食，他们藉由这样的旁道，以邪命营生——例如将有月食、将有日食、将有星食，日月将顺行、日月将逆行，群星将顺行、群星将逆行，将有流星，将有天火，将有地震，将有天鼓，将有日月星辰的升降染净，月食将有如是果、日食将有如是果、星食将有如是果，日月顺行将有如是果、日月逆行将有如是果，群星顺行将有如是果、群星逆行将有如是果，流星将有如是果，天火将有如是果，地震将有如是果，天鼓将有如是果，日月星辰的升降染净将有如是果——如此，沙门乔达摩戒离藉由这样旁道的邪命』，诸比丘！说着如来赞誉的凡夫或如此说。}

\begin{enumerate}\item \textbf{月食}等，当知即以「在某某日，罗睺将遮掩月亮」来记说。另外，星辰与火星等的遮掩交会，即\textbf{星食}。\textbf{流星}，即流星从天掉落。\textbf{天火}，即天空昏暗，充满火光烟雾。\textbf{天鼓}，即干云轰鸣。\textbf{升}即生起，\textbf{降}即隐没，\textbf{染}即不净，\textbf{净}即清净。\textbf{将有如是果}，即为世间带来种种苦乐。\end{enumerate}

{\Large 25.} \textbf{或者，『好比有些沙门婆罗门尊者受用了信施之食，他们藉由这样的旁道，以邪命营生——例如将有好雨、将有恶雨，将有收获、将有饥馑，将有安稳、将有怖畏，将有疾疫、将有健康，计算、算数、合计、诗歌、顺世论——如此，沙门乔达摩戒离藉由这样旁道的邪命』，诸比丘！说着如来赞誉的凡夫或如此说。}

\begin{enumerate}\item \textbf{好雨}，即天恰当地带来滋益。\textbf{恶雨}，即干旱，即是说雨受阻。\textbf{计算}，即手算。\textbf{算数}是说算数无缺。\textbf{合计}，即以加法、成百等等算总数。若对此练达，则看到树便知有多少叶子。\textbf{诗歌}，即在\begin{quoting}诸比丘！有这四种诗人。哪四种？思诗人、闻诗人、义诗人、辩才诗人。{\bluefootnotesize 〔增支部第 4.231 段〕}\end{quoting}这四种诗人中，或依自己的所思，或依听闻「曾有国王名毗输安多罗」等的所闻，或依「其义为此，我将如是运用」等的意义，或依「无论见到什么，我将作与之相似的」等即兴的辩才，为生计而作诗。\textbf{顺世论}则已述\footnote{对顺世论的解释，在第 17 段的义注中曾有提及。}。\end{enumerate}

{\Large 26.} \textbf{或者，『好比有些沙门婆罗门尊者受用了信施之食，他们藉由这样的旁道，以邪命营生——例如娶嫁，合离，敛贷，祝好运、祝厄运，保胎，缚舌、合颚，手咒、耳咒，问镜、问童女、问天，侍日、侍大，喷火，召祥——如此，沙门乔达摩戒离藉由这样旁道的邪命』，诸比丘！说着如来赞誉的凡夫或如此说。}

\begin{enumerate}\item \textbf{娶}者，即「从某某家族，在某某节日，为此男孩带来女孩」，为作婚娶。\textbf{嫁}，即「请在某某节日将此女孩给予某某男孩，如此她将富有」，为作婚嫁。\textbf{合}，所谓合，即「今日之节甚好，就在今日和合，此后你们永无分离」，如是为作和合。\textbf{离}者，即「如果你们想要分离，就在今日分离，此后你们再无结合」，如是为作分离。\textbf{敛}，即「请在今日收取收入或债务，因为今日所收的能持久」，如是聚拢财富。\textbf{贷}，即「如果你们想以借款、放债等赚取财富，今日所赚的会是两倍、四倍」，如是赚取财富。\textbf{祝好运}，即令可喜、可爱，或令祥瑞。\textbf{祝厄运}，即与之相反。
\item \textbf{保胎}，即对治不正、已消、不住、已死之胎，意即为令不再亡失而给予药物。因为胎儿会因风、虫、业等三种原因亡失。这里，当因风亡失时，可予以令其平复、寒凉之药，当因虫亡失时，可治疗虫，而当因业亡失时，则连诸佛也无法阻止。
\item \textbf{缚舌}，即以咒术束缚舌头。\textbf{合颚}，即以缚口的咒术施以束缚，让颚不能动摇。\textbf{手咒}，即为了翻掌的念咒。据说，当站在七步之内念诵此咒时，另一人便会翻掌投掷。\textbf{耳咒}，即为了让耳朵听不到声音而念诵咒语。据说，念了之后，在裁判处想什么便说什么，对方既无法听到，便不能作出回答。
\item \textbf{问镜}，即让天人落至镜中而提问。\textbf{问童女}，即让天人落至童女身上而提问。\textbf{问天}，即让天人落至女奴身上而提问。\textbf{侍日}，即为了营生而敬事太阳。\textbf{侍大}，即同样敬事大梵。\textbf{喷火}，即以咒术从口中喷出火焰。\textbf{召祥}，即以「来！祥瑞！请到我头上来」等，以头召唤祥瑞。\end{enumerate}

{\Large 27.} \textbf{或者，『好比有些沙门婆罗门尊者受用了信施之食，他们藉由这样的旁道，以邪命营生——例如许愿、还愿，鬼业、土业，精术、阉术，地业、备地，漱口、沐浴、供奉，催吐、催泻，上泻、下泻、头泻，耳油、眼敷、鼻滴，眼膏、续眼膏，眼科、外科、儿科，上药、去药——如此，沙门乔达摩戒离藉由这样旁道的邪命』，诸比丘！说着如来赞誉的凡夫或如此说。}

\textbf{这，诸比丘！就是说着如来赞誉的凡夫所能说的小量、浅量的戒量。}

\begin{enumerate}\item \textbf{许愿}，即去到天人处说「要是我成就此事，就以彼彼供奉你」，在成就时应作的许愿允诺之事。而当已成就时，其事即\textbf{还愿}。\textbf{土业}，即住在土屋后，所获咒术之运用。在「\textbf{精术、阉术}」中，精指男子，阉指阉人。如此，让阉人成为男子即精术，让男子成为阉人即阉术，但行此事只能令其成为无欲者，而不能令性别消失。\textbf{地业}，即在未开发的地上建设房舍。\textbf{备地}，即说着「把这个那个拿来」，供奉土地。\textbf{漱口}，即用水令口清洁。\textbf{沐浴}，即为他人沐浴。\textbf{供奉}，即为了彼等而供奉火。
\item \textbf{催吐}，即给予药物令吐。\textbf{催泻}亦同理。\textbf{上泻}是往上泻出毒素。\textbf{下泻}是往下泻出毒素。\textbf{头泻}，即在头部泻出。\textbf{耳油}，即为了包裹双耳或消除创伤而调制药油。\textbf{眼敷}，即眼部热敷之油。\textbf{鼻滴}，即调油滴鼻。\textbf{眼膏}，即能剥离两或三层膜的碱水眼膏。\textbf{续眼膏}，即令其平复、寒凉的眼药膏。\textbf{眼科}，即眼科医疗。\textbf{外科}，即外科医疗。\textbf{儿科}，即照顾儿童的医疗。\textbf{上药}，即以此显示治疗身体。\textbf{去药}，即施了碱水等后，当创伤到了合适之时，除去彼等。\end{enumerate}

\section{前际想者}

{\Large 28.} \textbf{诸比丘！有其它的诸法，甚深、难见、难觉、寂静、胜妙、非寻思所行、微妙、智者所经验，如来经自身证知、作证而宣说彼等，藉此，说着如来如实赞誉者才能正确地说。诸比丘！什么是这些甚深、难见、难觉、寂静、胜妙、非寻思所行、微妙、智者所经验的诸法，如来经自身证知、作证而宣说彼等，藉此，说着如来如实赞誉者才能正确地说？}

\begin{enumerate}\item 如是，依梵赐所说赞誉之适用，详说了三种戒后，现在，依比丘僧团所说赞誉之适用，以「诸比丘！有其它的诸法，甚深、难见……」等方法，开始作「空性阐释」。
\item 这里的\textbf{法}——法字适用于功德、开示、圣典、分解有情等处，因为在\begin{quoting}法与非法，两者并非是相等的果报，\\非法导向地狱，法则得达善趣。{\bluefootnotesize 〔长老偈第 304 颂〕}\end{quoting}等处法字为功德，在\begin{quoting}诸比丘！我将对你们开示初善……的法。{\bluefootnotesize 〔中部第 3.420 段〕}\end{quoting}等处为开示，在\begin{quoting}于此，比丘精通法——契经、应颂……{\bluefootnotesize 〔增支部第 5.73 段〕}\end{quoting}等处为圣典，在\begin{quoting}此时，有诸法，有诸蕴……{\bluefootnotesize 〔法集论第 121 段〕}\end{quoting}等处为分解有情，而此处适于功德，所以当视此中之义为「诸比丘！有其它的如来功德」。
\item \textbf{甚深}，即好似大海，除如来外，以他人蚊嘴之针般的智不可得住。由甚深故，即\textbf{难见}。由难见故，即\textbf{难觉}。由熄灭一切热恼故，即\textbf{寂静}，亦由对寂静所缘转起而寂静。以不令厌足之义为\textbf{胜妙}，好似美味的食物一般。由为最上智的境域故，非寻思所能行，即\textbf{非寻思所行}。\textbf{微妙}，即由自性之精致、微细故。由非愚人的境域故，唯智者所知，即\textbf{智者所经验}。
\item \textbf{如来经自身证知、作证而宣说彼等}，即如来不经他人引导，唯由自身，以超胜之智现量后，宣说、显明、谈论、阐明彼等诸法之义。\textbf{藉此}，即藉由这些功德之法。\textbf{如实}，即如实。\textbf{说着赞誉者才能正确地说}，即欲说如来赞誉者能正确地说，不遗漏而说之义。那是哪些法受世尊如是赞叹呢？一切知智。假设如是，为什么用复数表达？由多样的心的组合，以及多样的所缘故。因为在四种智相应大唯作心中可得，且无一法非其所缘。如说：\begin{quoting}以知晓过去一切，为一切知智，以对此无有障碍，为无碍智。{\bluefootnotesize 〔无碍解道第 1.120 段〕}\end{quoting}等等。如此，由多样的心的组合，以及再再生起的多样的所缘故，用复数表达。
\item 而\textbf{其它的}，这是此处令作差别之语，即以「其它的、非戒离杀生等等，为甚深，非自明」等，与一切词相连。因为声闻波罗蜜之智甚深，但辟支菩提之智较之更深，故非此处的差别，而一切知智又较此（辟支菩提之智）更深，故也非此处的差别，然而，没有其它较此（一切知智）更深者，所以可得「甚深」之差别。难见、难觉等一切当知也同样。
\item \textbf{诸比丘！什么是这些……}，此即欲谈论这些法的发问。「\textbf{诸比丘！一些沙门婆罗门……}」等是问题的解答。若问：那为什么如是开始呢？因为遇到四处，诸佛便轰鸣大作，诉诸于智，佛智之大得以显现，开示便成甚深、触及三相、关涉空性。何等为四？律之施设、地之差别、缘之行相、教义之差别。所以，「此轻、此重，此可救药、此不可救药，此有犯、此无犯，此趣于断、此趣于出罪、此趣于忏悔\footnote{这里的「趣于 \textit{gāminī}」为「需要……、以……赎罪」之义。}，此为世间罪、此为施设罪，于此事应施设此」，在如是堕落之事处，便有学处之施设，他人对此无势无力，非他人之境域，唯是如来之境域。如此，遇到律之施设，诸佛便轰鸣大作，诉诸于智……关涉空性。
\item 同样，「这些名四念处……名八支圣道，名五蕴，名十二处，名十八界，名四圣谛，名二十二根，名九因，名四食，名七触，名七受，名七想，名七思，名七心，其中如许名欲界法，如许名色界、无色界所系法，如许名世间法，如许名出世间法」，分别了二十四全部发趣、无尽理法的阿毗达摩藏而作谈论，他人对此无势无力，非他人之境域，唯是如来之境域。如此，遇到地之差别的界定，诸佛便轰鸣大作，诉诸于智……关涉空性。
\item 同样，此无明对诸行以九行相为缘：生起为缘，转起、相、积累、连接、障碍、集、因、缘为缘，诸行等对识等也同样。如说：\begin{quoting}如何于缘摄受之慧为法住智？无明为诸行的生起住、转起住、相住、积累住、连接住、障碍住、集住、因住、缘住，以此九行相，无明为缘，诸行为缘生，此二皆为缘生法，如此，则于缘摄受之慧为法住智。亦于过去时，亦于未来时，无明为诸行的生起住……生为老死的生起住……缘住，以此九行相，生为缘，老死为缘生，此二皆为缘生法，如此，则于缘摄受之慧为法住智。{\bluefootnotesize 〔无碍解道第 1.45 段〕}\end{quoting}如是分别了这依彼彼法如是如是之缘相而转起的三轮转、三时、三连结、四要略、二十行相的缘起而作谈论，他人对此无势无力，非他人之境域，唯是如来之境域。如此，遇到缘之行相，诸佛便轰鸣大作，诉诸于智……关涉空性。
\item 同样，「四人名为常论者，四部分常论者，四有边无边者，四无尽闪烁者，二偶然发生者，十六有想论者，八无想论者，八非有想非无想论者，七断论者，五人名为现法涅槃论者，他们依于此，执持此」，摧破了六十二见，解开结缚、脱诸丛薮而作谈论，他人对此无势无力，非他人之境域，唯是如来之境域。如此，遇到教义之差别，诸佛便轰鸣大作，诉诸于智，佛智之大得以显现，开示便成甚深、触及三相、关涉空性。
\item 而在此处，可得教义之差别，所以，为了显明一切知智之大，并为了阐明开示之「空性阐释」，法王乃诉诸教义之差别，以「诸比丘！一些沙门婆罗门……」等开始作问题的解答。\end{enumerate}

{\Large 29.} \textbf{诸比丘！有一些沙门婆罗门为前际想者、前际随见者，就前际以十八事表达多样胜解之句。那些沙门婆罗门尊者凭藉什么、就什么而为前际想者、前际随见者，就前际以十八事表达多样胜解之句？}

\begin{enumerate}\item 这里，\textbf{有}，即存在、可见、可得。\textbf{诸比丘}是呼格。\textbf{一些}，即部分。\textbf{沙门婆罗门}，以步入出家之相为沙门，以出生为婆罗门，或者被世间共许为沙门或婆罗门。设想、意向前际而执持者，为\textbf{前际想者}，或有前际想者为前际想者。
\item 这里，「际」字见于小肠、内部、界限、低劣、对立面、部分等处。因为在\begin{quoting}充以小肠，充以胃……{\bluefootnotesize 〔经集第 197 颂〕}\end{quoting}等处，际字为小肠，在\begin{quoting}他们裹以包装行于世间，内里不净，外在光鲜。{\bluefootnotesize 〔相应部第 1.122 段〕}\end{quoting}等处为内部，在\begin{quoting}腰带之端老损。{\bluefootnotesize 〔小品第 278 段〕}\end{quoting}\begin{quoting}它（到达）草边，或路边，或岩边，或水边。{\bluefootnotesize 〔中部第 1.304 段〕}\end{quoting}等处为界限，在\begin{quoting}诸比丘！这是活命的尽头，即乞食。{\bluefootnotesize 〔相应部第 3.80 段〕}\end{quoting}等处为低劣，在\begin{quoting}此即苦边。{\bluefootnotesize 〔相应部第 2.51 段〕}\end{quoting}等处为对立面，在\begin{quoting}有身，朋友！是一边。{\bluefootnotesize 〔增支部第 6.61 段〕}\end{quoting}等处为部分。此处则适合于部分。
\item 「想」字也在\begin{quoting}尊者！请世尊住一劫！{\bluefootnotesize 〔长部第 2.167 段〕}\end{quoting}\begin{quoting}最好去躺下。{\bluefootnotesize 〔增支部第 8.80 段〕}\end{quoting}\begin{quoting}应以得许可者缝合未许可者。{\bluefootnotesize 〔波逸提第 371 段〕}\end{quoting}等适合于寿劫、藉口、律之许可等众多意义\footnote{「想 \textit{kappa}」字的训释另可参见\textbf{经集}·吉祥经的义注。}。当知此处则适合于爱、见。如说：\begin{quoting}想，有二想，爱想与见想。{\bluefootnotesize 〔大义释第 28 段〕}\end{quoting}所以，当知此中之义为：依于爱、见，设想、遍计过去的蕴部分而住。对如是设想前际而住者，以再再生起而唯追随前际之见，为\textbf{前际随见者}。这些如是之见者，\textbf{就}、凭藉、依据此\textbf{前际}，为令他人也持有见，\textbf{以十八事表达多样胜解之句}。
\item 这里，\textbf{多样}，即多种。\textbf{胜解之句}，即名称之句。或者，由超过实际之义、未能如其自性地把握而转起故，诸见被称为诸胜解，诸胜解之句即胜解之句，意即阐述见的言词。\textbf{十八事}，即十八种原因。\end{enumerate}

\subsection{常论者}

{\Large 30.} \textbf{诸比丘！有一些沙门婆罗门为常论者，以四事施设自我和世间是常。那些沙门婆罗门尊者凭藉什么、就什么而为常论者，以四事施设自我和世间是常？}

\begin{enumerate}\item 现在，对于欲谈论以之作表达的十八事的发问，以「\textbf{那些……尊者}」等方法问了之后，为分别、显明诸事，说了「\textbf{诸比丘！有……}」等等。
\item 这里，以之而说者为论，即见的同义语。彼等有常论，即\textbf{常论者}，意即常见者。当知此后其它类似句义均同此法。\textbf{自我和世间是常}，即执持「色等中的某一为自我、为世间」已，\textbf{施设}其为常、不死、恒常、坚牢。如说：\begin{quoting}以「色是自我、世间和常」施设自我和世间，同样，以「受、想、诸行、识是自我、世间和常」施设自我和世间。\end{quoting}\end{enumerate}

{\Large 31.} \textbf{于此，诸比丘！部分沙门或婆罗门随着热勤、随着精勤、随着致力、随着不放逸、随着正作意，触证这样的心定，当心如是等持时，忆念多样的宿住。例如：一生、二生、三生、四生、五生、十生、二十生、三十生、四十生、五十生、百生、千生、百千生、数百生、数千生、数百千生——我在彼处曾作如是名、如是族姓、如是肤色\footnote{肤色 \textit{vaṇṇa} 亦可作容貌、种姓解，这里据\textbf{清净道论}·说神通品第十三第 71 段，译作肤色。}、如是食、如是经验苦乐、如是寿限，从此殁已，我便生于彼处，我在彼处亦曾作如是名、如是族姓、如是肤色、如是食、如是经验苦乐、如是寿限，从此殁已，投生于此——如此忆念有概观、有行相的多样的宿住。}

\textbf{他便如是说：『自我和世间是常、贫瘠、峰耸、柱立，且这些有情转世、轮回、亡殁、投生，但确有与常物相等者。其因为何？我随着热勤…………如此忆念有概观、有行相的多样的宿住。我由此而知：「如此，自我和世间是常、贫瘠、峰耸、柱立，且这些有情转世、轮回、亡殁、投生，但确有与常物相等者。」』}

\textbf{诸比丘！此即第一根由，凭藉此、就此，一些沙门婆罗门为常论者，施设自我和世间是常。}

\begin{enumerate}\item 在「随着热勤」等中，精进以烧焦烦恼被称为\textbf{热勤}，以精勤于此为\textbf{精勤}，以反复从事为\textbf{致力}，意即\textbf{随着}、凭藉、依据如是三类精进。\textbf{不放逸}是说念之不离。\textbf{正作意}，即方法作意、路径作意，意义上即是说「智」。因为对住立于此作意者，宿住随念智得以成就，这即此处作意之意趣。所以，此中的略义为：凭藉精进、念与智。\textbf{这样的}，即如是种类的。\textbf{心定}，即心定。\textbf{触证}，即寻得、获得。\textbf{当心如是等持时}，即以此定完全置于、善放置于心中。\textbf{忆念多样的宿住}等之义，在「清净道论」中已述\footnote{即\textbf{清净道论}·说神通品第十三第 13 段及以下。}。
\item \textbf{他便如是说}，即他既具足如是禅那的威力，持见者便如是说。\textbf{贫瘠}，即如不孕之畜、不实之多罗树一般无果，一无所生。以此拒斥执持自我和世间之禅那等的能生色等之相。如山峰一般耸立为\textbf{峰耸}。\textbf{柱立}，如城门柱一般竖立为柱立。好比深埋的门柱，不动而立，如是竖立之义。以此二者表明世间的不败坏相。而有人将文本读为 īsikaṭṭhāyiṭṭhito，并说如萱草内的草芯般站立。对此，其意趣为：即此所说的「出生」，便如草芯本来存在，从萱草中伸出一般。且因为草芯般站立，所以\textbf{这些有情转世}，意即从此去到别处。
\item \textbf{轮回}，即持续周流。\textbf{亡殁}，即如是得称。\textbf{投生}也同样。但在义注中，先说了「自我和世间是常」，现在，此持见者又以「这些有情转世」等语亲自打破自己的论点，持见者的知见并不固定，好似埋在稃堆里的树桩般晃动，好似疯人的篮子里饼块、粪便、牛粪等，即是说其中良莠不齐。在\textbf{但确有与常物相等者}中，常物，即因恒常地现存，认为如同大地，也如须弥山与日月，随后认为自我与这些相等，而说「但确有与常物相等者」。
\item 现在，为了证明「自我和世间是常」等之宗\footnote{此「宗 \textit{paṭiññā}」即因明中的宗。}而显明因，说了「\textbf{其因为何？我随着热勤……}」等等。这里，\textbf{我由此而知}，即以此所证的殊胜，我由现量而知，显示「我不仅仅出于信而说」。而此中的 \textbf{m} 字是为起连接词的作用而说。
\item \textbf{诸比丘！此即第一根由}，即在以「事」字说作「以四事」的四个根由中，此即第一根由，意即此忆念百千生之量为第一原因。
\item 以下二番也同理。只是此番以忆念数百千生而说，另二以忆念十、四十成坏劫。因为钝慧的外道忆念数百千生，中慧的十成坏劫，利慧的四十，无过此者。\end{enumerate}

{\Large 32.} \textbf{第二，沙门婆罗门尊者凭藉什么、就什么而为常论者，施设自我和世间是常？于此，诸比丘！部分沙门或婆罗门…………一成坏、二成坏、三成坏、四成坏、五成坏、十成坏…………如此忆念有概观、有行相的多样的宿住。他便如是说：…………}

\textbf{诸比丘！此即第二根由，凭藉此、就此，一些沙门婆罗门为常论者，施设自我和世间是常。}

{\Large 33.} \textbf{第三，沙门婆罗门尊者凭藉什么、就什么而为常论者，施设自我和世间是常？于此，诸比丘！部分沙门或婆罗门…………十成坏、二十成坏、三十成坏、四十成坏…………如此忆念有概观、有行相的多样的宿住。他便如是说：…………}

\textbf{诸比丘！此即第三根由，凭藉此、就此，一些沙门婆罗门为常论者，施设自我和世间是常。}

{\Large 34.} \textbf{第四，沙门婆罗门尊者凭藉什么、就什么而为常论者，施设自我和世间是常？于此，诸比丘！部分沙门或婆罗门是寻思者、考察者。他以寻思推导，以考察竞逐，由自身理解，便如是说：『自我和世间是常、贫瘠、峰耸、柱立，且这些有情转世、轮回、亡殁、投生，但确有与常物相等者。』}

\textbf{诸比丘！此即第四根由，凭藉此、就此，一些沙门婆罗门为常论者，施设自我和世间是常。}

\begin{enumerate}\item 在第四番中，以寻思为\textbf{寻思者}，或其有寻思为寻思者，即经寻思、寻而持见者的同义语。具足考察为\textbf{考察者}。考察，即衡量、选择、认可。因为好比男子以杖考察水后进入，如是，若经衡量、选择、认可而执持见，当知他即考察者。\textbf{以寻思推导}，意即以各种方法寻思。\textbf{以考察竞逐}，即以所说品类之考察竞逐。\textbf{由自身理解}，即源于自己的理解。\textbf{便如是说}，即执持常见而如是说。
\item 这里，有四种寻思者：随闻者、忆念前生者、有所得者、纯思者。这里，若听闻了「曾有国王名毗输安多罗」等，便寻思「那么，毗输安多罗就是世尊\footnote{毗输安多罗是世尊的前前世。}，自我是常」而执持见，此名随闻者。忆念了二三生，便寻思「我过去在某处名某某，所以自我是常」，此名忆念前生之寻思者。若以所得而寻思「如我现在之自我快乐，过去也曾如是，将来也会如是」而执持见，此名有所得之寻思者。仅以「当如是时便有此」之寻思而执持，名纯思者。\end{enumerate}

{\Large 35.} \textbf{藉此，诸比丘！这些沙门婆罗门为常论者，以四事施设自我和世间是常。诸比丘！因为任何沙门或婆罗门为常论者，施设自我和世间是常，他们全都以此四事，或以其中某个，无外于此。}

\begin{enumerate}\item \textbf{或以其中某个}，即以此四事之中的某一，或二或三。\textbf{无外于此}，即此四事之外，无有其它一因可作常之施设，以此作不可逆转的狮子吼。\end{enumerate}

{\Large 36.} \textbf{诸比丘！如来了知此：『这些见的根由被如是执持、如是执取，便成如是之趣、如是之将来。』如来了知此，还了知较此更上者，且不执取此了知，由不执取，便独自知晓寂灭。如实知晓了诸受的集、没、味、患、离而无取解脱者，诸比丘！即如来。}

\begin{enumerate}\item \textbf{诸比丘！如来了知此}，即诸比丘！如来从种种品类了知这四种见。随后，为显示了知此的行相，说了「这些见的根由」等等。
\item 这里，见即\textbf{见的根由}，另外，见的原因也是见的根由。如说：\begin{quoting}何等为八见的根由？诸蕴是见的根由，无明、触、想、寻、不如理作意、恶知识、外来之声是见的根由。\\诸蕴为因，诸蕴为缘，以依于见的根由而等起之义，如是诸蕴为见的根由。无明为因……恶知识为因。外来之声为因，外来之声为缘，以依于见的根由而等起之义，如是外来之声为见的根由。{\bluefootnotesize 〔无碍解道第 1.124 段〕}\end{quoting}
\item 首先，当见的根由被称为见，\textbf{如是执持}意为如是执持、采纳、转起「自我和世间是常」。\textbf{如是执取}，即以无疑之心反复捉摸、执取，总结以「唯此为真，余皆虚妄」。而当见的根由被称为原因，则正如被执持者等起诸见，如是被所缘、被转起、被习行所\textbf{执持}。以不见过患，被反复执持所\textbf{执取}。\textbf{如是之趣}，即如是地狱、畜生、饿鬼趣中的某一趣。\textbf{如是之将来}，即前词的异名，即是说如是种类的他世间。
\item \textbf{如来了知此}，即如来不仅了知持见者的原因和趣向，还了知所有，\textbf{还了知较此更上}的戒、定和一切知智。\textbf{且不执取此了知}，尽管了知如是种类的无上殊胜，不以「我了知」，不以爱、见、慢之执取而执取此。\textbf{由不执取，便独自知晓寂灭}，由如是不执取，以及不执取之缘，便亲身、自身知晓这些执取烦恼的寂灭。即显示：诸比丘！如来的涅槃显而易见。
\item 现在，由如是行道之如来所证得之寂灭，为显示其行道，宣说依于这些受的业处——外道贪染于这些受，以「在此处我们会快乐、在此中我们会快乐」而进入见的密林——说了「诸受的集」等等。
\item 这里，\textbf{如实知晓}，即\begin{quoting}「无明集而有受集」，以缘集之义见受蕴的生起，「爱集而有受集」，以缘集之义见受蕴的生起，「业集而有受集」，以缘集之义见受蕴的生起，「触集而有受集」，以缘集之义见受蕴的生起，乃至见到转生之相，也见受蕴的生起。{\bluefootnotesize 〔无碍解道第 1.50 段〕}\end{quoting}依于此五特相，如实知晓了\textbf{诸受的集}，\begin{quoting}「无明灭而有受灭」，以缘灭之义见受蕴的灭去，「爱灭而有受灭」，以缘灭之义见受蕴的灭去，「业灭而有受灭」，以缘灭之义见受蕴的灭去，「触灭而有受灭」，以缘灭之义见受蕴的灭去，乃至见到变易之相，也见受蕴的灭去。{\bluefootnotesize 〔无碍解道第 1.50 段〕}\end{quoting}依于此五特相，如实知晓了诸受的\textbf{没}，以\begin{quoting}缘受生起喜乐，这即受的味。{\bluefootnotesize 〔相应部第 3.26 段〕}\end{quoting}如实知晓了\textbf{味}，以\begin{quoting}受为无常、苦、变易法，这即受的患。\end{quoting}如实知晓了\textbf{患}，以\begin{quoting}于受调伏欲贪、舍弃欲贪，这即受的离。\end{quoting}如实知晓了\textbf{离}，以离去贪染而为无取著者。\textbf{无取解脱者，诸比丘！即如来}，当有此取时，若取著任何，则由已取著故，便有蕴\footnote{蕴 \textit{kandho}：PTS 本作「束缚 \textit{baddho}」。}，由无有此，不取著任何法，诸比丘！如来便解脱。\end{enumerate}

{\Large 37.} \textbf{这些，诸比丘！诸法甚深、难见、难觉、寂静、胜妙、非寻思所行、微妙、智者所经验，如来经自身证知、作证而宣说彼等，藉此，说着如来如实赞誉者才能正确地说。}

\begin{enumerate}\item \textbf{这些，诸比丘……}，即凡我所问的「诸比丘！什么是这些甚深……的诸法」，诸比丘！当知这些以「如来了知此，还了知较此更上者」所指示的一切知智之\textbf{诸法甚深、难见……智者所经验}。\textbf{藉此}，既非凡夫，也非须陀洹等中的任一能如实地说赞誉，而唯有如来是\textbf{说着如来如实赞誉者}，\textbf{才能正确地说}。
\item 如是，提问者所问的即一切知智，回答者所答的也正是此，而其中间则是诸见的分别。\end{enumerate}

\subsection{部分常论者}

{\Large 38.} \textbf{诸比丘！有一些沙门婆罗门为部分常、部分非常者，以四事施设自我和世间部分是常、部分非常。那些沙门婆罗门尊者凭藉什么、就什么而为部分常、部分非常者，以四事施设自我和世间部分是常、部分非常？}

\begin{enumerate}\item \textbf{部分常}者，即部分常论者。他们有两种：有情部分常者、行部分常者。此处也包摄了两种。\end{enumerate}

{\Large 39.} \textbf{诸比丘！会有一时，当某时、有时，在经历了长时之后，此世间毁坏。而世间毁坏时，多数有情转至光音天。他们在那里由意所成，以喜为食，自身光芒，行于空中，住于净处，长久而住。}

\begin{enumerate}\item \textbf{当} \textit{yaṃ}，只是不变词。\textbf{某时}，即某个时间。\textbf{有时}只是其异名。\textbf{长时}，即长时间。\textbf{经历}，即经过。\textbf{毁坏}，即消亡。\textbf{多数}，即指除去转生至更高的梵世间或无色界者而言\footnote{光音天在二禅天，其上更有「少净、无量净、遍净」等三禅天和「广果、无想有情、净居」等四禅天，以及无色界。}。
\item 由禅那之意转生故，为\textbf{由意所成}。喜是他们的所食、食物，为\textbf{以喜为食}。他们的光芒唯向自身，为\textbf{自身光芒}。在空中而行，为\textbf{行于空中}。住在洁净的庭园、宫殿、如意树下，为\textbf{住于净处}，或者，以洁净、悦意的衣装璎珞而住，为住于净处。\textbf{长久而住}，即最多八劫。\end{enumerate}

{\Large 40.} \textbf{诸比丘！会有一时，当某时、有时，在经历了长时之后，此世间回还。而世间回还时，空的梵宫显现。于是，某个有情由寿尽或福尽，从光音天众殁已，投生至空的梵宫。他在那里由意所成，以喜为食，自身光芒，行于空中，住于净处，长久而住。}

\begin{enumerate}\item \textbf{回还}，即安立。\textbf{空的梵宫}，因没有自然转生的有情而为空，意即梵众之地形成。它并无造作者或令造作者，而是按「清净道论」所说之法，由业缘时节等起\footnote{见\textbf{清净道论}·说道非道智见清净品第二十第 29 段。}而在宝地上形成。且仍在此自然形成之处，庭园、如意树等形成。于是，众有情出于天性，对曾住之处生起欣求，他们修习了初禅，随后降下，所以说了「\textbf{于是，某个有情……}」等等。
\item \textbf{由寿尽或福尽}，若造了显赫的福业而转生至任何少寿的天界，他们不能以自身的福力而住，只能按此天界的寿量而殁，被称为寿尽而殁。若造了有限的福业而转生至长寿的天界，也不能尽寿而住，中道便殁，被称为福尽而殁。\textbf{长久而住}，即一劫或半劫。\end{enumerate}

{\Large 41.} \textbf{他由长时独自居住故，生起不喜、焦渴：「哎！让其他有情也来此间吧！」于是，其他有情也由寿尽或福尽，从光音天众殁已，投生至梵宫，与此有情为伴。他们在那里也由意所成，以喜为食，自身光芒，行于空中，住于净处，长久而住。}

\begin{enumerate}\item \textbf{不喜}，即愿求后面的有情来到。在梵界中没有与嗔恚相应的不满。\textbf{焦渴}，即惊慌、战栗，它有恐惧之焦渴、爱之焦渴、见之焦渴、智之焦渴等四种。这里，\begin{quoting}缘生而有怖畏、畏惧、颤抖、身毛竖立、心的恐惧。缘老……病……死而有……恐惧。{\bluefootnotesize 〔分别论第 921 段〕}\end{quoting}名为\textbf{恐惧之焦渴}。\begin{quoting}哎！让其他有情也来此处吧！{\bluefootnotesize 〔长部第 3.38 段〕}\end{quoting}名为\textbf{爱之焦渴}。\begin{quoting}唯焦渴和战栗。\end{quoting}名为\textbf{见之焦渴}。\begin{quoting}他们听完如来法的开示，多数生起怖畏、悚惧、恐惧。{\bluefootnotesize 〔增支部第 4.33 段〕}\end{quoting}名为\textbf{智之焦渴}。而此处适合爱的焦渴和见的焦渴。\textbf{梵宫}，于此，以有最初转生者，不说为「空」。\textbf{投生}，即依投生而来。\textbf{为伴}，即一起。\end{enumerate}

{\Large 42.} \textbf{于此，诸比丘！这最初投生的有情作是念：『我是梵、大梵、胜者、不为所胜、普见者、自在天、主宰、作者、创造者、最胜者、指定者、掌控者、已生将生者之父。这些有情由我创造。其因为何？因为我先前作是念：「哎！让其他有情也来此间吧！」我的心愿如此，且这些有情到了此间。』}

\textbf{而那些后来投生的有情，他们也作是念：『此尊是梵、大梵、胜者、不为所胜、普见者、自在天、主宰、作者、创造者、最胜者、指定者、掌控者、已生将生者之父。我们由此尊梵创造。其因为何？因为我们看到他最初投生在此，我们则是后来投生。』}

\begin{enumerate}\item \textbf{胜者}，即征服已而住，说「我是长者」。\textbf{不为所胜}，即不为他人所胜。\textbf{普}是确然之语的不变词，以见为\textbf{见者}，意即「我见到一切」。\textbf{自在天}，即我对所有人行使权力。\textbf{主宰、作者、创造者}，即我是世上的主宰，我是世间的作者、创造者，大地、雪山、须弥山、轮围山、大海、日月都由我所造。\textbf{最胜者、指定者}，我是世间的最高和指定者，以「让你名为刹帝利，你为婆罗门、吠舍、首陀罗、在家人、出家人，乃至让你成为骆驼、成为公牛」，如是，认为「我是有情的分配者」。\textbf{掌控者、已生将生者之父}，即认为「我因曾行掌控而为掌控者，我是已生和将生者之父」。这里，卵生和胎生的有情，在蛋壳和胎膜内名为将生，从出离在外之时起名为已生。湿生者在第一心刹那为将生，从第二开始为已生。化生者在第一威仪路时为将生，从第二开始当知为已生。他以「他们全都是我的孩子」之想，认为「我是已生将生者之父」。
\item 现在，欲从原因上证明，在以「\textbf{这些有情由我创造}」立宗之后，说了「\textbf{其因为何}」等等。\textbf{此处}，即这样的状态，意即梵的状态。\textbf{我们由此……}，尽管依自身的业而死、而投生，唯凭妄想认为「我们由他创造」，如曲隙中的曲楔一般，跪倒并趋于其足下。\end{enumerate}

{\Large 43.} \textbf{于此，诸比丘！这最初投生的有情，他更长寿、更美貌、更大威德，而那些后来投生的有情，他们更短寿、更丑貌、更少威德。}

\begin{enumerate}\item \textbf{更美貌}，更美貌，意即英俊、明净。\textbf{更大威德}，即依主宰和随从而名望更大。\end{enumerate}

{\Large 44.} \textbf{有这可能，诸比丘！某一有情从彼众中殁后，来到此间。当来到此间时，从家出至非家。当从家出至非家时，随着热勤、随着精勤、随着致力、随着不放逸、随着正作意，触证这样的心定，当心如是等持时，忆念此宿住，未忆念此前。}

\textbf{他便如是说：『此尊是梵、大梵、胜者、不为所胜、普见者、自在天、主宰、作者、创造者、最胜者、指定者、掌控者、已生将生者之父。我们由此尊梵创造，他恒常、坚牢、常、非变易法、与常物相等，仍将如是而住。而我们曾由此尊梵创造，我们无常、不坚牢、短寿、殁法，来到此间。』}

\textbf{诸比丘！此即第一根由，凭藉此、就此，一些沙门婆罗门为部分常、部分非常者，施设自我和世间部分是常、部分非常。}

\begin{enumerate}\item \textbf{有这可能}，即有这原因。他从彼殁后，不去别处，唯来此处——就此而说。\textbf{从家}，即从俗家。\textbf{非家}，即出家。因为出家没有家中耕田、牧牛等的事务，所以被称为非家。\textbf{出至}，即参与。\textbf{未忆念此前}，即未忆念此宿住之前，当不能忆念，便停留在此而执持见。
\item \textbf{恒常}等中，当未见其投生，便说是\textbf{恒常}，未见死为\textbf{坚牢}，由永恒之相故为\textbf{常}，即便老去也无变易，为\textbf{非变易法}。第一番中，余皆自明。\end{enumerate}

{\Large 45.} \textbf{第二，沙门婆罗门尊者凭藉什么、就什么而为部分常、部分非常者，施设自我和世间部分是常、部分非常？诸比丘！有一些天人名为戏毁，他们过度地具足戏笑、游戏之乐法而住。当他们过度地具足戏笑、游戏之乐法而住时，念便忘失。由念的忘失，这些天人从彼众中亡殁。}

{\Large 46.} \textbf{有这可能，诸比丘！某一有情从彼众中殁后，来到此间。当来到此间时，从家出至非家。当从家出至非家时，随着热勤、随着精勤、随着致力、随着不放逸、随着正作意，触证这样的心定，当心如是等持时，忆念此宿住，未忆念此前。}

\textbf{他便如是说：『那些尊天非戏毁者，他们不过度地具足戏笑、游戏之乐法而住。当他们不过度地具足戏笑、游戏之乐法而住时，念便不忘失。由念的不忘失，那些天人不从彼众中亡殁，恒常、坚牢、常、非变易法、与常物相等，仍将如是而住。而我们曾是戏毁者，我们曾过度地具足戏笑、游戏之乐法而住。当我们过度地具足戏笑、游戏之乐法而住时，念便忘失。由念的忘失，我们便如是从彼众中亡殁，无常、不坚牢、短寿、殁法，来到此间。』}

\textbf{诸比丘！此即第二根由，凭藉此、就此，一些沙门婆罗门为部分常、部分非常者，施设自我和世间部分是常、部分非常。}

\begin{enumerate}\item 在第二番中，因游戏而污毁、亡失，为\textbf{戏毁}。文本也写作 padūsikā，在义注中无此。\textbf{过度}，即超过时间，意即极久。\textbf{具足戏笑、游戏之乐法}，即具足、从事戏笑之乐法和游戏之乐法，从事娱乐戏笑之乐和身语嬉戏之乐，意即据有所说品类的乐法而住。
\item \textbf{念便忘失}，即在嚼食、主食处，念便忘失。据说，他们因所证福德殊胜之巨大，当以自身的灿烂丰盈嬉戏星宿时，因成就之巨大，甚至不知道「我们是否吃过食物了」。于是，从错过一顿食物起，即便连续又吃又喝，仍然亡殁，不能存活。为什么？由业生火强、业生身弱故。因为人类的业生火弱，业生身强，由他们的火弱、业生身强故，即便过了七天，也能以温水清粥维持依处。而天人的火强，业生却弱，只错过一顿饭的时间，便不能安立。正好比夏天中午，在晒热的岩石上放置的红莲或青莲，在哺时即便以百罐浇洒，也不能恢复原状，仍然凋亡，如是，即便后来连续又吃又喝，仍然亡殁，不能存活。因此说「\textbf{由念的忘失，这些天人从彼众中亡殁}」。
\item 那他们是哪些天人呢？义注中并未考察「是这些天人」，但由以「天人的业生火强，业生却弱」作无差别的述说故，当知任何依靠段食的天人若如是行，他们就会亡殁。而有人说「这些天人是化乐天、他化自在天」，则他们只是以游戏之污毁被称为「戏毁」。当知此中其余仍如前法。\end{enumerate}

{\Large 47.} \textbf{第三，沙门婆罗门尊者凭藉什么、就什么而为部分常、部分非常者，施设自我和世间部分是常、部分非常？诸比丘！有一些天人名为意毁，他们过度地凝视彼此。当他们过度地凝视彼此时，便对彼此污染其心。他们对彼此心怀污染，身心疲惫。这些天人从彼众中亡殁。}

{\Large 48.} \textbf{有这可能，诸比丘！某一有情从彼众中殁后，来到此间。当来到此间时，从家出至非家。当从家出至非家时，随着热勤、随着精勤、随着致力、随着不放逸、随着正作意，触证这样的心定，当心如是等持时，忆念此宿住，未忆念此前。}

\textbf{他便如是说：『那些尊天非意毁者，他们不过度地彼此凝视。当他们不过度地彼此凝视时，便不对彼此污染其心。他们不对彼此心怀污染，身心不疲。那些天人不从彼众中亡殁，恒常、坚牢、常、非变易法、与常物相等，仍将如是而住。而我们曾是意毁者，我们曾过度地彼此凝视。当我们过度地彼此凝视时，便对彼此污染其心。我们对彼此心怀污染，身心疲惫。我们便如是从彼众中亡殁，无常、不坚牢、短寿、殁法，来到此间。』}

\textbf{诸比丘！此即第三根由，凭藉此、就此，一些沙门婆罗门为部分常、部分非常者，施设自我和世间部分是常、部分非常。}

\begin{enumerate}\item 在第三番中，以心意污毁、亡失，为\textbf{意毁}，他们是四大王天。据说，他们中一位天子以「我要嬉戏星宿」，带着随从登车上路，而另一位也在出发，看到他走在前面，想「这可怜人」，像先前没见过他一样，怒道：「他走起来像是开心得胀爆一样！」而走在前面的也回头看到那忿怒的——忿怒总是易被知晓——了知其忿怒之相后，反怒道：「你还发怒，你能对我作什么？这成就都是由我依施、戒等而得，并未依你。」因为当一人忿怒，另一不怒者尚能守护，但当双方忿怒，则一人之怒为另一之缘，彼之怒又为另一之缘，双双在悲泣的侍女中\footnote{侍女 \textit{orodha}：指上述二位四大王天天子的侍女。}亡殁。这是此处的法性。当知其余仍如前述。\end{enumerate}

{\Large 49.} \textbf{第四，沙门婆罗门尊者凭藉什么、就什么而为部分常、部分非常者，施设自我和世间部分是常、部分非常？于此，诸比丘！部分沙门或婆罗门是寻思者、考察者。他以寻思推导，以考察竞逐，由自身理解，便如是说：『凡被称为眼、耳、鼻、舌、身的，其自我为无常、不坚牢、不常、变易法。而凡被称为心、意、识的，其自我为恒常、坚牢、常、非变易法、与常物相等，仍将如是而住。』}

\textbf{诸比丘！此即第四根由，凭藉此、就此，一些沙门婆罗门为部分常、部分非常者，施设自我和世间部分是常、部分非常。}

\begin{enumerate}\item 在寻思论中，他看见眼等的断坏，而心，因为前前作了后后之缘便灭去，所以尽管比眼等的断坏更有力，他却看不到心的断坏。当他看不到此，便如是执持而说：「正好比飞鸟舍了一树后棲止于另一，如是当此自体断坏，心便去了别处。」当知此中其余仍如前述。\end{enumerate}

{\Large 50.} \textbf{藉此，诸比丘！这些沙门婆罗门为部分常、部分非常者，以四事施设自我和世间部分是常、部分非常。诸比丘！因为任何沙门或婆罗门为部分常、部分非常者，施设自我和世间部分是常、部分非常，他们全都以此四事，或以其中某个，无外于此。}

{\Large 51.} \textbf{诸比丘！如来了知此：『这些见的根由被如是执持、如是执取，便成如是之趣、如是之将来。』如来了知此，还了知较此更上者，且不执取此了知，由不执取，便独自知晓寂灭。如实知晓了诸受的集、没、味、患、离而无取解脱者，诸比丘！即如来。}

{\Large 52.} \textbf{这些，诸比丘！诸法甚深、难见、难觉、寂静、胜妙、非寻思所行、微妙、智者所经验，如来经自身证知、作证而宣说彼等，藉此，说着如来如实赞誉者才能正确地说。}

\subsection{有边无边论者}

{\Large 53.} \textbf{诸比丘！有一些沙门婆罗门为有边无边者，以四事施设世间之有边无边。那些沙门婆罗门尊者凭藉什么、就什么而为有边无边者，以四事施设世间之有边无边？}

\begin{enumerate}\item \textbf{有边无边者}，即有边无边论者，意即就有边、或无边、或有边无边、或非有边非无边而转起之论。\end{enumerate}

{\Large 54.} \textbf{于此，诸比丘！部分沙门或婆罗门随着热勤、随着精勤、随着致力、随着不放逸、随着正作意，触证这样的心定，当心如是等持时，对世间住于有边想。}

\textbf{他便如是说：『此世间具有边、具有边界。其因为何？因为我随着热勤…………对世间住于有边想。我由此而知：「如此，此世间具有边、具有边界。」』}

\textbf{诸比丘！此即第一根由，凭藉此、就此，一些沙门婆罗门为有边无边者，施设世间之有边无边。}

{\Large 55.} \textbf{第二，沙门婆罗门尊者凭藉什么、就什么而为有边无边者，施设世间之有边无边？于此，诸比丘！部分沙门或婆罗门随着热勤、随着精勤、随着致力、随着不放逸、随着正作意，触证这样的心定，当心如是等持时，对世间住于无边想。}

\textbf{他便如是说：『此世间无边、无边界。凡沙门婆罗门作如是说「此世间具有边、具有边界」，彼等虚妄，此世间无边、无边界。其因为何？因为我随着热勤…………对世间住于无边想。我由此而知：「如此，此世间无边、无边界。」』}

\textbf{诸比丘！此即第二根由，凭藉此、就此，一些沙门婆罗门为有边无边者，施设世间之有边无边。}

{\Large 56.} \textbf{第三，沙门婆罗门尊者凭藉什么、就什么而为有边无边者，施设世间之有边无边？于此，诸比丘！部分沙门或婆罗门随着热勤、随着精勤、随着致力、随着不放逸、随着正作意，触证这样的心定，当心如是等持时，对世间住于上下有边想，四旁无边想。}

\textbf{他便如是说：『此世间既有边又无边。凡沙门婆罗门作如是说「此世间具有边、具有边界」，彼等虚妄，凡沙门婆罗门作如是说「此世间无边、无边界」，彼等亦虚妄，此世间既有边又无边。其因为何？因为我随着热勤…………对世间住于上下有边想，四旁无边想。我由此而知：「如此，此世间既有边又无边。」』}

\textbf{诸比丘！此即第三根由，凭藉此、就此，一些沙门婆罗门为有边无边者，施设世间之有边无边。}

{\Large 57.} \textbf{第四，沙门婆罗门尊者凭藉什么、就什么而为有边无边者，施设世间之有边无边？于此，诸比丘！部分沙门或婆罗门是寻思者、考察者。他以寻思推导，以考察竞逐，由自身理解，便如是说：『此世间既非有边，亦非无边。凡沙门婆罗门作如是说「此世间具有边、具有边界」，彼等虚妄，凡沙门婆罗门作如是说「此世间无边、无边界」，彼等亦虚妄，凡沙门婆罗门作如是说「此世间既有边又无边」，彼等亦虚妄，此世间既非有边，亦非无边。』}

\textbf{诸比丘！此即第四根由，凭藉此、就此，一些沙门婆罗门为有边无边者，施设世间之有边无边。}

\begin{enumerate}\item \textbf{对世间住于有边想}，即不增长似相至轮围边际，执持其为「世间」，对世间住于有边想。而以轮围为边际增长似相，则为\textbf{无边想}。不增长上下，增长四旁，则为\textbf{上下有边想，四旁无边想}。寻思论当知仍如前述。且此四者亦唯以忆念自身先前所见而执持见故，归于前际想。\end{enumerate}

{\Large 58.} \textbf{藉此，诸比丘！这些沙门婆罗门为有边无边者，以四事施设世间之有边无边。诸比丘！因为任何沙门或婆罗门为有边无边者，施设世间之有边无边，他们全都以此四事，或以其中某个，无外于此。}

{\Large 59.} \textbf{诸比丘！如来了知此：『这些见的根由被如是执持、如是执取，便成如是之趣、如是之将来。』如来了知此，还了知较此更上者，且不执取此了知，由不执取，便独自知晓寂灭。如实知晓了诸受的集、没、味、患、离而无取解脱者，诸比丘！即如来。}

{\Large 60.} \textbf{这些，诸比丘！诸法甚深、难见、难觉、寂静、胜妙、非寻思所行、微妙、智者所经验，如来经自身证知、作证而宣说彼等，藉此，说着如来如实赞誉者才能正确地说。}

\subsection{无尽闪烁论者}

{\Large 61.} \textbf{诸比丘！有一些沙门婆罗门为无尽闪烁者，当在各处被提问时，以四事言辞闪烁、无尽闪烁。那些沙门婆罗门尊者凭藉什么、就什么而为无尽闪烁者，当在各处被提问时，以四事言辞闪烁、无尽闪烁？}

\begin{enumerate}\item 不死即\textbf{无尽}。这是什么？即以「如是于我则否」等方法，持见者的没有底线的见和言辞。种种散乱即\textbf{闪烁}，以无尽的见和言辞而闪烁，即无尽闪烁，有此等者，即\textbf{无尽闪烁者}。
\item 另一法。无尽为一鱼种之名\footnote{词典或解释作鳗鱼。}，它以出没等方式在水中穿梭，难被捕获，如是，此论亦四处穿梭，以不能靠近把握，被称为无尽闪烁，有此等者，即无尽闪烁者。\end{enumerate}

{\Large 62.} \textbf{于此，诸比丘！部分沙门或婆罗门不如实了知『此是善、此是不善』。他便如是想：『我不如实了知「此是善、此是不善」，若当我不如实了知「此是善、此是不善」，却去解答「此是善、此是不善」，我便成虚妄。令我虚妄者，便成我的困扰。令我困扰者，便成我的障碍。』如此，他由畏惧妄语、嫌厌妄语，既不解答『此是善』，也不解答『此是不善』，当在各处被提问时，言辞闪烁、无尽闪烁：『如是于我则否，这样于我则否，别样于我则否，不然于我则否，非不然于我则否。』}

\textbf{诸比丘！此即第一根由，凭藉此、就此，一些沙门婆罗门为无尽闪烁者，当在各处被提问时，言辞闪烁、无尽闪烁。}

\begin{enumerate}\item \textbf{不如实了知「此是善」}，意即不如实了知十善业道。\textbf{不善}则意指十不善业道。\textbf{便成我的困扰}，即因以「我所说为虚妄」生起后悔故，便成我的困扰，意即将成为苦。\textbf{便成我的障碍}，即成为我天界或道的障碍。\textbf{畏惧妄语、嫌厌妄语}，即对妄语既愧又惭。\textbf{言辞闪烁}，即以言辞闪烁。像什么？无尽闪烁，意即没有底线的闪烁。
\item 在「如是于我则否」等中，\textbf{如是于我则否}，即不确定之闪烁。\textbf{这样于我则否}，即拒斥以「自我和世间是常」所说的常论。\textbf{别样于我则否}，即拒斥与常论别样而说的部分常论。\textbf{不然于我则否}，即拒斥以「如来死后不存在」所说的断论。\textbf{非不然于我则否}，即拒斥以「既非是，也非不是」所说的寻思论。
\item 而他自己当被问到「此是善还是不善」时，则不作解答。被问到「此是善」时，说「如是于我则否」。随后，说到「难道是不善」时，说「这样于我则否」。说到「难道与两者有别」时，说「别样于我则否」。随后，说到「三种都不是，你主张什么」时，说「不然于我则否」。随后，说到「难道非不然是你的主张」时，说「非不然于我则否」，如是唯陷入混乱，不立于一方。\end{enumerate}

{\Large 63.} \textbf{第二，沙门婆罗门尊者凭藉什么、就什么而为无尽闪烁者，当在各处被提问时，言辞闪烁、无尽闪烁？于此，诸比丘！部分沙门或婆罗门不如实了知『此是善、此是不善』。他便如是想：『我不如实了知…………却去解答「此是善、此是不善」，我对此有欲或贪、嗔或恚。若我对此有欲或贪、嗔或恚，我便有取著。令我取著者，便成我的困扰。令我困扰者，便成我的障碍。』如此，他由畏惧取著、嫌厌取著…………}

\textbf{诸比丘！此即第二根由，凭藉此、就此，一些沙门婆罗门为无尽闪烁者，当在各处被提问时，言辞闪烁、无尽闪烁。}

\begin{enumerate}\item \textbf{欲或贪}，意即尽管不了知，仍遽然说善为「善」、说不善为「不善」后，去问其他智者「我对某某如是解答，此是善解否」，当他们说「善解，贤首！你唯对善答作善，对不善答作不善」，便以「没有和我一样的智者」而对此有欲或贪。且此中欲为弱力的贪，贪为强力的贪。\textbf{嗔或恚}，意即说善为「不善」或不善为「善」后，去问其他智者，当他们说「你错答」时，以「我连这点也不知道」而对此有嗔或恚。此处，嗔为弱力的忿，恚为强力的忿。
\item \textbf{我便有取著……便成我的困扰}，即这欲、贪二者是我的取著，嗔、恚二者是困扰。或者，两者以强烈的抓取为取著，以伤害为困扰。因为贪以不想放手而抓取所缘，好似水蛭，嗔以欲求消灭，好似毒蛇，且两者都以灼热之义，唯行伤害，被称为取著和困扰。其余仍与第一番相同。\end{enumerate}

{\Large 64.} \textbf{第三，沙门婆罗门尊者凭藉什么、就什么而为无尽闪烁者，当在各处被提问时，言辞闪烁、无尽闪烁？于此，诸比丘！部分沙门或婆罗门不如实了知『此是善、此是不善』。他便如是想：『我不如实了知…………却去解答「此是善、此是不善」。但有一些沙门婆罗门智者微妙，研究过他人的论点，好似劈开毫尖者，我想他们在劈时，以慧周遍诸见，他们对此会与我诘难、追究、商议。若他们对此与我诘难、追究、商议，我便不能解答。而我不能解答他们，便成我的困扰。令我困扰者，便成我的障碍。』如此，他由畏惧诘问、嫌厌诘问…………}

\textbf{诸比丘！此即第三根由，凭藉此、就此，一些沙门婆罗门为无尽闪烁者，当在各处被提问时，言辞闪烁、无尽闪烁。}

\begin{enumerate}\item \textbf{智者}，即具足智性。\textbf{微妙}，即有精致、微细的觉知，能通达微细之义的差别。\textbf{研究过他人的论点}，即已知晓他人的论点，且与他人探讨过论点。\textbf{好似劈开毫尖者}，即如同劈开毫尖的射手。\textbf{我想他们在劈时}，意即如同劈开毫尖者之于毫发，即便他人的见极其微细，也能\textbf{以}自己的\textbf{慧}劈开、\textbf{周遍诸见}。
\item \textbf{他们对此会与我}，即这些沙门婆罗门会与我对这些善、不善。\textbf{诘难}，即以「什么是善，什么是不善，请说你自己的主张」问主张。\textbf{追究}，即当说了「是此」时，以「因何原因执持此义」问原因。\textbf{商议}，即当说了「因此原因」时，先显明原因的过失，再以「你不知此却执持此，请回答此」而作诘难。\textbf{我便不能解答}，即我便不能圆满，意即不能圆满地谈论。\textbf{便成我的困扰}，即尽管反复地说，却不能解答，便成我的困扰，意即唯令唇颚舌喉干渴受苦而已。此中其余仍与第一番相同。\end{enumerate}

{\Large 65.} \textbf{第四，沙门婆罗门尊者凭藉什么、就什么而为无尽闪烁者，当在各处被提问时，言辞闪烁、无尽闪烁？于此，诸比丘！部分沙门或婆罗门迟钝、愚钝。他由迟钝、愚钝故，当在各处被提问时，言辞闪烁、无尽闪烁。}

\textbf{假设你问我『存在另一世间吗』，假设我认为『存在另一世间』，我应对你解答『存在另一世间』，但却说『如是于我则否，这样于我则否，别样于我则否，不然于我则否，非不然于我则否』。}\\
\textbf{……『不存在另一世间吗』……}\\
\textbf{……『存在且不存在另一世间吗』……}\\
\textbf{……『非存在非不存在另一世间吗』……}\\
\textbf{……『存在化生有情吗』……}\\
\textbf{……『不存在化生有情吗』……}\\
\textbf{……『存在且不存在化生有情吗』……}\\
\textbf{……『非存在非不存在化生有情吗』……}\\
\textbf{……『存在善恶之业的果、异熟吗』……}\\
\textbf{……『不存在善恶之业的果、异熟吗』……}\\
\textbf{……『存在且不存在善恶之业的果、异熟吗』……}\\
\textbf{……『非存在非不存在善恶之业的果、异熟吗』……}\\
\textbf{……『如来死后存在吗』……}\\
\textbf{……『如来死后不存在吗』……}\\
\textbf{……『如来死后存在且不存在吗』……}\\
\textbf{……『如来死后非存在非不存在吗』……}

\textbf{诸比丘！此即第四根由，凭藉此、就此，一些沙门婆罗门为无尽闪烁者，当在各处被提问时，言辞闪烁、无尽闪烁。}

\begin{enumerate}\item \textbf{迟钝}，即钝慧，唯无慧之名。\textbf{愚钝}，即极度愚昧。在「\textbf{如来存在}」等中，「如来」意指有情。此中余皆自明。且此四者亦唯以忆念先前转起之法而执持见故，归于前际想。\end{enumerate}

{\Large 66.} \textbf{藉此，诸比丘！这些沙门婆罗门为无尽闪烁者，当在各处被提问时，以四事言辞闪烁、无尽闪烁。诸比丘！因为任何沙门或婆罗门为无尽闪烁者，当在各处被提问时，言辞闪烁、无尽闪烁，他们全都以此四事，或以其中某个，无外于此……藉此，说着如来如实赞誉者才能正确地说。}

\subsection{偶然发生论者}

{\Large 67.} \textbf{诸比丘！有一些沙门婆罗门为偶然发生者，以二事施设自我和世间是偶然发生。那些沙门婆罗门尊者凭藉什么、就什么而为偶然发生者，以二事施设自我和世间是偶然发生？}

\begin{enumerate}\item 「自我和世间是偶然发生」之知见，即偶然发生，有此等者，即\textbf{偶然发生者}。偶然发生，即无原因之发生。\end{enumerate}

{\Large 68.} \textbf{诸比丘！有一些天人名为无想有情。而由想的生起，这些天人从彼众中亡殁。有这可能，诸比丘！某一有情从彼众中殁后，来到此间。当来到此间时，从家出至非家。当从家出至非家时，随着热勤、随着精勤、随着致力、随着不放逸、随着正作意，触证这样的心定，当心如是等持时，忆念想的生起，未忆念此前。}

\textbf{他便如是说：『自我和世间是偶然发生。其因为何？我先前不存在，我经历了不存在，现今变成了存在。』}

\textbf{诸比丘！此即第一根由，凭藉此、就此，一些沙门婆罗门为偶然发生者，施设自我和世间是偶然发生。}

\begin{enumerate}\item \textbf{无想有情}，这是开示的纲目，意即没有心的生起，仅仅是色的自体。当知他们的缘起如是——有些人在外道处出家后，于风遍作了遍作，转起第四禅，从禅出起后，看见心的过失「当有心时，便有断手等之苦，以及所有怖畏，这心真是够了，唯无心之状为寂静」，如是看到心的过失后，以不退失禅那而死，在无想有情中投生，心便随其死心之灭，在此停止，唯在彼处显现色蕴。正好比被弓弦之势射出的箭，随弓弦之势而在空中飞行，如是，他们在彼处被禅那之势射出，投生之后，随禅那之势而停留如许时间，但当禅那之势退失时，色蕴便在彼处消隐，而在此处生起结生想。因为以这在此处生起的想，他们在彼处之死可知，所以说「\textbf{而由想的生起，这些天人从彼众中亡殁}」。\textbf{存在}，即存在之状。此中其余自明。\end{enumerate}

{\Large 69.} \textbf{第二，沙门婆罗门尊者凭藉什么、就什么而为偶然发生者，施设自我和世间是偶然发生？于此，诸比丘！部分沙门或婆罗门是寻思者、考察者。他以寻思推导，以考察竞逐，由自身理解，便如是说：『自我和世间是偶然发生。』}

\textbf{诸比丘！此即第二根由，凭藉此、就此，一些沙门婆罗门为偶然发生者，施设自我和世间是偶然发生。}

\begin{enumerate}\item 寻思论当知仍同前述之法。\end{enumerate}

{\Large 70.} \textbf{藉此，诸比丘！这些沙门婆罗门为偶然发生者，以二事施设自我和世间是偶然发生。诸比丘！因为任何沙门或婆罗门为偶然发生者，施设自我和世间是偶然发生，他们全都以此二事，或以其中某个，无外于此……藉此，说着如来如实赞誉者才能正确地说。}

{\Large 71.} \textbf{藉此，诸比丘！这些沙门婆罗门为前际想者、前际随见者，就前际以十八事表达多样胜解之句。诸比丘！因为任何沙门或婆罗门为前际想者、前际随见者，就前际表达多样胜解之句，他们全都以此十八事，或以其中某个，无外于此。}

{\Large 72.} \textbf{诸比丘！如来了知此：『这些见的根由被如是执持、如是执取，便成如是之趣、如是之将来。』如来了知此，还了知较此更上者，且不执取此了知，由不执取，便独自知晓寂灭。如实知晓了诸受的集、没、味、患、离而无取解脱者，诸比丘！即如来。}

{\Large 73.} \textbf{这些，诸比丘！诸法甚深、难见、难觉、寂静、胜妙、非寻思所行、微妙、智者所经验，如来经自身证知、作证而宣说彼等，藉此，说着如来如实赞誉者才能正确地说。}

\section{后际想者}

{\Large 74.} \textbf{诸比丘！有一些沙门婆罗门为后际想者、后际随见者，就后际以四十四事表达多样胜解之句。那些沙门婆罗门尊者凭藉什么、就什么而为后际想者、后际随见者，就后际以四十四事表达多样胜解之句？}

\begin{enumerate}\item 如是，在显示了十八前际想者后，现在，为显示四十四后际想者，说了「诸比丘！有一些……」等等。这里，设想后际而执持者，为\textbf{后际想者}，或有后际想者为后际想者。如是当知其余也同先前所说品类之法。\end{enumerate}

\subsection{有想论者}

{\Large 75.} \textbf{诸比丘！有一些沙门婆罗门为死后有想论者，以十六事施设自我死后有想。那些沙门婆罗门尊者凭藉什么、就什么而为死后有想论者，以十六事施设自我死后有想？}

\begin{enumerate}\item \textbf{死后者}，死亡被称为死，他们说自我为死后，即「死后者」。转起「有想」之论，即有想论，有此等者，为\textbf{有想论者}。\end{enumerate}

{\Large 76.} \textbf{他们施设其为：『自我有色，死后无病、有想。』}\\
\textbf{他们施设其为：『自我无色，死后无病、有想。』}\\
\textbf{……『自我有色且无色……』}\\
\textbf{……『自我非有色非无色……』}\\
\textbf{……『自我有边……』}\\
\textbf{……『自我无边……』}\\
\textbf{……『自我有边且无边……』}\\
\textbf{……『自我非有边非无边……』}\\
\textbf{……『自我为同一想……』}\\
\textbf{……『自我为差异想……』}\\
\textbf{……『自我为有限想……』}\\
\textbf{……『自我为无量想……』}\\
\textbf{……『自我一向是乐……』}\\
\textbf{……『自我一向是苦……』}\\
\textbf{……『自我有苦有乐……』}\\
\textbf{他们施设其为：『自我非苦非乐，死后无病、有想。』}\\

\begin{enumerate}\item 在「自我有色」等中，执取遍处之色为自我，且对此转起之想为想，或者如邪命者等仅以寻思，施设其为「自我有色，死后无病、有想」。这里，\textbf{无病}，即恒常。
\item 而执取无色定之相为自我，且定想为想，或者如尼乾陀等仅以寻思，施设其为「自我无色，死后无病、有想」。第三见以混合执取而转起。第四则仅以寻思。
\item 第二组的四个，当知如「有边无边论」所述。
\item 在第三组的四个中，当知依入定者为\textbf{同一想}，依未入定者为\textbf{差异想}，依有限的遍处为\textbf{有限想}，依广大的遍处为\textbf{无量想}。
\item 而在第四组的四个中，以天眼看到转生于三、四禅地\footnote{三、四禅地：\textbf{长部戒蕴品新疏}说，就四法和五法而说三、四禅地 \textit{catukkanayaṃ pañcakanayañ ca sandhāya tika-catukka-jjhānabhūmiyan ti vuttaṃ}。}，便执持\textbf{一向是乐}。看到转生于地狱，则执持\textbf{一向是苦}。看到转生于人，则执持\textbf{有苦有乐}。看到转生于广果天，则执持\textbf{非苦非乐}。差别在于，得宿住随念智者成为前际想者，天眼者成为后际想者。\end{enumerate}

{\Large 77.} \textbf{藉此，诸比丘！这些沙门婆罗门为死后有想论者，以十六事施设自我死后有想。诸比丘！因为任何沙门或婆罗门为死后有想论者，施设自我死后有想，他们全都以此十六事，或以其中某个，无外于此……藉此，说着如来如实赞誉者才能正确地说。}

\subsection{无想论者}

{\Large 78.} \textbf{诸比丘！有一些沙门婆罗门为死后无想论者，以八事施设自我死后无想。那些沙门婆罗门尊者凭藉什么、就什么而为死后无想论者，以八事施设自我死后无想？}

\begin{enumerate}\item \textbf{无想论}，依在「有想论」开头所说的两组四个可知。\textbf{非有想非无想论}亦然。只是彼处为执持「自我有想」者的见，此处为执持「无想」及「非有想非无想」者的。这里，不应一味地寻觅原因，因为持见者的执取，如同疯人的篮子一般\footnote{疯人的篮子：见第 31 段的义注。}。\end{enumerate}

{\Large 79.} \textbf{他们施设其为：『自我有色，死后无病、无想。』}\\
\textbf{他们施设其为：『自我无色，死后无病、无想。』}\\
\textbf{……『自我有色且无色……』}\\
\textbf{……『自我非有色非无色……』}\\
\textbf{……『自我有边……』}\\
\textbf{……『自我无边……』}\\
\textbf{……『自我有边且无边……』}\\
\textbf{他们施设其为：『自我非有边非无边，死后无病、无想。』}

{\Large 80.} \textbf{藉此，诸比丘！这些沙门婆罗门为死后无想论者，以八事施设自我死后无想。诸比丘！因为任何沙门或婆罗门为死后无想论者，施设自我死后无想，他们全都以此八事，或以其中某个，无外于此……藉此，说着如来如实赞誉者才能正确地说。}

\subsection{非有想非无想论者}

{\Large 81.} \textbf{诸比丘！有一些沙门婆罗门为死后非有想非无想论者，以八事施设自我死后非有想非无想。那些沙门婆罗门尊者凭藉什么、就什么而为死后非有想非无想论者，以八事施设自我死后非有想非无想？}

{\Large 82.} \textbf{他们施设其为：『自我有色，死后无病、非有想非无想。』}\\
\textbf{……『自我无色……』}\\
\textbf{……『自我有色且无色……』}\\
\textbf{……『自我非有色非无色……』}\\
\textbf{……『自我有边……』}\\
\textbf{……『自我无边……』}\\
\textbf{……『自我有边且无边……』}\\
\textbf{他们施设其为：『自我非有边非无边，死后无病、非有想非无想。』}

{\Large 83.} \textbf{藉此，诸比丘！这些沙门婆罗门为死后非有想非无想论者，以八事施设自我死后非有想非无想。诸比丘！因为任何沙门或婆罗门为死后非有想非无想论者，施设自我死后非有想非无想，他们全都以此八事……藉此，说着如来如实赞誉者才能正确地说。}

\subsection{断论者}

{\Large 84.} \textbf{诸比丘！有一些沙门婆罗门为断论者，以七事施设现存有情的断绝、消亡、无有。那些沙门婆罗门尊者凭藉什么、就什么而为断论者，以七事施设现存有情的断、消亡、无有？}

\begin{enumerate}\item 在断论中，\textbf{现存}，即当前存在。\textbf{断绝}即断坏，\textbf{消亡}即不可见，\textbf{无有}即状态的离去，这一切都是彼此的异名。
\item 这里，有二人执持断见：有所得者和无所得者。有所得者以阿罗汉的天眼看到死，却不见投生，或者只能看到死，而非投生，他便执持断见。而无所得者，或以「谁知他世有否」而贪求欲乐，或以「好比叶子从树上掉落，不再生长，有情亦如是」等的寻思，而执持断。而在此处，当知依于爱、见，唯意向这样和别样，而生起此七见。\end{enumerate}

{\Large 85.} \textbf{于此，诸比丘！部分沙门或婆罗门如是论、如是见：『先生！由于这自我有色，属四大种，父母所生，身坏后便断绝、消亡，死后不存在，至此，先生！这自我便完全断绝。』如此，一些人施设现存有情的断绝、消亡、无有。}

\begin{enumerate}\item 这里，\textbf{有色}，即具有色。\textbf{属四大种}，即四大种所造。父母有此，即父母之所有。这是什么？精血。在父母之所有中生成、产生，即\textbf{父母所生}。如此，以色身为纲目，说人的自体为「自我」。\textbf{如此，一些人}，意即「如是，一些人」。\end{enumerate}

{\Large 86.} \textbf{另有人如是说：『先生！存在这自我，你所说的，我不说它不存在，但是，先生！这自我并不至此便完全断绝。先生！存在另一自我，天界，有色，欲界，段食为食。你对其不知不见，而我对其知见。先生！这自我由于身坏后便断绝、消亡，死后不存在，至此，先生！这自我便完全断绝。』如此，一些人施设现存有情的断绝、消亡、无有。}

\begin{enumerate}\item 第二（见）在拒斥此之后，说天的自体。\textbf{天界}，即在天界中生成。\textbf{欲界}，即系属于六欲界天。吃段食之食，即\textbf{段食为食}。\end{enumerate}

{\Large 87.} \textbf{另有人如是说：『先生！存在这自我，你所说的，我不说它不存在，但是，先生！这自我并不至此便完全断绝。先生！存在另一自我，天界，有色，由意所成，肢体俱全，诸根无缺。你对其不知不见，而我对其知见…………}

\begin{enumerate}\item \textbf{由意所成}，即由禅那之意而转生。\textbf{肢体俱全}，即具备一切肢体。\textbf{诸根无缺}，即诸根圆满。这是依梵界之中的存在以及依其他的形态而说的。\end{enumerate}

{\Large 88.} \textbf{另有人如是说：『先生！存在这自我，你所说的，我不说它不存在，但是，先生！这自我并不至此便完全断绝。先生！存在另一自我，由完全超越色想，由有对想的隐没，由种种想的不作意，以「虚空无尽」经验到空无边处。你对其不知不见，而我对其知见…………}

\begin{enumerate}\item \textbf{由完全超越色想}等之义，在「清净道论」中已说\footnote{见\textbf{清净道论}·说无色品第十第 12 段及以下。}。而在「\textbf{经验到空无边处}」等处，当知其义为经验到空无边处之有。此中其余自明。\end{enumerate}

{\Large 89.} \textbf{另有人如是说：『先生！存在这自我，你所说的，我不说它不存在，但是，先生！这自我并不至此便完全断绝。先生！存在另一自我，完全超越了空无边处，以「识无尽」经验到识无边处。你对其不知不见，而我对其知见…………}

{\Large 90.} \textbf{另有人如是说：『先生！存在这自我，你所说的，我不说它不存在，但是，先生！这自我并不至此便完全断绝。先生！存在另一自我，完全超越了识无边处，以「无有任何」经验到无所有处。你对其不知不见，而我对其知见…………}

{\Large 91.} \textbf{另有人如是说：『先生！存在这自我，你所说的，我不说它不存在，但是，先生！这自我并不至此便完全断绝。先生！存在另一自我，完全超越了无所有处，以「此为寂静、此为胜妙」经验到非想非非想处。你对其不知不见，而我对其知见…………}

{\Large 92.} \textbf{藉此，诸比丘！这些沙门婆罗门为断论者，以七事施设现存有情的断绝、消亡、无有。诸比丘！因为任何沙门或婆罗门为断论者，施设现存有情的断绝、消亡、无有，他们全都以此七事……藉此，说着如来如实赞誉者才能正确地说。}

\subsection{现法涅槃论者}

{\Large 93.} \textbf{诸比丘！有一些沙门婆罗门为现法涅槃论者，以五事施设现存有情的最上现法涅槃。那些沙门婆罗门尊者凭藉什么、就什么而为现法涅槃论者，以五事施设现存有情的最上现法涅槃？}

\begin{enumerate}\item 在现法涅槃论中，现法，是说现量之法，即在各处所得自体的同义语。在现法中涅槃，即现法涅槃，意即就在此自体中止息了苦。作此论说者，即\textbf{现法涅槃论者}。\textbf{最上现法涅槃}，即最上的现法涅槃，意即最高。\end{enumerate}

{\Large 94.} \textbf{于此，诸比丘！部分沙门或婆罗门如是论、如是见：『先生！由于这自我被付诸、据有种种五欲而娱乐，至此，先生！这自我便到达最上现法涅槃。』如此，一些人施设现存有情的最上现法涅槃。}

\begin{enumerate}\item \textbf{种种五欲}，即可意之色等的五种爱欲或束缚。\textbf{被付诸}，即被完全吸引、黏著。\textbf{据有}，即具足。\textbf{娱乐}，即在这些种种爱欲中随其所乐地放肆、极放肆诸根，处处供给。或者，即游戏、喜乐、嬉戏。
\item 且此中有两种爱欲：人类的和天界的。人类的，当视之如曼达多的种种爱欲，天界的，则如他化自在天天王的种种爱欲。因为对经验这般的爱欲者，他们才施设现法涅槃的成就。\end{enumerate}

{\Large 95.} \textbf{另有人如是说：『先生！存在这自我，你所说的，我不说它不存在，但是，先生！这自我并不至此便到达最上现法涅槃。其因为何？因为，先生！爱欲为无常、苦、变易法，由彼等的变易、别样之状，生起忧悲苦忧愁。先生！由于这自我已离爱欲，已离不善法，进入并住于有寻有伺、因远离而生喜乐的初禅，至此，先生！这自我便到达最上现法涅槃。』如此，一些人施设现存有情的最上现法涅槃。}

\begin{enumerate}\item 在第二番中，当知以存在后的不存在之义为\textbf{无常}，以逼迫之义为\textbf{苦}，以舍弃本然之义为\textbf{变易法}。\textbf{由彼等的变易、别样之状}，即由被称为彼等爱欲的变易的别样之状。以「那曾是我的，我却已无有」所说之法，\textbf{生起忧悲苦忧愁}。
\item 这里，内里焦灼之相为\textbf{忧}，依于此的悲叹之相为\textbf{悲}，身体逼迫之相为\textbf{苦}，心意困扰之相为\textbf{忧}，消沉之相为\textbf{愁}。「\textbf{已离爱欲}」等之义，在在「清净道论」中已说\footnote{见\textbf{清净道论}·说地遍品第四第 79 段及以下。}。\end{enumerate}

{\Large 96.} \textbf{另有人如是说：『先生！存在这自我，你所说的，我不说它不存在，但是，先生！这自我并不至此便到达最上现法涅槃。其因为何？在彼处仍有所寻所伺，它便因此被称为粗的。先生！由于这自我平息了寻伺，进入并住于内在洁净、心念专一、无寻无伺、因定而生喜乐的二禅，至此，先生！这自我便到达最上现法涅槃。』如此，一些人施设现存有情的最上现法涅槃。}

\begin{enumerate}\item \textbf{所寻}，即以植入转起的寻。\textbf{所伺}，即以细思转起的伺。\textbf{它便因此}，即这初禅便因此寻与伺，被认为是\textbf{粗的}，像有刺一般。\end{enumerate}

{\Large 97.} \textbf{另有人如是说：『先生！存在这自我，你所说的，我不说它不存在，但是，先生！这自我并不至此便到达最上现法涅槃。其因为何？在彼处仍有与喜俱行的心之飘然，它便因此被称为粗的。先生！由于这自我从喜离贪，放舍而住，念且正知，以身感受乐，进入并住于如圣者们宣说为「放舍、具念、乐住」的三禅，至此，先生！这自我便到达最上现法涅槃。』如此，一些人施设现存有情的最上现法涅槃。}

\begin{enumerate}\item \textbf{与喜俱行}，即喜。\textbf{心之飘然}，即令心起飘然之状。\end{enumerate}

{\Large 98.} \textbf{另有人如是说：『先生！存在这自我，你所说的，我不说它不存在，但是，先生！这自我并不至此便到达最上现法涅槃。其因为何？在彼处仍有「乐」为心的受用，它便因此被称为粗的。先生！由于这自我已断除乐，已断除苦，由先前的喜忧的隐没，进入并住于无苦无乐、舍念遍净的四禅，至此，先生！这自我便到达最上现法涅槃。』如此，一些人施设现存有情的最上现法涅槃。}

\begin{enumerate}\item \textbf{心的受用}，即从禅那出起，心对此乐反复地受用、作意、存念。在「现法涅槃论」中，其余自明。
\item 至此，所有六十二见均已论述。其中，仅七为断见，其余为常见。\end{enumerate}

{\Large 99.} \textbf{藉此，诸比丘！这些沙门婆罗门为现法涅槃论者，以五事施设现存有情的最上现法涅槃。诸比丘！因为任何沙门或婆罗门为现法涅槃论者，施设现存有情的最上现法涅槃，他们全都以此五事……藉此，说着如来如实赞誉者才能正确地说。}

{\Large 100.} \textbf{藉此，诸比丘！这些沙门婆罗门为后际想者、后际随见者，就后际以四十四事表达多样胜解之句。诸比丘！因为任何沙门或婆罗门为后际想者、后际随见者，就后际表达多样胜解之句，他们全都以此四十四事……藉此，说着如来如实赞誉者才能正确地说。}

{\Large 101.} \textbf{藉此，诸比丘！这些沙门婆罗门为前际想者、后际想者、前后际想者、前后际随见者，就前后际以六十二事表达多样胜解之句。}

\begin{enumerate}\item 现在，以此番之「藉此，诸比丘！这些」总结了所有后际想者，解答了一切知智，又以「藉此，诸比丘！这些」等番总结了所有这些前后际想者，仍然解答此智。如此，在「诸比丘！什么是这些……诸法」等处提问也罢，问了一切知智后解答也罢，如同以秤衡量有情的意乐，如同从须弥山脚下提取沙子，在提取了六十二见后，仍然解答一切知智。如是，这开示是以次第为连接而展开。
\item 因为有三种经的连接：以问为连接，以意乐为连接，以次第为连接。这里，\begin{quoting}如是说已，某比丘便对世尊说：「尊者！何为此岸？何为彼岸？何为在中间沉沦？何为在高地隆起？何为被人执取？何为被非人执取？何为被回流执取？何为内里腐败？」{\bluefootnotesize 〔相应部第 4.241 段〕}\end{quoting}如是，依世尊对所问的解答之经，当知为\textbf{以问为连接}。\begin{quoting}于是，某比丘心中便起如是寻思：「据说，色无我，受、想、行、识无我，那么，无我之业会触及哪个自我呢？」于是，世尊以心了知此比丘心中的寻思，便告诸比丘：「诸比丘！有这可能，于此，一些痴人无知，无明，被爱主导，过度推导大师的教法，会认为『据说，色无我……触及哪个自我呢？』于汝意云何？诸比丘！色是常还是无常？」{\bluefootnotesize 〔中部第 3.90 段〕}\end{quoting}如是，依世尊在知晓他人的意乐后所说的经，当知为\textbf{以意乐为连接}。
\item 而在起始处以某法被提出的开示，则对于此法，或以顺适之法，或以相反，在这些经中展开更多开示，当知依于彼等，为\textbf{以次第为连接}。例如：在「若他希望经」中，先以戒提出开示，而后展开六神通。在「布经」中，先以烦恼提出开示，而后展开梵住。在「㤭赏弥经」中，先以争吵提出，而后展开应忆念之法。在「锯喻经」中，先以不可忍耐提出，而后展开锯喻。在此「梵网」中，也是先以见提出开示，而后展开「空性阐释」。因此说「如是，这开示是以次第为连接而展开」。\end{enumerate}

{\Large 102.} \textbf{诸比丘！因为任何沙门或婆罗门为前际想者、后际想者、前后际想者、前后际随见者，就前后际表达多样胜解之句，他们全都以此六十二事，或以其中某个，无外于此。}

{\Large 103.} \textbf{诸比丘！如来了知此：『这些见的根由被如是执持、如是执取，便成如是之趣、如是之将来。』如来了知此，还了知较此更上者，且不执取此了知，由不执取，便独自知晓寂灭。如实知晓了诸受的集、没、味、患、离而无取解脱者，诸比丘！即如来。}

{\Large 104.} \textbf{这些，诸比丘！诸法甚深、难见、难觉、寂静、胜妙、非寻思所行、微妙、智者所经验，如来经自身证知、作证而宣说彼等，藉此，说着如来如实赞誉者才能正确地说。}

\section{施设自我和世间之事由}

\subsection{不安与战栗节}

{\Large 105.} \textbf{此处，诸比丘！凡沙门婆罗门为常论者，以四事施设自我和世间是常，这只是这些沙门婆罗门尊者不知、不见、与渴爱俱行的感受，只是不安与战栗。}

\begin{enumerate}\item 现在，为了显示界限与区分，发起以「此处，诸比丘」为起始的开示。\textbf{这只是这些沙门婆罗门尊者不知、不见、与渴爱俱行的感受，只是不安与战栗}，即因见之味、见之乐、见之感受，他们生起喜悦，以四事施设自我和世间是常，这只是这些沙门婆罗门尊者不如实知见诸法自性的感受，只是与渴爱俱行、唯与渴爱俱行的感受，且这只是不安与战栗。因被称为见和爱的不安而战栗、动摇、震动，和埋在稃堆里的树桩一样，显示非如须陀洹之知见般不动。于「部分常见论」等中，亦同此理。\end{enumerate}

{\Large 106.} \textbf{此处，诸比丘！凡沙门婆罗门为部分常、部分非常者，以四事施设自我和世间部分是常、部分非常，这只是……不安与战栗。}

{\Large 107.} \textbf{此处，诸比丘！凡沙门婆罗门为有边无边者，以四事施设世间之有边无边，这只是……不安与战栗。}

{\Large 108.} \textbf{此处，诸比丘！凡沙门婆罗门为无尽闪烁者，当在各处被提问时，以四事言辞闪烁、无尽闪烁，这只是……不安与战栗。}

{\Large 109.} \textbf{此处，诸比丘！凡沙门婆罗门为偶然发生者，以二事施设自我和世间是偶然发生，这只是……不安与战栗。}

{\Large 110.} \textbf{此处，诸比丘！凡沙门婆罗门为前际想者、前际随见者，就前际以十八事表达多样胜解之句，这只是……不安与战栗。}

{\Large 111.} \textbf{此处，诸比丘！凡沙门婆罗门为死后有想论者，以十六事施设自我死后有想，这只是……不安与战栗。}

{\Large 112.} \textbf{此处，诸比丘！凡沙门婆罗门为死后无想论者，以八事施设自我死后无想，这只是……不安与战栗。}

{\Large 113.} \textbf{此处，诸比丘！凡沙门婆罗门为死后非有想非无想论者，以八事施设自我死后非有想非无想，这只是……不安与战栗。}

{\Large 114.} \textbf{此处，诸比丘！凡沙门婆罗门为断论者，以七事施设现存有情的断绝、消亡、无有，这只是……不安与战栗。}

{\Large 115.} \textbf{此处，诸比丘！凡沙门婆罗门为现法涅槃论者，以五事施设现存有情的最上现法涅槃，这只是……不安与战栗。}

{\Large 116.} \textbf{此处，诸比丘！凡沙门婆罗门为后际想者、后际随见者，就后际以四十四事表达多样胜解之句，这只是……不安与战栗。}

{\Large 117.} \textbf{此处，诸比丘！凡沙门婆罗门为前际想者、后际想者、前后际想者、前后际随见者，就前后际以六十二事表达多样胜解之句，这只是……不安与战栗。}

\subsection{以触为缘节}

{\Large 118.} \textbf{此处，诸比丘！凡沙门婆罗门为常论者，以四事施设自我和世间是常，这只是以触为缘。}

\begin{enumerate}\item 复次，为了显示辗转之缘，发起「此处，诸比丘！凡沙门婆罗门为常论者」等等。这里，\textbf{这只是以触为缘}，即因见之味、见之乐、见之感受，他们生起喜悦，以四事施设自我和世间是常，显示这被爱、见所战栗的感受，只是以触为缘。在一切处均同此理。\end{enumerate}

{\Large 119.} \textbf{此处，诸比丘！凡沙门婆罗门为部分常、部分非常者，以四事施设自我和世间部分是常、部分非常，这只是以触为缘。}

{\Large 120.} \textbf{此处，诸比丘！凡沙门婆罗门为有边无边者，以四事施设世间之有边无边，这只是以触为缘。}

{\Large 121.} \textbf{此处，诸比丘！凡沙门婆罗门为无尽闪烁者，当在各处被提问时，以四事言辞闪烁、无尽闪烁，这只是以触为缘。}

{\Large 122.} \textbf{此处，诸比丘！凡沙门婆罗门为偶然发生者，以二事施设自我和世间是偶然发生，这只是以触为缘。}

{\Large 123.} \textbf{此处，诸比丘！凡沙门婆罗门为前际想者、前际随见者，就前际以十八事表达多样胜解之句，这只是以触为缘。}

{\Large 124.} \textbf{此处，诸比丘！凡沙门婆罗门为死后有想论者，以十六事施设自我死后有想，这只是以触为缘。}

{\Large 125.} \textbf{此处，诸比丘！凡沙门婆罗门为死后无想论者，以八事施设自我死后无想，这只是以触为缘。}

{\Large 126.} \textbf{此处，诸比丘！凡沙门婆罗门为死后非有想非无想论者，以八事施设自我死后非有想非无想，这只是以触为缘。}

{\Large 127.} \textbf{此处，诸比丘！凡沙门婆罗门为断论者，以七事施设现存有情的断绝、消亡、无有，这只是以触为缘。}

{\Large 128.} \textbf{此处，诸比丘！凡沙门婆罗门为现法涅槃论者，以五事施设现存有情的最上现法涅槃，这只是以触为缘。}

{\Large 129.} \textbf{此处，诸比丘！凡沙门婆罗门为后际想者、后际随见者，就后际以四十四事表达多样胜解之句，这只是以触为缘。}

{\Large 130.} \textbf{此处，诸比丘！凡沙门婆罗门为前际想者、后际想者、前后际想者、前后际随见者，就前后际以六十二事表达多样胜解之句，这只是以触为缘。}

\subsection{无有是处节}

{\Large 131.} \textbf{此处，诸比丘！凡沙门婆罗门为常论者，以四事施设自我和世间是常，他们若离于触而能感受，无有是处。}

\begin{enumerate}\item 现在，为了显示此缘在见的感受中的有力之相，再次说了「此处，诸比丘！凡沙门婆罗门为常论者」等等。这里，\textbf{他们若离于触}，即这些沙门婆罗门若没有触而\textbf{能感受}此受，无有此因。因为好比支撑正在倒塌的房舍的柱子，为有力之缘，无此柱的支撑，则不能存立，如是，触也是受的有力之缘，显示没有它，便没有此见的感受。在一切处均同此理。\end{enumerate}

{\Large 132.} \textbf{此处，诸比丘！凡沙门婆罗门为部分常、部分非常者，以四事施设自我和世间部分是常、部分非常，他们若离于触而能感受，无有是处。}

{\Large 133.} \textbf{此处，诸比丘！凡沙门婆罗门为有边无边者，以四事施设世间之有边无边，他们若离于触而能感受，无有是处。}

{\Large 134.} \textbf{此处，诸比丘！凡沙门婆罗门为无尽闪烁者，当在各处被提问时，以四事言辞闪烁、无尽闪烁，他们若离于触而能感受，无有是处。}

{\Large 135.} \textbf{此处，诸比丘！凡沙门婆罗门为偶然发生者，以二事施设自我和世间是偶然发生，他们若离于触而能感受，无有是处。}

{\Large 136.} \textbf{此处，诸比丘！凡沙门婆罗门为前际想者、前际随见者，就前际以十八事表达多样胜解之句，他们若离于触而能感受，无有是处。}

{\Large 137.} \textbf{此处，诸比丘！凡沙门婆罗门为死后有想论者，以十六事施设自我死后有想，他们若离于触而能感受，无有是处。}

{\Large 138.} \textbf{此处，诸比丘！凡沙门婆罗门为死后无想论者，以八事施设自我死后无想，他们若离于触而能感受，无有是处。}

{\Large 139.} \textbf{此处，诸比丘！凡沙门婆罗门为死后非有想非无想论者，以八事施设自我死后非有想非无想，他们若离于触而能感受，无有是处。}

{\Large 140.} \textbf{此处，诸比丘！凡沙门婆罗门为断论者，以七事施设现存有情的断绝、消亡、无有，他们若离于触而能感受，无有是处。}

{\Large 141.} \textbf{此处，诸比丘！凡沙门婆罗门为现法涅槃论者，以五事施设现存有情的最上现法涅槃，他们若离于触而能感受，无有是处。}

{\Large 142.} \textbf{此处，诸比丘！凡沙门婆罗门为后际想者、后际随见者，就后际以四十四事表达多样胜解之句，他们若离于触而能感受，无有是处。}

{\Large 143.} \textbf{此处，诸比丘！凡沙门婆罗门为前际想者、后际想者、前后际想者、前后际随见者，就前后际以六十二事表达多样胜解之句，他们若离于触而能感受，无有是处。}

\subsection{论基于见的流转}

{\Large 144.} \textbf{此处，诸比丘！凡沙门婆罗门为常论者，以四事施设自我和世间是常，}\\
\textbf{凡沙门婆罗门为部分常、部分非常者……}\\
\textbf{凡沙门婆罗门为有边无边者……}\\
\textbf{凡沙门婆罗门为无尽闪烁者……}\\
\textbf{凡沙门婆罗门为偶然发生者……}\\
\textbf{凡沙门婆罗门为前际想者……}\\
\textbf{凡沙门婆罗门为死后有想论者……}\\
\textbf{凡沙门婆罗门为死后无想论者……}\\
\textbf{凡沙门婆罗门为死后非有想非无想论者……}\\
\textbf{凡沙门婆罗门为断论者……}\\
\textbf{凡沙门婆罗门为现法涅槃论者……}\\
\textbf{凡沙门婆罗门为后际想者……}\\
\textbf{凡沙门婆罗门为前际想者、后际想者、前后际想者、前后际随见者，就前后际以六十二事表达多样胜解之句，他们全都由六触入处，经触、触而后感受，以彼等的受为缘而爱，以爱为缘而取，以取为缘而有，以有为缘而生，以生为缘而老死，生起忧悲苦忧愁。}

\begin{enumerate}\item 现在，以「\textbf{凡沙门婆罗门为常论者，以四事施设自我和世间是常，凡沙门婆罗门为部分常、部分非常者……}」等方法，混同所有见的感受。为什么？为了之后归入触中。如何？\textbf{他们全都由六触入处，经触、触而后感受}。这里，六触入处者，即眼触入处、耳触入处、鼻触入处、舌触入处、身触入处、意触入处等六。
\item 此「入处」字，用于出生、会合、原因、仅作施设等意义中。这里，\begin{quoting}剑浮阇为群马之处，南方为群牛（之处）。\end{quoting}即用于出生，意即出生之地。\begin{quoting}悦意之处，诸鸟同栖。{\bluefootnotesize 〔增支部第 5.38 段〕}\end{quoting}即用于会合。\begin{quoting}均在处中。{\bluefootnotesize 〔增支部第 3.102 段〕}\end{quoting}即用于原因。\begin{quoting}在林野处的茅屋内安住。{\bluefootnotesize 〔相应部第 1.255 段〕}\end{quoting}即用于仅作施设。此处，它适合出生等的三种意义。因为在眼等之中，以触为第五的诸法出生、会合，且这些入处即彼等的原因。但在此处，以\begin{quoting}缘眼、色生起眼识，三者的和合为触。{\bluefootnotesize 〔相应部第 2.43 段〕}\end{quoting}的方法，唯以触为首引入开示，为显示以触为始的缘的辗转，而说「触入处」等。
\item \textbf{经触、触而后感受}，此中，虽然像是在说诸入处的触的作用，但应知它们并不如此有触的作用。因为诸入处并不去触，唯是触在触彼彼所缘，而诸入处在比照触后才显现，所以，当知此中之义为：这一切由六触入处生成的触，经触了色等所缘而后感受此见的感受。
\item 在「\textbf{以彼等的受为缘而爱}」等中，受，即六触入处生成的受。它对色爱等类的爱，为亲依止方面的缘，因此说「以彼等的受为缘而爱」。而这（爱）对四种取，为亲依止方面与俱生方面的缘。取对有也同样。有对生，为亲依止方面的缘。而此中的\textbf{生}，当视为有变化的五蕴，生对老死以及忧等，为亲依止方面的缘。这于此为略说，其详在「清净道论·说缘起」中已述。此处则仅须知其目的。
\item 因为世尊在谈流转论时，或以无明为首，如\begin{quoting}诸比丘！不知无明的前际，诸比丘！便如是说：「此前没有无明，然后才生成。」但是可知「缘此而有无明」。{\bluefootnotesize 〔增支部第 10.61 段〕}\end{quoting}或以爱为首，如\begin{quoting}诸比丘！不知有爱的前际……但是可知「缘此而有有爱」。{\bluefootnotesize 〔增支部第 10.62 段〕}\end{quoting}或以见为首，如\begin{quoting}诸比丘！不知有见的前际……但是可知「缘此而有有见」。\end{quoting}而作论述。在此则以见为首而作论述，在论述了因受的贪染而生起的诸见后，便论述以受为根本的缘起。
\item 以此显示：如是，这些持见者在执持此知见后，于三有、四生、五趣、七识住、九有情居中，从此至彼、从彼至此地转世、轮回，像轭在器械上的牛，像栓在柱子上的狗，像随风飘逝的船，只是随转于流转之苦，连从流转之苦中抬起头来，也不能够。\end{enumerate}

\section{论还灭等}

{\Large 145.} \textbf{诸比丘！由于比丘如实了知六触入处的集、没、味、患、离，他便了知较此一切更上者。}

\begin{enumerate}\item 如是，在论述了基于持见者的流转后，现在，基于从事修行的比丘，为显示还灭，说了「诸比丘！由于比丘」等等。这里，\textbf{由于}，即当……之时。\textbf{六触入处}，即由此六触入处，持见者经触后而感受，发生轮回，即此六触入处。
\item 在「\textbf{集}」等中，当以「由无明集而眼集」等在受业处所说之法，而知触入处的集等。然而，在彼处说为「由触集、由触灭」，在此处，对于眼等，当知其为「由食集、由食灭」，对于意入处为「由名色集、由名色灭」。
\item \textbf{了知更上者}，即持见者唯了知见，而他却了知见，以及较见更上的戒、定、慧、解脱，乃至阿罗汉。谁如是了知？漏尽者了知，阿那含、斯陀含、须陀洹、多闻者、持书的比丘了知，开始作观者了知。而开示则以阿罗汉为顶点而终结。\end{enumerate}

{\Large 146.} \textbf{诸比丘！因为任何沙门或婆罗门为前际想者、后际想者、前后际想者、前后际随见者，就前后际表达多样胜解之句，他们全都被此六十二事囊括，唯依附于此而出现，唯系属、囊括于此而出现。}

\textbf{譬如，诸比丘！熟练的渔夫，或渔夫弟子，用细目的网撒在小水塘上，他如是想：『任何在此水塘里的粗大生类，它们全都被囊括，唯依附于此而出现，唯系属、囊括于此而出现。』如是，诸比丘！任何沙门或婆罗门为前际想者、后际想者、前后际想者、前后际随见者，就前后际表达多样胜解之句，他们全都被此六十二事囊括，唯依附于此而出现，唯系属、囊括于此而出现。}

\begin{enumerate}\item 如是论述完还灭，现在，为了显示「没有持见者能解脱于开示之网」，又开始说「诸比丘！因为任何」。这里，\textbf{囊括}，即只能在我这开示之网以内。\textbf{唯依附于此}，即唯依附、依止、终结于我这开示之网。\textbf{出现}是说什么？他们下沉也罢，上浮也罢，唯依附于我的开示之网而下沉、上浮。\textbf{系属于此}，即系属于我的这开示之网，为其绑定、囊括而出现，因为没有不被包摄的持见者。
\item \textbf{细目}，意即纲目细密、网眼细小。世尊如渔夫，开示如网，一万世界如小水，持六十二见者如粗大生类。如同站在此岸边的观察者，见到粗大生类被囊括之相，世尊则见到一切持见者在开示之网以内，如是当知此中譬喻的对照。\end{enumerate}

{\Large 147.} \textbf{已切断有的导向，诸比丘！如来的身体住立，只要其身体还住立，人天便能看到他，在身坏命尽之后，人天便看不到他。}

\textbf{譬如，诸比丘！茎干被切断的芒果串，任何与茎干相连的芒果，全都随着它。如是，诸比丘！已切断有的导向，如来的身体住立，只要其身体还住立，人天便能看到他，在身坏命尽之后，人天便看不到他。}

\begin{enumerate}\item 如是，由一切见为此六十二见所摄故，在显示了一切持见者系属于此开示之网后，现在，为显示自身无所系属之状，说了「已切断有的导向，诸比丘！如来的身体」等等。
\item 这里，以之引领者即\textbf{导向}。引领者，即套住脖子后牵引，为绳索之名。但此处，由与导向相似故，以有爱为导向。因为它套住大众的脖子，引领、带领至彼彼有，而为\textbf{有的导向}。以阿罗汉道之剑切断其有的导向，即\textbf{已切断有的导向}。\textbf{身坏之后}，即从身坏以后。\textbf{命尽}，即命的彻底穷尽、灭尽，意即不再结生之状。\textbf{便看不到他}，他即如来。天或者人将看不见，意即至于不可施设之状。
\item 而在「\textbf{譬如，诸比丘}」的譬喻中，对照如下。芒果树，如如来的身体，树上所生的大茎干，如依于此的先前转起的渴爱，依附于此茎干的五颗熟果、十二颗熟果、十八颗熟果之量的芒果串，如渴爱尚存时，依附于渴爱的将来转生的五蕴、十二处、十八界。然而，好比此茎干被切断，所有芒果全都随着它，跟随此茎干，意即由茎干的切断而被切断，如是，由有的导向之茎干未被切断故，在将来会生起的五蕴、十二处、十八界，所有这些法都随着它，跟随有的导向，意即由彼的切断而被切断。
\item 然而，好比当此树由于曼荼迦刺之毒触，渐渐枯萎而死时，只留下「此处曾有名为这般的树」之言说，无人可见此树，如是，由于圣道之触，渴爱黏腻的穷尽故，当此身仿佛渐渐枯萎而破坏时，身坏命尽之后，他们将看不见他，人天将看不见如来，只留下「据说，这是名为这般的大师的教法」之言说。如此，便得达无余依涅槃界，完成了开示。\end{enumerate}

{\Large 148.} \textbf{如是说已，阿难尊者便对世尊说：「希有！尊者！未曾有！尊者！此法门为何名？尊者！」}

\textbf{「所以，阿难！对此法门，你应以『义网』持之，应以『法网』持之，应以『梵网』持之，应以『见网』持之，应以『无上的战斗胜利』持之！」世尊说罢。}

\begin{enumerate}\item \textbf{如是说已，阿难尊者}，如是，当世尊说完此经，长老从头存念了整部经，如是显明了佛力后，想到「所说的经名世尊还没取，噫！让我来请名」，便对世尊说了这些。
\item 在「所以……」等中，其意义的连结为：\textbf{阿难}！因为在此法门中，此世的义利、来世的义利均已分别，\textbf{所以，对此法门，你应以「义网」持之}，而因为其中讲了许多经典之法，所以\textbf{应以「法网」持之}，又因为其中以最胜义分别了梵的一切知智，所以\textbf{应以「梵网」持之}，又因为其中分别了六十二见，所以\textbf{应以「见网」持之}，而因为听闻此法门后，堪能粉碎天子魔、蕴魔、死魔、烦恼魔，所以\textbf{应以「无上的战斗胜利」持之}！
\item \textbf{世尊说罢}，即世尊为阐明以他人之慧不可得住、最上甚深的一切知智，为驱散见的大黑暗，如同太阳驱散黑暗，而说了这从因缘的终了\footnote{依据义注，本经的第 1-4 段为因缘部分。}开始，直到「应以『无上的战斗胜利』持之」的整部经。\end{enumerate}

{\Large 149.} \textbf{彼众比丘心满意足，欢喜于世尊之所说。当在说此记说时，一万世界发生震动。}

\begin{enumerate}\item \textbf{彼众比丘心满意足}，即彼众比丘心满意足、愉悦，即是说以对佛陀的喜而心生踊跃。\textbf{世尊之所说}，即如是与多种理法开示之优雅相应的此经，是世尊以迦陵频伽之鸣般的美妙悦耳，以甘露灌顶于智者心中般所说的话语。
\item \textbf{欢喜}，即随喜和领受。因为这「欢喜」一字在\begin{quoting}欢喜、欢迎。{\bluefootnotesize 〔相应部第 3.5 段〕}\end{quoting}等处是关于渴爱，在\begin{quoting}天人两者都很欢喜食物。{\bluefootnotesize 〔相应部第 1.43 段〕}\end{quoting}等处则关于趋近，在\begin{quoting}久客异乡者，自远处安归，亲友与知识，欢喜而迎彼。{\bluefootnotesize 〔法句第 219 颂〕}\end{quoting}等处则关于领受，在\begin{quoting}欢喜、随喜。{\bluefootnotesize 〔中部第 1.205 段〕}\end{quoting}等处则关于随喜。它在此处适合随喜和领受，因此说「欢喜，即随喜和领受」。\begin{quoting}对如如者的善说善谈，众比丘以\\「善哉善哉」而随喜，顶礼领受。\end{quoting}
\item \textbf{此记说}，即此无颂之经。由无偈颂故，它被称为记说\footnote{参见本经论因缘·九分的叙述。}。\textbf{一万世界}，即一万轮围之量的世界。\textbf{发生震动}，当知不只是在经的终了才发生震动。因为是说「\textbf{在说}」，所以，当知于六十二见展开被开示时，在一一见终了，计六十二处，发生震动。
\item 这里，当知以八种原因，大地震动：以界的扰动、神变的威力、菩萨的入胎、从母胎出生、证等觉、转法轮、舍寿行、般涅槃。彼等的抉择，我们将在「大般涅槃」中对出现的经文——\begin{quoting}阿难！此八因八缘，致使大地的震动显现。\end{quoting}作注释时再说。
\item 而此大地在另八处也发生震动：在大出离时、前往菩提座时、取粪扫衣时、洗粪扫衣时、黑园经时\footnote{黑园经：见\textbf{增支部}第 4.24 段。}、乔达摩经时\footnote{乔达摩支提经：见\textbf{增支部}第 3.126 段。}、毗输安多罗本生时\footnote{毗输安多罗本生：即\textbf{本生}的最后一则。}、此梵网时。这里，在大出离、前往菩提座时，以精进力而震动。在取粪扫衣时，在舍弃了二千小洲围绕的四大洲后，出家去往塚间，世尊以拾取粪扫衣而行难行之事，由希有的力量的冲击而震动。在洗粪扫衣、毗输安多罗本生时，以非时之震动而震动。在黑园、乔达摩经时，以「我作证为世尊」的作证之相而震动。而在此梵网时，当知于六十二见经解开、厘清而被开示时，因给予鼓掌而震动。
\item 且大地不仅在这些处震动，还在三次结集时，在摩诃摩哂陀长老来到此岛、坐在光辉林\footnote{光辉林 \textit{Jotivana}：\textbf{巴利专名佛教辞典} \textit{DPPN} 中说，原名欢喜林 \textit{Nandanavana}，在摩哂陀长老在那里开示后改名。它靠近阿耨楼陀城的南门内，隶属大寺，后来，大军 \textit{Mahāsena} 在光辉林中建了祇园寺。}中开示法的那天也震动过。在善好寺\footnote{善好寺 \textit{Kalyāṇiyavihāra}：待考。}中，乞食长老打扫完支提的庭院，便于彼处坐下，在以佛陀为所缘得喜之后，开始诵此经，在经的终了，以水为边界而震动。铜殿\footnote{铜殿 \textit{Lohapāsāda}：也在阿耨楼陀城内。}之中，有一处名为东芒果树园，诵长部的长老们便在那里开始梵网经，在他们诵习的终了，大地也以水为边界而震动。\begin{quoting}如是，以这自成者开示的最胜之经的威力，大地数数震动。\\于此，智者应恭敬地学习此梵网的法与义，而如理地行道！\end{quoting}\end{enumerate}
